% Panafricanisme et développement de l’Afrique — Philosophie Tle Tech — Cours signé OnBuch+
% Programme officiel MINESEC. Compiler avec : tectonic N20-panafricanisme-developpement-afrique.tex
\newcommand{\DOCMATIERE}{Philosophie}
\newcommand{\DOCNIVEAU}{Tle Tech}
\newcommand{\DOCMODULE}{Philosophie – classe de Terminale technique}
\newcommand{\DOCLECON}{N20}
\newcommand{\DOCTITRE}{Panafricanisme et développement de l’Afrique}
% OnBuch+ — préambule « cours signés » ENRICHI (compilé avec tectonic / XeLaTeX)
% Système visuel « Pop » d'OnBuch+ : fond crème (jamais blanc), bordures encre
% épaisses, ombres « dures » décalées, titres Archivo Black en capitales.
% Version MATHÉMATIQUES : graphes (pgfplots/tikz), boîtes pédagogiques
% (définition, propriété, méthode, attention, le savais-tu, exercice résolu,
% exercice, TP, bilan), page de garde et signature OnBuch+ ancrée en bas.
%
% Macros attendues dans main.tex AVANT \input{preamble} :
%   \DOCMATIERE  \DOCNIVEAU  \DOCMODULE  \DOCLECON (numéro)  \DOCTITRE
\documentclass[11pt]{article}

\usepackage{fontspec}
\usepackage[french]{babel}
\usepackage[a4paper,margin=2cm,top=3.4cm,bottom=2.8cm,headheight=1.4cm,footskip=1.3cm]{geometry}
\usepackage{amsmath,amssymb,amsfonts}
\usepackage{mathtools} % \xmapsto, \coloneqq... souvent utilisés par les modèles
\usepackage{cancel} % simplifications barrées (bilans de forces)
% \abs{z} (module, valeur absolue) et \norm{u} (norme), utilisés par les modèles
\providecommand{\abs}{}\let\abs\relax\DeclarePairedDelimiter{\abs}{\lvert}{\rvert}
\providecommand{\norm}{}\let\norm\relax\DeclarePairedDelimiter{\norm}{\lVert}{\rVert}
\usepackage[table]{xcolor} % [table] : fournit aussi \rowcolors (alternance de lignes)
\usepackage{pagecolor}
\usepackage{tikz}
\usetikzlibrary{arrows.meta,positioning,calc,shapes.geometric,shapes.misc,shapes.arrows,decorations.pathmorphing,decorations.pathreplacing,decorations.markings,patterns,shadows,mindmap,trees,fit,backgrounds,matrix,intersections,angles,quotes}
\usepackage{pgfplots}
\usepackage[version=4]{mhchem} % équations nucléaires \ce{^{235}_{92}U}
\usepackage[european,siunitx]{circuitikz} % schémas électriques
\pgfplotsset{compat=1.18}
\usepgfplotslibrary{fillbetween,groupplots} % aires entre courbes (calcul intégral)
\usepackage{enumitem}
\usepackage{fancyhdr}
% [most] charge toutes les bibliothèques tcolorbox (listings, minted, raster,
% documentation...) et gonfle inutilement la mémoire TeX sur un document long
% (~300 boîtes) : on ne charge que ce qui est réellement utilisé.
\usepackage[skins,breakable]{tcolorbox}
\usepackage{graphicx}
\usepackage{colortbl}
\usepackage{booktabs}
\usepackage{tabularx}
\usepackage{multirow}
\usepackage{diagbox} % case d'en-tête diagonale des tableaux à double entrée
% Colonnes extensibles centrée / gauche / droite (C, L, R), souvent utilisées par les modèles
\newcolumntype{C}{>{\centering\arraybackslash}X}
\newcolumntype{L}{>{\raggedright\arraybackslash}X}
\newcolumntype{R}{>{\raggedleft\arraybackslash}X}
\usepackage{array}
\usepackage{multicol}
\usepackage{siunitx}
% parse-numbers=false : accepte aussi \SI{2,5 \times 10^{-3}}{...}
\sisetup{locale=FR,output-decimal-marker={,},group-minimum-digits=4,parse-numbers=false}
\DeclareSIUnit\ohm{\ensuremath{\symup{\Omega}}} % U+2126 absent de la police
\DeclareSIUnit\division{div} % réglages d'oscilloscope (ms/div, V/div)
\DeclareSIUnit\tr{tr} % tours (vitesse de rotation, ex. \SI{1500}{\tr\per\minute})
\DeclareSIUnit\molar{M} % concentration molaire (ex. \SI{3}{\molar}, \SI{1}{\micro\molar})
\DeclareSIUnit\an{an} % année (le modèle confond parfois avec \year, un compteur TeX) — ex. \SI{5}{\kilo\meter\per\an}
\DeclareSIUnit\FCFA{FCFA} % franc CFA (ex. \SI{500}{\FCFA\per\kilo\gram})
\DeclareSIUnit\degreeBrix{\textdegree Brix} % teneur en sucre d'un sirop/jus (réfractométrie)
\DeclareSIUnit\month{mois} % le modèle utilise parfois \month, un compteur TeX, comme unité de durée
\DeclareSIUnit\microliter{\micro\liter} % le modèle écrit parfois \microliter en un seul token
\DeclareSIUnit\million{millions} % dénombrement (ex. \SI{15}{\million\per\mL} spermatozoïdes)
\usepackage{qrcode}
% \degree : commande fréquemment utilisée par les modèles pour un angle ou une
% température en dehors de \SI{}{} (non fournie par les paquets chargés).
\providecommand{\degree}{^{\circ}}
% \qed (fin de démonstration), utilisé par les modèles sans amsthm
\providecommand{\qed}{\hfill$\square$}
% \barre : double barre des valeurs interdites dans un tableau de variations (tkz-tab)
\providecommand{\barre}{\ensuremath{\|}}
% environnement proof (amsthm non chargé) utilisé par les modèles
\ifdefined\proof\else\newenvironment{proof}[1][Démonstration]{\par\noindent\textit{#1.}\ }{\qed\par}\fi
\usepackage{lastpage}
\usepackage{hyperref}
\hypersetup{hidelinks}

% ============================================================
% Palette « Pop » OnBuch+ — mêmes hex que la classe P (plus_ui.dart)
% ============================================================
\definecolor{poporange}{HTML}{F59321}
\definecolor{popdark}{HTML}{E07A0C}
\definecolor{popcream}{HTML}{FFF6E8}
\definecolor{popink}{HTML}{1C1714}
\definecolor{poppurple}{HTML}{7A5AE0}
\definecolor{popgreen}{HTML}{2E9E6A}
\definecolor{poppink}{HTML}{E0457B}
\definecolor{popblue}{HTML}{2F7DE1}
\definecolor{popgold}{HTML}{C98A12}
\definecolor{popmuted}{HTML}{8A7F76}
\definecolor{popline}{HTML}{EADFD2}
\definecolor{popgray}{HTML}{8A7F76}
\colorlet{popgrey}{popgray}
% Alias tolérants (le modèle nomme parfois les couleurs différemment)
\colorlet{popwhite}{white}
\colorlet{popblack}{popink}
\colorlet{popred}{poppink}
\colorlet{popyellow}{popgold}
% Teintes claires (fonds de tableaux/encadrés) — le modèle nomme parfois
% les couleurs avec un suffixe L (ex. popblueL) sans les avoir définies.
\colorlet{poporangeL}{poporange!20}
\colorlet{popdarkL}{popdark!20}
\colorlet{poppurpleL}{poppurple!20}
\colorlet{popgreenL}{popgreen!20}
\colorlet{poppinkL}{poppink!20}
\colorlet{popblueL}{popblue!20}
\colorlet{popgoldL}{popgold!20}
\colorlet{popredL}{popred!20}
\colorlet{popgrayL}{popgray!20}
% Teintes claires pour fonds de tableaux / zones
\colorlet{popblueL}{popblue!10!white}
\colorlet{poporangeL}{poporange!14!white}
\colorlet{popgreenL}{popgreen!12!white}
\colorlet{poppurpleL}{poppurple!12!white}
\colorlet{poppinkL}{poppink!12!white}
\colorlet{popgoldL}{popgold!14!white}

\pagecolor{popcream}

% ============================================================
% Typographie
% ============================================================
\defaultfontfeatures{Ligatures=TeX}
\setmainfont{PlusJakartaSans-Medium.ttf}[Path=../../../fonts/,BoldFont=PlusJakartaSans-Bold.ttf]
% Maths en sans-serif (Fira Math) avec chiffres et lettres droites de Plus
% Jakarta, pour que les formules se fondent dans le texte.
\usepackage[math-style=ISO,bold-style=ISO]{unicode-math}
\setmathfont{FiraMath-Regular.otf}[Path=../../../fonts/]
\setmathfont{PlusJakartaSans-Medium.ttf}[Path=../../../fonts/,range={up/num,up/{Latin,latin},\mathup}]
\usepackage{icomma} % virgule décimale sans espace en mode math
% Caractères mathématiques Unicode absents des polices de texte (Plus Jakarta,
% Archivo Black) : rendus par leur équivalent mathématique, en texte comme en
% formule — sinon ils s'affichent en carré vide (ex. « Arithmétique dans ℤ »).
\usepackage{newunicodechar}
\newunicodechar{ℝ}{\ensuremath{\mathbb{R}}}
\newunicodechar{ℤ}{\ensuremath{\mathbb{Z}}}
\newunicodechar{ℂ}{\ensuremath{\mathbb{C}}}
\newunicodechar{ℕ}{\ensuremath{\mathbb{N}}}
\newunicodechar{ℚ}{\ensuremath{\mathbb{Q}}}
\newunicodechar{𝔻}{\ensuremath{\mathbb{D}}}
\newunicodechar{α}{\ensuremath{\alpha}}
\newunicodechar{β}{\ensuremath{\beta}}
\newunicodechar{γ}{\ensuremath{\gamma}}
\newunicodechar{δ}{\ensuremath{\delta}}
\newunicodechar{ε}{\ensuremath{\varepsilon}}
\newunicodechar{θ}{\ensuremath{\theta}}
\newunicodechar{λ}{\ensuremath{\lambda}}
\newunicodechar{μ}{\ensuremath{\mu}}
\newunicodechar{σ}{\ensuremath{\sigma}}
\newunicodechar{τ}{\ensuremath{\tau}}
\newunicodechar{φ}{\ensuremath{\varphi}}
\newunicodechar{ψ}{\ensuremath{\psi}}
\newunicodechar{ω}{\ensuremath{\omega}}
\newunicodechar{Σ}{\ensuremath{\Sigma}}
% majuscules produites par \MakeUppercase dans les titres (α-aminés -> Α)
\newunicodechar{Α}{\ensuremath{\alpha}}
\newunicodechar{Β}{\ensuremath{\beta}}
\newunicodechar{Γ}{\ensuremath{\Gamma}}
\newunicodechar{Δ}{\ensuremath{\Delta}}
\newunicodechar{Ε}{\ensuremath{\varepsilon}}
\newunicodechar{Θ}{\ensuremath{\Theta}}
\newunicodechar{Λ}{\ensuremath{\Lambda}}
\newunicodechar{Μ}{\ensuremath{\mu}}
\newunicodechar{Τ}{\ensuremath{\tau}}
\newunicodechar{Φ}{\ensuremath{\Phi}}
\newunicodechar{Ψ}{\ensuremath{\Psi}}
\newunicodechar{Ω}{\ensuremath{\Omega}}
\newunicodechar{↦}{\ensuremath{\mapsto}}
\newunicodechar{∈}{\ensuremath{\in}}
\newunicodechar{∉}{\ensuremath{\notin}}
\newunicodechar{∧}{\ensuremath{\wedge}}
\newunicodechar{⇔}{\ensuremath{\Leftrightarrow}}
\newunicodechar{⟺}{\ensuremath{\Longleftrightarrow}}
\newunicodechar{⇒}{\ensuremath{\Rightarrow}}
\newunicodechar{⟹}{\ensuremath{\Longrightarrow}}
\newunicodechar{∘}{\ensuremath{\circ}}
\newunicodechar{∪}{\ensuremath{\cup}}
\newunicodechar{∩}{\ensuremath{\cap}}
\newunicodechar{≡}{\ensuremath{\equiv}}
\newunicodechar{⊂}{\ensuremath{\subset}}
\newunicodechar{⊥}{\ensuremath{\perp}}
\newunicodechar{∃}{\ensuremath{\exists}}
\newunicodechar{∀}{\ensuremath{\forall}}
\newunicodechar{∛}{\ensuremath{\sqrt[3]{\,}}}
\newunicodechar{∜}{\ensuremath{\sqrt[4]{\,}}}
\newunicodechar{⃗}{\ensuremath{{}^{\to}}}
\newunicodechar{ⁿ}{\ifmmode^{n}\else\textsuperscript{\ensuremath{n}}\fi}
\newunicodechar{ˣ}{\ifmmode^{x}\else\textsuperscript{\ensuremath{x}}\fi}
\newunicodechar{ᵇ}{\ifmmode^{b}\else\textsuperscript{\ensuremath{b}}\fi}
\newunicodechar{ʳ}{\ifmmode^{r}\else\textsuperscript{\ensuremath{r}}\fi}
\newunicodechar{ᵘ}{\ifmmode^{u}\else\textsuperscript{\ensuremath{u}}\fi}
\newunicodechar{ʸ}{\ifmmode^{y}\else\textsuperscript{\ensuremath{y}}\fi}
\newunicodechar{ᵃ}{\ifmmode^{a}\else\textsuperscript{\ensuremath{a}}\fi}
\newunicodechar{ᵉ}{\ifmmode^{e}\else\textsuperscript{\ensuremath{e}}\fi}
\newunicodechar{ᵏ}{\ifmmode^{k}\else\textsuperscript{\ensuremath{k}}\fi}
\newunicodechar{ᵖ}{\ifmmode^{p}\else\textsuperscript{\ensuremath{p}}\fi}
\newunicodechar{ᴺ}{\ifmmode^{N}\else\textsuperscript{\ensuremath{N}}\fi}
\newunicodechar{ᶜ}{\ifmmode^{c}\else\textsuperscript{\ensuremath{c}}\fi}
\newunicodechar{ʰ}{\ifmmode^{h}\else\textsuperscript{\ensuremath{h}}\fi}
\newunicodechar{ᵠ}{\ifmmode^{q}\else\textsuperscript{\ensuremath{q}}\fi}
\newunicodechar{ₙ}{\ifmmode_{n}\else\textsubscript{\ensuremath{n}}\fi}
\newunicodechar{ₐ}{\ifmmode_{a}\else\textsubscript{\ensuremath{a}}\fi}
\newunicodechar{ᵢ}{\ifmmode_{i}\else\textsubscript{\ensuremath{i}}\fi}
\newunicodechar{ᵣ}{\ifmmode_{r}\else\textsubscript{\ensuremath{r}}\fi}
\newunicodechar{ₖ}{\ifmmode_{k}\else\textsubscript{\ensuremath{k}}\fi}
\newunicodechar{ᵦ}{\ifmmode_{\beta}\else\textsubscript{\ensuremath{\beta}}\fi}
\newunicodechar{ₓ}{\ifmmode_{x}\else\textsubscript{\ensuremath{x}}\fi}
\newfontfamily\popdisplay{ArchivoBlack-Regular.ttf}[Path=../../../fonts/]
\newfontfamily\poplogo{SpaceGrotesk-Bold.ttf}[Path=../../../fonts/]
\newfontfamily\popsemi{PlusJakartaSans-SemiBold.ttf}[Path=../../../fonts/]
\newfontfamily\popxbold{PlusJakartaSans-ExtraBold.ttf}[Path=../../../fonts/]
\newfontfamily\popxboldls{PlusJakartaSans-ExtraBold.ttf}[Path=../../../fonts/,LetterSpace=3.0]

\setlist[enumerate]{leftmargin=1.7em,itemsep=5pt,topsep=4pt}
\setlist[itemize]{leftmargin=1.7em,itemsep=4pt,topsep=4pt,label={\color{poporange}\textbullet}}
\setlength{\parindent}{0pt}
\setlength{\parskip}{5pt}
\renewcommand{\arraystretch}{1.25}

% pgfplots : style Pop par défaut
\pgfplotsset{
  every axis/.append style={axis line style={popink,line width=1pt},tick style={popink},
    grid style={popline,line width=0.5pt},label style={font=\small},tick label style={font=\footnotesize}},
  popcurve/.style={poporange,line width=1.8pt,no markers,smooth},
}
% Même style disponible dans un \draw TikZ simple (hors pgfplots)
\tikzset{popcurve/.style={poporange,line width=1.8pt}}
\tikzset{popgrid/.style={popline,very thin}}
\colorlet{popcurve}{poporange} % le nom du style est parfois utilisé comme couleur

% ============================================================
% Pill / squiggle
% ============================================================
\newcommand{\poppill}[3]{%
  \tikz[baseline=-0.55ex]\node[draw=popink,line width=1.2pt,rounded corners=2.6pt,%
    fill=#2,inner xsep=6.5pt,inner ysep=3pt]{%
    \popxboldls\fontsize{7}{7}\selectfont\color{#3}\MakeUppercase{#1}};%
}
\newcommand{\popsquiggle}[1]{%
  \tikz[baseline=-0.4ex]{
    \draw[#1,line width=2.6pt,line cap=round]
      plot [smooth] coordinates {(0,0.09) (0.34,0.01) (0.68,0.21) (1.02,0.01) (1.36,0.10)};
  }%
}
% Mot-clé surligné (marqueur orange sous le texte)
\newcommand{\cle}[1]{{\popxbold\color{popdark}#1}}

% ============================================================
% En-tête / pied
% ============================================================
\pagestyle{fancy}
\fancyhf{}
\renewcommand{\headrulewidth}{0pt}
\renewcommand{\footrulewidth}{0pt}
\lhead{\raisebox{-0.15ex}{\poplogo\fontsize{13}{13}\selectfont\color{popink}OnBuch\color{poporange}+}}
\rhead{\poppill{\DOCMATIERE\ \textbullet\ \DOCNIVEAU\ \textbullet\ Leçon \DOCLECON}{white}{popink}}
\lfoot{\popxbold\fontsize{7.6}{7.6}\selectfont\color{popmuted}\MakeUppercase{Cours signé OnBuch+}\ \textbullet\ {\popsemi\DOCTITRE}}
\rfoot{\tikz[baseline=-0.6ex]\node[rounded corners=8.5pt,fill=popink,minimum height=17pt,minimum width=17pt,inner xsep=5pt,inner sep=0]{\popxbold\fontsize{8}{8}\selectfont\color{popcream}\thepage\,/\,\pageref*{LastPage}};}

% ============================================================
% Titres
% ============================================================
\newcommand{\chaptitre}[2]{%
  \begin{tcolorbox}[enhanced,colback=poporange,colframe=popink,boxrule=2.6pt,arc=11pt,
      drop shadow=popink,left=15pt,right=15pt,top=12pt,bottom=11pt,before skip=3pt,after skip=18pt]
    \poppill{#2}{popink}{popcream}\par\vspace{9pt}
    {\raggedright\hyphenpenalty=10000\exhyphenpenalty=10000\popdisplay\fontsize{22}{24}\selectfont\color{popcream}\MakeUppercase{#1}\par}
  \end{tcolorbox}
}
% Section numérotée : pastille orange avec le numéro + titre Archivo
\newcounter{popsec}
\newcommand{\coursec}[1]{%
  \stepcounter{popsec}\setcounter{popsub}{0}%
  \par\vspace{16pt}\noindent
  \begin{minipage}{\linewidth}
  \tikz[baseline=-0.7ex]\node[draw=popink,line width=1.6pt,fill=poporange,rounded corners=4pt,minimum size=22pt,inner sep=0]{\popdisplay\fontsize{12}{12}\selectfont\color{popcream}\thepopsec};%
  \hspace{8pt}{\popdisplay\fontsize{15}{17}\selectfont\color{popink}\MakeUppercase{#1}}\par
  \vspace{4pt}\noindent{\color{poporange}\rule{36pt}{3.6pt}}
  \end{minipage}\par\nopagebreak\vspace{8pt}
}
\newcounter{popsub}
\newcommand{\courssub}[1]{%
  \stepcounter{popsub}%
  \par\vspace{9pt}\noindent{\popxbold\fontsize{11.5}{13}\selectfont\color{popink}{\color{poporange}\thepopsec.\thepopsub}\hspace{6pt}#1}\par\nopagebreak\vspace{4pt}
}

% ============================================================
% Boîtes pédagogiques — même recette : fond blanc, bordure épaisse colorée,
% ombre dure encre, badge plein. Titre optionnel [#1].
% ============================================================
\tcbset{popbase/.style={breakable,enhanced,colback=white,boxrule=2.3pt,arc=9pt,
  shadow={3pt}{-3pt}{0pt}{popink},left=14pt,right=14pt,top=11pt,bottom=11pt,before skip=11pt,after skip=15pt,
  fontupper=\popsemi\fontsize{10.6}{14.5}\selectfont\color{popink},
  fontlower=\popsemi\fontsize{10.6}{14.5}\selectfont\color{popink}}}
% Coche dessinée (le glyphe ✓ n'existe pas dans les polices du document)
\newcommand{\popcheck}[1]{\tikz[baseline=-0.5ex]\draw[#1,line width=1.6pt,line cap=round,line join=round](-0.13,0.0)--(-0.04,-0.09)--(0.13,0.1);}
\providecommand{\coche}{\popcheck{popgreen}}
\providecommand{\resultat}[1]{\par\smallskip\noindent\fcolorbox{popgreen}{popgreenL}{\parbox{\dimexpr\linewidth-2\fboxsep-2\fboxrule\relax}{\textbf{Résultat.} #1}}\par\smallskip}
\newcommand{\popboxhead}[3]{\poppill{#1}{#2}{white}\ifx\relax#3\relax\else\hspace{6pt}{\popxbold\fontsize{10.6}{12}\selectfont\color{popink}#3}\fi\par\vspace{7pt}}

\newtcolorbox{aretenir}[1][]{popbase,colframe=popgreen,before upper={\popboxhead{À retenir}{popgreen}{#1}}}
\newtcolorbox{exemplebox}[1][]{popbase,colframe=poporange,before upper={\popboxhead{Exemple}{poporange}{#1}}}
\newtcolorbox{definition}[1][]{popbase,colframe=popblue,before upper={\popboxhead{Définition}{popblue}{#1}}}
\newtcolorbox{propriete}[1][]{popbase,colframe=poppurple,before upper={\popboxhead{Propriété}{poppurple}{#1}}}
\newtcolorbox{methode}[1][]{popbase,colframe=popgold,before upper={\popboxhead{Méthode}{popgold}{#1}}}
\newtcolorbox{attention}[1][]{popbase,colframe=poppink,before upper={\popboxhead{Attention}{poppink}{#1}}}
\newtcolorbox{experience}[1][]{popbase,colframe=popblue,colback=popblueL,before upper={\popboxhead{Expérience}{popblue}{#1}}}
% « Le savais-tu ? » : carte violette pleine (contexte, culture, Cameroun)
\newtcolorbox{savaistu}[1][]{popbase,colback=poppurple,colframe=popink,
  fontupper=\popsemi\fontsize{10.4}{14.2}\selectfont\color{white},
  before upper={\poppill{Le savais-tu\,?}{white}{poppurple}\ifx\relax#1\relax\else\hspace{6pt}{\popxbold\color{white}#1}\fi\par\vspace{7pt}}}
% Exercice résolu : énoncé en haut, \tcblower puis la solution sur fond crème
\newcounter{popexo}\newcounter{popexr}
\newtcolorbox{exoresolu}[1][]{popbase,colframe=poporange,colbacklower=popcream,
  segmentation style={popink,line width=1.2pt,dashed},
  before upper={\stepcounter{popexr}\popboxhead{Exercice résolu \thepopexr}{poporange}{#1}},
  before lower={\poppill{Solution}{popink}{popcream}\par\vspace{6pt}}}
% Exercice (non résolu en place) — la correction est dans la partie Corrigés
\newtcolorbox{exercice}[1][]{popbase,colframe=popink,
  before upper={\stepcounter{popexo}\popboxhead{Exercice \thepopexo}{popink}{#1}}}
% Corrigé
\newtcolorbox{corrige}[1][]{popbase,colframe=popgreen,colback=popgreenL,
  before upper={\popboxhead{Corrigé}{popgreen}{#1}}}
% Objectifs / prérequis / bilan
\newtcolorbox{objectifs}{popbase,colframe=popink,colback=poporangeL,
  before upper={\popboxhead{Objectifs de la leçon}{popink}{}}}
\newtcolorbox{prerequis}{popbase,colframe=popink,colback=popblueL,
  before upper={\popboxhead{Prérequis}{popblue}{}}}
\newtcolorbox{bilan}{popbase,colframe=popink,colback=white,boxrule=2.8pt,
  before upper={\popboxhead{Fiche bilan}{poporange}{L'essentiel à connaître pour l'examen}}}
% Figure encadrée avec légende
\newtcolorbox{popfigure}[1][]{enhanced,breakable=false,colback=white,colframe=popline,boxrule=1.4pt,arc=8pt,
  left=10pt,right=10pt,top=10pt,bottom=8pt,before skip=10pt,after skip=12pt,halign=center,
  lower separated=false,halign lower=center,
  fontlower=\popsemi\fontsize{9}{11.5}\selectfont\color{popmuted},
  #1}
\newcommand{\legende}[1]{\par\vspace{6pt}{\popsemi\fontsize{9}{11.5}\selectfont\color{popmuted}#1}}

% ============================================================
% Signature OnBuch+
% ============================================================
\newcommand{\poplogoinline}[1]{{\poplogo\fontsize{#1}{#1}\selectfont\color{popink}OnBuch\color{poporange}+}}
\newcommand{\signatureonbuch}{%
  \begin{tcolorbox}[enhanced,colback=popink,colframe=popink,boxrule=2.2pt,arc=10pt,
      drop shadow=poporange,left=14pt,right=14pt,top=10pt,bottom=10pt,before skip=0pt,after skip=0pt]
    \tikz[baseline=-0.7ex]\node[circle,fill=popgreen,draw=popcream,line width=1.3pt,minimum size=22pt,inner sep=0]{\popcheck{white}};%
    \hspace{9pt}\begin{minipage}[c]{0.62\linewidth}
      {\popxbold\fontsize{12}{13}\selectfont\color{popcream}Signé\ \textbullet\ {\poplogo OnBuch\color{poporange}+}}\par\vspace{2pt}
      {\raggedright\popsemi\fontsize{8}{10.5}\selectfont\color{popline}\MakeUppercase{Équipe pédagogique OnBuch+}\\\MakeUppercase{Conforme au programme officiel MINESEC}\par}
    \end{minipage}\hfill
    {\poplogo\fontsize{22}{22}\selectfont\color{popcream}OnBuch\color{poporange}+}
  \end{tcolorbox}%
}

% ============================================================
% Page de garde
% #1 titre · #2 matière · #3 niveau · #4 module · #5 n° leçon · #6 durée ·
% #7 description / objectifs (texte ou itemize)
% ============================================================
\newcommand{\pagegarde}[7]{%
  \thispagestyle{empty}
  \newgeometry{a4paper,margin=1.7cm,top=1.6cm,bottom=1.4cm}%
  \begin{tikzpicture}[remember picture,overlay]
    \fill[poporange!18] ([xshift=-3.2cm,yshift=-3.6cm]current page.north east) circle (4.6cm);
    \fill[poppurple!14] ([xshift=2.4cm,yshift=7.6cm]current page.south west) circle (3.8cm);
  \end{tikzpicture}%
  {\poplogo\fontsize{30}{30}\selectfont\color{popink}OnBuch\color{poporange}+}\hfill
  \raisebox{0.6ex}{\poppill{Cours signé \textbullet\ Premium}{popink}{popcream}}\par
  \vspace{1pt}\popsquiggle{poppurple}\par
  \vspace{8pt}
  \begin{tcolorbox}[enhanced,colback=poporange,colframe=popink,boxrule=2.8pt,arc=13pt,
      drop shadow=popink,left=18pt,right=18pt,top=12pt,bottom=13pt,after skip=10pt]
    \poppill{#2 \textbullet\ #3}{popink}{popcream}\hspace{5pt}\poppill{#4}{white}{popink}\par\vspace{8pt}
    {\popdisplay\fontsize{14}{14}\selectfont\color{popink}LEÇON #5}\par\vspace{4pt}
    {\raggedright\hyphenpenalty=10000\exhyphenpenalty=10000\popdisplay\fontsize{25}{27}\selectfont\color{popcream}\MakeUppercase{#1}\par}
    \vspace{8pt}
    {\popxbold\fontsize{9.5}{9.5}\selectfont\color{popink}\MakeUppercase{Durée conseillée : #6}}
  \end{tcolorbox}
  \begin{tcolorbox}[enhanced,colback=white,colframe=popink,boxrule=2.2pt,arc=10pt,
      drop shadow=popink,left=16pt,right=16pt,top=10pt,bottom=10pt,after skip=8pt]
    \poppill{Dans cette leçon}{poppurple}{white}\par\vspace{5pt}
    \popsemi\fontsize{9.3}{12.6}\selectfont\color{popink}#7
  \end{tcolorbox}
  \noindent\begin{minipage}[t]{0.487\linewidth}
    \begin{tcolorbox}[enhanced,colback=popgreen,colframe=popink,boxrule=2.2pt,arc=10pt,
        drop shadow=popink,left=11pt,right=11pt,top=10pt,bottom=10pt,height=3.1cm,valign=center]
      \tcbox[colback=white,colframe=popink,boxrule=1.3pt,arc=5pt,boxsep=0pt,nobeforeafter,
        left=3pt,right=3pt,top=3pt,bottom=3pt,valign=center]{\qrcode[height=1.9cm]{https://play.google.com/store/apps/details?id=com.luvvix.onbuch&hl=fr}}%
      \hspace{8pt}\begin{minipage}[c]{0.5\linewidth}
        {\popdisplay\fontsize{11}{12.5}\selectfont\color{white}\MakeUppercase{Télécharge OnBuch}}\par\vspace{4pt}
        {\raggedright\popsemi\fontsize{8.4}{11}\selectfont\color{white}Gratuit \textbullet\ Google Play\\Tuteur IA, annales, concours, résultats\par}
      \end{minipage}
    \end{tcolorbox}
  \end{minipage}\hfill
  \begin{minipage}[t]{0.487\linewidth}
    \begin{tcolorbox}[enhanced,colback=poppurple,colframe=popink,boxrule=2.2pt,arc=10pt,
        drop shadow=popink,left=11pt,right=11pt,top=10pt,bottom=10pt,height=3.1cm,valign=center]
      \tikz[baseline=-0.7ex]\node[circle,fill=white,draw=popink,line width=1.3pt,minimum size=1.6cm,inner sep=0]{\popdisplay\fontsize{24}{24}\selectfont\color{poppurple}+};%
      \hspace{8pt}\begin{minipage}[c]{0.6\linewidth}
        {\popdisplay\fontsize{11}{12.5}\selectfont\color{white}\MakeUppercase{Passe à OnBuch+}}\par\vspace{4pt}
        {\raggedright\popsemi\fontsize{8.4}{11}\selectfont\color{white}Tous les cours signés, Tuteur IA illimité, fiches PDF — onglet OnBuch+ de l'app\par}
      \end{minipage}
    \end{tcolorbox}
  \end{minipage}
  \vspace{8pt}
  \popcontenu
  \vfill
  \signatureonbuch
  \restoregeometry
  \newpage
}

% Rangée « Ce cours contient » (page de garde)
\newcommand{\poptile}[3]{% #1 couleur #2 gros texte #3 légende
  \begin{tcolorbox}[enhanced,colback=#1,colframe=popink,boxrule=2pt,arc=9pt,shadow={2.5pt}{-2.5pt}{0pt}{popink},
      left=6pt,right=6pt,top=9pt,bottom=9pt,halign=center,valign=center,height=1.75cm,width=\linewidth,nobeforeafter]
    {\popdisplay\fontsize{12.5}{13}\selectfont\color{white}\MakeUppercase{#2}}\par\vspace{3pt}
    {\centering\popsemi\fontsize{7.8}{9.5}\selectfont\color{white}#3\par}
  \end{tcolorbox}}
\newcommand{\popcontenu}{%
  \par\noindent{\popxboldls\fontsize{7.5}{7.5}\selectfont\color{popmuted}CE COURS SIGNÉ CONTIENT}\par\vspace{7pt}
  \noindent\begin{minipage}[t]{0.235\linewidth}\poptile{popblue}{Cours}{complet, illustré, pas à pas}\end{minipage}\hfill
  \begin{minipage}[t]{0.235\linewidth}\poptile{poporange}{Méthodes}{et exercices résolus}\end{minipage}\hfill
  \begin{minipage}[t]{0.235\linewidth}\poptile{poppink}{Exercices}{type examen + corrigés}\end{minipage}\hfill
  \begin{minipage}[t]{0.235\linewidth}\poptile{popgreen}{Fiche bilan}{l'essentiel à retenir}\end{minipage}\par}

% Sommaire
\newtcolorbox{sommaire}{enhanced,colback=white,colframe=popink,boxrule=2.2pt,arc=10pt,
  drop shadow=popink,left=18pt,right=18pt,top=13pt,bottom=13pt,before skip=6pt,after skip=20pt,
  before upper={\poppill{Sommaire}{poppurple}{white}\par\vspace{9pt}\popsemi\fontsize{10.6}{14.5}\selectfont\color{popink}}}

% Signature de fin de cours — poussée tout en bas de la dernière page
\newcommand{\findecours}{%
  \par\vfill\nopagebreak
  \begin{center}\popsquiggle{poporange}\end{center}\vspace{4pt}
  \signatureonbuch
}
\AtBeginDocument{\def\llbracket{\mathopen{[\mkern-3.5mu[}}\def\rrbracket{\mathclose{]\mkern-3.5mu]}}} % intervalles d'entiers (glyphe ⟦ absent de la police)
% Glyphes absents de Fira Math : redéfinis à partir de glyphes disponibles
\AtBeginDocument{%
  \def\setminus{\mathbin{\backslash}}\let\smallsetminus\setminus
  \def\vdots{\mathord{\vbox{\baselineskip3.2pt\lineskiplimit0pt\kern4pt\hbox{.}\hbox{.}\hbox{.}}}}%
  \def\ddots{\mathinner{\mkern1mu\raise6.5pt\hbox{.}\mkern2mu\raise3.6pt\hbox{.}\mkern2mu\raise0.7pt\hbox{.}\mkern1mu}}%
  \def\triangle{\mathord{\scalebox{0.9}{$\symup{\Delta}$}}}%
  \def\nabla{\mathord{\raisebox{\depth}{\scalebox{1}[-1]{$\symup{\Delta}$}}}}%
  \def\checkmark{\popcheck{popdark}}%
  \def\star{\mathbin{\ast}}\def\bigstar{\mathord{\ast}}%
  \def\diamond{\mathbin{\scalebox{0.6}{$\Diamond$}}}%
  \def\bigcup{\mathop{\mathchoice{\raisebox{-0.3em}{\scalebox{1.7}{$\cup$}}}{\scalebox{1.25}{$\cup$}}{\cup}{\cup}}\displaylimits}%
  \def\bigcap{\mathop{\mathchoice{\raisebox{-0.3em}{\scalebox{1.7}{$\cap$}}}{\scalebox{1.25}{$\cap$}}{\cap}{\cap}}\displaylimits}%
  \def\lesseqqgtr{\mathrel{\lessgtr}}\def\bigoplus{\mathop{\oplus}\displaylimits}%
}
\providecommand{\ssi}{si et seulement si} % abréviation fréquente
\AtBeginDocument{\providecommand{\wideparen}{\overparen}} % arc de cercle

%% FIGFIX-BEGIN v5 — correctifs globaux des figures (docs/onbuch-generation/tools/figfix)
\usepackage{adjustbox}\usepackage{etoolbox}\tcbuselibrary{raster}
% toute figure TikZ/pgfplots plus large que la ligne est réduite à la largeur disponible
\BeforeBeginEnvironment{tikzpicture}{\begin{adjustbox}{max width=\linewidth}}
\AfterEndEnvironment{tikzpicture}{\end{adjustbox}}
% légendes pgfplots lisibles par-dessus les courbes
\pgfplotsset{every axis/.append style={legend style={fill=white,fill opacity=0.92,text opacity=1,draw=black!15,inner sep=3pt}}}
% étiquettes d'arêtes (syntaxe edge["..."]) avec fond blanc
\tikzset{every edge quotes/.append style={fill=white,inner sep=1.2pt}}
%% FIGFIX-END
\begin{document}

\pagegarde{Panafricanisme et développement de l’Afrique}{Philosophie}{Tle Tech}{Philosophie – classe de Terminale technique}{N20}{2 séances (fiche de progression harmonisée)}{Cette leçon définit les concepts de développement et de sous-développement, puis retrace l'historique, les objectifs et les difficultés du panafricanisme comme philosophie du développement de l'Afrique. Elle invite l'élève à appliquer ces notions à des situations concrètes du Cameroun et à construire une argumentation rigoureuse en vue du Baccalauréat technique.
\vspace{4pt}
\begin{multicols}{2}\small
\begin{itemize}
\item À la fin de cette leçon, je sais définir les concepts de développement et de sous-développement
\item À la fin de cette leçon, je sais distinguer croissance économique et développement humain à partir d'indicateurs chiffrés
\item À la fin de cette leçon, je sais retracer les grandes étapes historiques du panafricanisme (congrès, OUA, UA)
\item À la fin de cette leçon, je sais identifier les objectifs politiques, économiques et culturels du panafricanisme
\item À la fin de cette leçon, je sais analyser les difficultés internes et externes qui freinent l'unité africaine
\item À la fin de cette leçon, je sais évaluer la pertinence du panafricanisme comme philosophie du développement de l'Afrique
\end{itemize}\end{multicols}}

\begin{sommaire}
\begin{multicols}{2}
\begin{enumerate}
\item Le développement : définition et dimensions
\item Le sous-développement : concept et manifestations en Afrique
\item Historique du panafricanisme
\item Les objectifs du panafricanisme comme philosophie du développement
\item Les difficultés du panafricanisme
\item Synthèse : le panafricanisme peut-il développer l'Afrique ?
\item Activité d'intégration
\item Méthodes et exercices résolus
\item Exercices
\item Corrigés des exercices
\item Fiche bilan
\end{enumerate}
\end{multicols}
\end{sommaire}

\begin{objectifs}
\begin{itemize}
\item À la fin de cette leçon, je sais définir les concepts de développement et de sous-développement
\item À la fin de cette leçon, je sais distinguer croissance économique et développement humain à partir d'indicateurs chiffrés
\item À la fin de cette leçon, je sais retracer les grandes étapes historiques du panafricanisme (congrès, OUA, UA)
\item À la fin de cette leçon, je sais identifier les objectifs politiques, économiques et culturels du panafricanisme
\item À la fin de cette leçon, je sais analyser les difficultés internes et externes qui freinent l'unité africaine
\item À la fin de cette leçon, je sais évaluer la pertinence du panafricanisme comme philosophie du développement de l'Afrique
\item À la fin de cette leçon, je sais appliquer les notions de développement et d'unité à des situations concrètes de la vie courante au Cameroun
\item À la fin de cette leçon, je sais rédiger et justifier une démarche argumentative (thèse, antithèse, synthèse) sur un sujet de dissertation
\end{itemize}
\end{objectifs}

\begin{prerequis}
\begin{itemize}
\item Notion de colonisation et de décolonisation de l'Afrique (histoire de 1ère)
\item Notion de culture et d'aliénation culturelle (leçons de philosophie de Terminale)
\item Notion de liberté et de responsabilité (début d'année de Terminale)
\item Connaissance générale des institutions : ONU, OUA/UA, CEEAC
\end{itemize}
\end{prerequis}

\newpage

\coursec{Le développement : définition et dimensions}

\textbf{Situation d'accroche.} En 2024, un jeune diplômé camerounais de Douala constate que le cacao récolté à Obala est exporté brut vers l'Europe, puis revient sous forme de chocolat vendu dix fois plus cher dans les supermarchés de Bonanjo. Pendant ce temps, les dirigeants africains se réunissent à Addis-Abeba pour parler d'unité et de développement du continent. Pourquoi, soixante ans après les indépendances, l'Afrique reste-t-elle perçue comme « sous-développée » malgré ses immenses richesses ? Le panafricanisme, né du rêve d'une Afrique unie et prospère, peut-il encore constituer une philosophie du développement ? C'est à ces questions que cette leçon tente de répondre.

\textbf{Problématique.} Mais avant de demander si le panafricanisme peut développer l'Afrique, encore faut-il savoir de quoi l'on parle : qu'est-ce que le développement ? Se réduit-il à l'argent, aux usines, aux routes ? Ou bien engage-t-il toute la vie d'un peuple ? Telle est l'interrogation de cette première section.

\courssub{Étymologie et définition générale du développement}

Le mot \cle{développement} vient du verbe « développer », formé du préfixe \emph{dé-} (qui marque l'ouverture, le déploiement) et de « envelopper » : développer, c'est littéralement \emph{dés-envelopper}, c'est-à-dire déplier ce qui était replié, faire sortir ce qui était contenu en puissance. On pense au photographe qui « développe » une pellicule : l'image était déjà là, cachée ; le développement la rend visible. On pense aussi à la graine qui se développe en plante : l'arbre était en puissance dans la semence.

Cette intuition étymologique est précieuse : le développement n'est pas l'ajout de quelque chose qui viendrait de l'extérieur, mais le \emph{déploiement de potentialités internes}. Une société se développe lorsqu'elle fait fructifier ce qu'elle porte déjà en elle : ses ressources, ses savoirs, ses hommes.

\begin{definition}[Développement]
Le \cle{développement} est l'ensemble des transformations économiques, sociales, politiques et culturelles qui font passer une société d'un état inférieur à un état supérieur et qui améliorent \emph{durablement} les conditions de vie d'une population.
\end{definition}

Trois éléments de cette définition doivent être retenus et mobilisés dans une dissertation :

\begin{enumerate}
\item Le développement est un \textbf{processus} (un passage, un mouvement), non un état figé : on ne « possède » pas le développement, on se développe.
\item Il va d'un état jugé \textbf{inférieur} (pauvreté, dépendance, ignorance) vers un état jugé \textbf{supérieur} (aisance, autonomie, savoir) : il suppose donc une idée de \emph{progrès}.
\item Il doit être \textbf{durable} : un enrichissement passager (par exemple lié à la hausse momentanée du prix du pétrole) n'est pas un développement s'il ne transforme pas en profondeur la société.
\end{enumerate}

\begin{attention}[Piège de définition]
Ne définis jamais le développement comme « le fait d'être riche ». Un pays peut être riche en ressources (pétrole, cacao, cobalt) sans être développé, si cette richesse ne profite ni à la santé, ni à l'éducation, ni à la liberté de sa population. Le développement est un concept \emph{qualitatif} et \emph{global}, pas seulement quantitatif et monétaire.
\end{attention}

\courssub{Le développement économique : croissance, industrialisation, transformation structurelle}

La dimension la plus visible du développement est la dimension \cle{économique}. Elle repose sur trois mécanismes qu'il faut distinguer.

\textbf{a) La croissance du PIB.} Le Produit Intérieur Brut (PIB) mesure la valeur des richesses produites par un pays en une année. La \cle{croissance économique} est l'augmentation de ce PIB. C'est une condition du développement : difficile d'améliorer durablement la vie d'une population si le pays ne produit pas davantage de richesses.

\textbf{b) L'industrialisation.} Un pays développé ne se contente pas d'exporter des matières premières brutes : il les \emph{transforme} sur place. Le cacao d'Obala devient du chocolat fabriqué à Douala ; le coton du Nord-Cameroun devient du tissu et des vêtements. L'industrie crée des emplois qualifiés, des savoir-faire et de la valeur ajoutée qui reste au pays.

\textbf{c) La transformation structurelle.} C'est le passage d'une économie fondée sur l'agriculture de subsistance et l'exportation de produits bruts (économie dite « primaire ») à une économie diversifiée, où l'industrie et les services occupent une place croissante. Le développement économique n'est donc pas seulement « produire plus », c'est \emph{produire autrement}.

\begin{exemplebox}[Le paradoxe du cacao camerounais]
Le Cameroun est l'un des premiers producteurs mondiaux de cacao. Pourtant, l'essentiel de la fève est exportée brute. Or la valeur d'un produit augmente à chaque étape de transformation : la fève brute rapporte peu au planteur, tandis que le chocolat fini se vend dix fois plus cher. Tant que la transformation se fait à l'étranger, la plus grande partie de la richesse créée par le cacao camerounais profite à d'autres. Voilà pourquoi l'industrialisation locale est un enjeu central du développement économique africain.
\end{exemplebox}

\courssub{Le développement humain et social : santé, éducation, IDH}

Un pays peut voir son PIB augmenter sans que la vie quotidienne de ses habitants s'améliore. C'est pourquoi les Nations Unies ont construit un indicateur plus large : l'\cle{Indice de Développement Humain} (IDH).

\begin{definition}[IDH]
L'\cle{IDH} est un indicateur composite créé par le Programme des Nations Unies pour le Développement (PNUD). Il combine trois dimensions : la \textbf{santé} (espérance de vie à la naissance), l'\textbf{éducation} (durée moyenne et durée attendue de scolarisation) et le \textbf{niveau de vie} (revenu national brut par habitant). Il varie entre 0 et 1.
\end{definition}

L'échelle de lecture est la suivante : un IDH supérieur à 0,80 correspond à un développement humain \emph{très élevé} ; entre 0,70 et 0,80, \emph{élevé} ; entre 0,55 et 0,70, \emph{moyen} ; en dessous de 0,55, \emph{faible}.

\begin{exemplebox}[L'IDH du Cameroun]
Avec un IDH d'environ 0,58 (Rapport PNUD 2023/2024), le Cameroun se situe dans la catégorie « développement humain \textbf{moyen} », loin des pays à IDH très élevé (supérieur à 0,80) comme la France (environ 0,90) ou la Corée du Sud (environ 0,93). Cela signifie concrètement que l'espérance de vie, la scolarisation et le revenu des Camerounais restent nettement inférieurs à ceux des pays développés, malgré les richesses naturelles du pays.
\end{exemplebox}

Le développement humain et social recouvre donc des réalités concrètes : pouvoir se soigner (hôpitaux accessibles, médicaments disponibles), pouvoir apprendre (écoles, universités, formation technique), vivre longtemps et dignement (eau potable, alimentation suffisante, logement décent). L'économiste Amartya Sen, prix Nobel d'économie en 1998, résume cette idée en affirmant que le développement est un processus d'\emph{élargissement des libertés réelles} dont jouissent les individus : être développé, c'est être \emph{capable} de vivre la vie que l'on a des raisons de valoriser.

\courssub{Le développement politique et culturel}

Le développement possède encore deux dimensions, trop souvent oubliées dans les copies.

\textbf{a) La dimension politique.} Un pays ne se développe durablement que si ses institutions fonctionnent : c'est ce qu'on appelle la \cle{bonne gouvernance}. Elle implique l'État de droit (la loi s'applique à tous, y compris aux puissants), la lutte contre la corruption, la transparence dans la gestion des deniers publics, la participation des citoyens aux décisions et la stabilité des institutions. Une richesse mal gouvernée est une richesse gaspillée ou détournée.

\textbf{b) La dimension culturelle.} Se développer, ce n'est pas copier servilement les modèles étrangers : c'est aussi s'épanouir dans ses \emph{valeurs propres} — langues, arts, savoirs traditionnels, formes de solidarité communautaire. Un développement qui détruit l'identité culturelle d'un peuple produit des sociétés déracinées. Comme le soulignait Léopold Sédar Senghor, l'Afrique doit s'approprier les techniques modernes sans renier son âme : « s'assimiler, non être assimilés ».

\begin{popfigure}
\begin{center}
\begin{tikzpicture}
\fill[popblueL] (0,0) circle (3.4);
\fill[popgreenL] (0,0) circle (2.5);
\fill[poporangeL] (0,0) circle (1.6);
\fill[poppinkL] (0,0) circle (0.75);
\node[font=\small\bfseries, popdark] at (0,0) {Culturel};
\node[font=\small\bfseries, popdark] at (0,1.15) {Politique};
\node[font=\small\bfseries, popdark] at (0,2.05) {Social};
\node[font=\small\bfseries, popdark] at (0,2.95) {Économique};
\draw[->, thick, popdark] (3.6,0) -- (4.6,0) node[right, font=\small, text width=4.2cm]{Le développement est un tout : les quatre dimensions se renforcent mutuellement.};
\end{tikzpicture}
\end{center}
\legende{Les quatre dimensions du développement, représentées en cercles emboîtés : la dimension économique (extérieure) englobe et soutient les dimensions sociale et politique, jusqu'au cœur culturel qui donne son identité au développement d'un peuple.}
\end{popfigure}

\courssub{Distinction capitale : croissance économique et développement}

C'est la distinction la plus importante de cette section, et la plus rentable à l'examen.

\begin{definition}[Croissance économique]
La \cle{croissance économique} est l'augmentation quantitative de la production de richesses d'un pays (mesurée par le PIB) sur une période donnée.
\end{definition}

La croissance est \emph{quantitative} : elle compte. Le développement est \emph{qualitatif} : il interroge la répartition et l'usage des richesses. Un pays peut connaître une forte croissance — par exemple grâce à l'exploitation pétrolière — sans que la majorité de sa population en bénéficie : si les revenus sont concentrés entre quelques mains, si les écoles et les hôpitaux restent défaillants, il y a \textbf{croissance sans développement}.

\begin{attention}[Erreur fréquente en dissertation]
Confondre croissance et développement est la faute la plus répandue. Retiens la formule : \textbf{la croissance est une condition du développement, mais elle n'en est pas une condition suffisante}. Une croissance sans redistribution (santé, éducation, emplois, justice) ne développe pas. Dans une introduction ou une définition, écris toujours les deux termes et précise leur différence : c'est la marque d'une copie solide.
\end{attention}

La figure suivante illustre ce décalage : elle compare le revenu par habitant et l'IDH de quatre pays. On observe que le classement par la richesse et le classement par le développement humain se recoupent globalement, mais que l'IDH ajoute une information essentielle : ce que la richesse \emph{devient} dans la vie des gens.

\begin{popfigure}
\begin{center}
\begin{tikzpicture}
\begin{axis}[
ybar, bar width=0.5cm,
width=0.95\linewidth, height=7cm,
symbolic x coords={Cameroun,Nigeria,France,Coree du Sud},
xtick=data,
ymin=0, ymax=45000,
ylabel={PIB par habitant (dollars, environ)},
legend style={at={(0.5,-0.25)}, anchor=north, legend columns=2, font=\small},
ymajorgrids=true, grid style={popline!40},
every axis plot/.append style={fill}
]
\addplot[popblue] coordinates {(Cameroun,1600) (Nigeria,2100) (France,44000) (Coree du Sud,33000)};
\legend{PIB par habitant}
\end{axis}
\end{tikzpicture}

\vspace{0.4cm}
\begin{tikzpicture}
\begin{axis}[
ybar, bar width=0.5cm,
width=0.95\linewidth, height=6cm,
symbolic x coords={Cameroun,Nigeria,France,Coree du Sud},
xtick=data,
ymin=0, ymax=1,
ylabel={IDH (entre 0 et 1, environ)},
ymajorgrids=true, grid style={popline!40},
]
\addplot[poporange, fill] coordinates {(Cameroun,0.58) (Nigeria,0.55) (France,0.90) (Coree du Sud,0.93)};
\end{axis}
\end{tikzpicture}
\end{center}
\legende{Comparaison du PIB par habitant (en haut) et de l'IDH (en bas) pour quatre pays (ordres de grandeur, sources PNUD et Banque mondiale, données récentes). Le Cameroun et le Nigeria restent dans la catégorie du développement humain moyen, tandis que la France et la Corée du Sud affichent un IDH très élevé.}
\end{popfigure}

\courssub{Le développement endogène : se développer à partir de soi-même}

Si le développement est un déploiement des potentialités internes (souvenir de l'étymologie), alors il doit d'abord partir de l'\emph{intérieur} : c'est l'idée de \cle{développement endogène} (du grec \emph{endon}, « en dedans », et \emph{genos}, « naissance » : « né de l'intérieur »).

\begin{definition}[Développement endogène]
Le \cle{développement endogène} est un développement fondé sur les ressources, les besoins et les valeurs propres d'une société, par opposition à un développement importé ou imposé de l'extérieur.
\end{definition}

Un développement endogène pour l'Afrique signifie : exploiter d'abord ses propres ressources pour répondre d'abord à ses propres besoins (nourrir ses populations avant d'exporter, transformer ses matières premières sur place), en s'appuyant sur ses propres valeurs (solidarité communautaire, savoirs locaux) plutôt qu'en copiant des modèles étrangers inadaptés. Cette idée, défendue notamment par l'économiste égyptien Samir Amin, sera au cœur du projet panafricaniste que nous étudierons dans la suite de la leçon.

\begin{savaistu}[Le miracle coréen : une leçon pour l'Afrique]
En 1960, la Corée du Sud avait un PIB par habitant comparable à celui du Cameroun : c'était alors un pays pauvre, agricole, sortant d'une guerre dévastatrice. Soixante ans plus tard, elle est l'une des premières puissances industrielles mondiales (électronique, automobile, chantiers navals). Son secret ? Un investissement massif dans l'éducation, une industrialisation volontariste et une stratégie pensée à partir de ses propres besoins. Cet exemple prouve que le « sous-développement » n'est pas une fatalité : c'est un état dont on peut sortir.
\end{savaistu}

\courssub{Exemple camerounais : la Stratégie Nationale de Développement 2020-2030}

Le Cameroun s'est doté d'une stratégie officielle : la \cle{Stratégie Nationale de Développement 2020-2030 (SND30)}, qui s'inscrit dans l'ambition affichée de faire du Cameroun un \emph{pays émergent à l'horizon 2035}. Ce plan vise notamment la transformation structurelle de l'économie (industrialisation, transformation locale des matières premières comme le cacao et le coton), l'amélioration des conditions de vie des populations (santé, éducation, emploi des jeunes) et le renforcement de la gouvernance. Il illustre concrètement les quatre dimensions du développement étudiées dans cette section : économique, sociale, politique et culturelle.

\begin{exoresolu}[Croissance ou développement ?]
\textbf{Énoncé.} Un pays africain exporte du pétrole et voit son PIB augmenter de 6 \% par an pendant dix ans. Pourtant, dans le même temps, le taux de chômage des jeunes augmente, les hôpitaux manquent de médicaments et l'espérance de vie stagne. Peut-on dire que ce pays se développe ? Justifie ta réponse en t'appuyant sur les notions du cours.
\tcblower
\textbf{Solution.} Non, ce pays connaît une \emph{croissance économique} (augmentation quantitative du PIB grâce au pétrole) mais pas un \emph{développement}. En effet, le développement est l'ensemble des transformations économiques, sociales, politiques et culturelles qui améliorent \emph{durablement} les conditions de vie de la population. Or ici, la richesse produite n'est pas redistribuée : ni la santé (hôpitaux défaillants, espérance de vie stagnante), ni l'emploi (chômage des jeunes en hausse) ne progressent. On parle de \emph{croissance sans développement}. Pour que la croissance devienne développement, il faudrait que les revenus pétroliers financent l'éducation, la santé, l'industrialisation locale et une gouvernance transparente. Cet exemple illustre la distinction capitale du cours : la croissance est une condition du développement, mais non une condition suffisante.
\end{exoresolu}

\begin{aretenir}[L'essentiel de la section]
\begin{itemize}
\item Le \cle{développement} est le passage durable d'un état inférieur à un état supérieur, dans toutes les dimensions de la vie d'une société.
\item Il possède \textbf{quatre dimensions} : économique (croissance, industrialisation, transformation structurelle), humaine et sociale (santé, éducation, IDH), politique (bonne gouvernance) et culturelle (épanouissement des valeurs propres).
\item \textbf{Croissance $\neq$ développement} : la croissance est quantitative (PIB), le développement est qualitatif (bien-être, redistribution, libertés réelles).
\item Le \cle{développement endogène} part des ressources, besoins et valeurs propres d'une société : c'est la voie que le panafricanisme proposera à l'Afrique.
\item L'exemple de la Corée du Sud montre que le sous-développement n'est pas une fatalité ; la SND30 montre que le Cameroun s'est fixé l'ambition de l'émergence en 2035.
\end{itemize}
\end{aretenir}

\coursec{Le sous-développement : concept et manifestations en Afrique}

Dans la section précédente, nous avons défini le \cle{développement} comme un processus global de transformation économique, sociale et culturelle qui améliore durablement les conditions de vie des populations. Mais ce constat s'impose aussitôt : tous les pays ne se situent pas au même niveau de ce processus. Certains, dits « pays développés », disposent d'industries puissantes, d'écoles et d'hôpitaux performants ; d'autres, dont la plupart des pays africains, accumulent les difficultés. Comment nommer et comprendre cet écart ? C'est toute la question du \cle{sous-développement}, que nous abordons maintenant.

\courssub{Définition du sous-développement}

Commençons par une intuition simple. Un élève qui n'a jamais eu accès aux livres, aux enseignants et au temps d'étude n'est pas « naturellement » moins intelligent que les autres : il est en \emph{retard} par rapport à ceux qui ont bénéficié de ces conditions, et ce retard s'explique par son histoire, non par sa nature. Il en va de même, à une autre échelle, pour les nations.

\begin{definition}[Sous-développement]
Le \cle{sous-développement} désigne l'état de retard économique, social et technique d'un pays par rapport aux pays industrialisés. Il se caractérise par la \textbf{pauvreté} généralisée, la faiblesse de l'appareil productif et la \textbf{dépendance extérieure} (alimentaire, financière, technologique).
\end{definition}

Cette définition appelle trois remarques. Premièrement, le sous-développement est une notion \emph{relative} : on est sous-développé \emph{par rapport à} un modèle, celui des pays industrialisés. Deuxièmement, il ne concerne pas seulement la richesse (le PIB), mais aussi la santé, l'éducation, les infrastructures : il est multidimensionnel, comme le développement lui-même. Troisièmement, le trait le plus profond du sous-développement est peut-être la \cle{dépendance} : le pays sous-développé ne décide pas seul de son destin économique ; il subit des prix, des règles et des décisions fixés ailleurs.

\courssub{Origine du concept : Truman et le « Tiers-Monde »}

Le mot « sous-développement » n'a pas toujours existé. Il entre dans le vocabulaire international le 20 janvier 1949, lors du discours d'investiture du président américain Harry \cle{Truman}. Celui-ci propose un programme d'aide technique aux pays « sous-développés » (\emph{underdeveloped areas}) : c'est le fameux « Point IV ». Pour la première fois, le monde est officiellement divisé entre pays développés et pays sous-développés, et le développement devient un objectif explicite de la politique internationale.

\begin{savaistu}[D'où vient l'expression « Tiers-Monde » ?]
En 1952, le démographe français Alfred \cle{Sauvy} forge l'expression « \cle{Tiers-Monde} » par analogie avec le Tiers-État de la Révolution française de 1789 : de même que le Tiers-État représentait la majorité du peuple mais ne comptait pour rien dans le pouvoir, de même les pays pauvres et colonisés ou décolonisés représentaient la majorité de l'humanité mais étaient exclus des décisions mondiales. Le Tiers-Monde désignait alors les pays qui n'appartenaient ni au bloc capitaliste (Premier Monde) ni au bloc socialiste (Deuxième Monde).
\end{savaistu}

Cette généalogie du vocabulaire est philosophiquement importante : les mots ne sont jamais neutres. Dire d'un pays qu'il est « sous-développé », c'est implicitement le situer \emph{en dessous} d'une norme occidentale, comme s'il n'était qu'une version inachevée de l'Europe ou de l'Amérique. C'est pourquoi on préfère aujourd'hui parler de « pays en développement » ou de « Sud global ».

\courssub{Critique du concept : le sous-développement n'est pas un état naturel}

La thèse la plus féconde pour penser le sous-développement africain vient de l'école dite « de la \cle{dépendance} ». Son intuition est la suivante : avant la colonisation, les sociétés africaines n'étaient pas « sous-développées » ; elles avaient leurs économies, leurs royaumes, leurs réseaux commerciaux (pensons aux empires du Mali et du Songhaï, aux royaumes Bamoun et Bamiléké). Le sous-développement est donc un \emph{produit de l'histoire}, non une donnée de la nature.

\begin{propriete}[Thèse des théoriciens de la dépendance]
Pour les théoriciens de la dépendance, notamment l'économiste égyptien Samir \cle{Amin} et le sociologue Andre Gunder \cle{Frank}, le sous-développement de l'Afrique est la conséquence historique de son insertion inégale dans le système capitaliste mondial. Frank parle du « \emph{développement du sous-développement} » : le même processus qui a enrichi l'Europe (traite négrière, colonisation, pillage des matières premières) a appauvri l'Afrique. Centre et périphérie se sont construits ensemble, dans une relation d'exploitation. (Thèse à connaître ; la démonstration exhaustive n'est pas exigible.)
\end{propriete}

Le raisonnement peut se reconstruire ainsi. Si le sous-développement était un état naturel ou originel, tous les pays auraient été autrefois sous-développés de la même manière, et les pays riches ne devraient leur richesse qu'à eux-mêmes. Or l'histoire montre que l'accumulation du capital européen s'est faite en grande partie grâce aux richesses arrachées aux colonies : or, caoutchouc, coton, cacao, et même travail humain avec la traite négrière. Le retard africain et l'avance européenne sont donc les deux faces d'un même processus historique. Conclusion : le sous-développement est une \emph{relation}, non une essence.

\begin{attention}[Piège à déconstruire absolument]
Il ne faut \textbf{jamais} présenter le sous-développement comme une fatalité ou comme le signe d'une infériorité naturelle des peuples africains. Ce préjugé, hérité de l'idéologie coloniale qui prétendait « civiliser » des peuples jugés inférieurs, est scientifiquement faux et moralement inacceptable. L'histoire des grandes civilisations africaines et les réussites récentes de pays autrefois très pauvres (Corée du Sud, Singapour, Rwanda) prouvent que le sous-développement est un état \emph{historique et réversible}, non une malédiction.
\end{attention}

\courssub{Les manifestations du sous-développement en Afrique}

Comment reconnaît-on concrètement le sous-développement ? À travers des indicateurs observables dans la vie quotidienne :

\begin{itemize}
\item la \textbf{pauvreté} : une large part de la population vit avec moins de deux dollars par jour ; l'accès à l'eau potable, à l'électricité et aux soins reste inégal ;
\item l'\textbf{analphabétisme} : des millions d'enfants ne sont pas scolarisés, et beaucoup d'adultes ne savent ni lire ni écrire, ce qui freine toutes les autres transformations ;
\item la \textbf{faible industrialisation} : l'économie repose sur l'agriculture vivrière et l'exportation de matières premières brutes, sans usines de transformation ;
\item la \textbf{dépendance alimentaire} : des pays qui pourraient nourrir leur population importent du riz, du blé, du lait en poudre ;
\item la \textbf{dépendance financière} : recours massif à l'endettement extérieur et à l'aide internationale ;
\item l'\textbf{exode rural} et la \textbf{fuite des cerveaux} : les jeunes quittent les campagnes pour des villes surpeuplées, et les médecins, ingénieurs et enseignants formés en Afrique partent exercer en Europe ou en Amérique, privant le continent de ses compétences.
\end{itemize}

\begin{exemplebox}[Un paradoxe révélateur]
L'Afrique détient environ \textbf{30 \% des réserves mondiales de minerais} (or, diamant, cobalt, coltan, uranium, pétrole...), mais elle ne capte qu'une \textbf{faible part de la valeur ajoutée mondiale}, estimée à quelques pourcents seulement. Autrement dit : le continent possède les richesses du sous-sol, mais ce sont d'autres qui les transforment, les commercialisent et en tirent l'essentiel des profits. La richesse potentielle ne devient richesse réelle que si elle est transformée sur place.
\end{exemplebox}

\courssub{Les causes du sous-développement}

\textbf{Causes externes.} La première est la \textbf{colonisation} : pendant près d'un siècle, les économies africaines ont été organisées pour servir les métropoles — on produisait ce dont l'Europe avait besoin (cacao, café, coton), pas ce dont les Africains avaient besoin. Les frontières, les ports et les chemins de fer ont été tracés pour évacuer les matières premières vers la mer, non pour relier les Africains entre eux. Viennent ensuite les \textbf{échanges inégaux} : l'Afrique vend des produits bruts dont les prix, fixés sur les marchés de Londres ou de New York, baissent durablement, tandis qu'elle achète des produits manufacturés dont les prix montent. S'y ajoutent l'\textbf{endettement} (les intérêts de la dette absorbent l'argent qui pourrait financer écoles et hôpitaux) et les \textbf{règles du commerce mondial}, souvent défavorables aux pays pauvres (subventions agricoles occidentales qui écrasent les producteurs africains, par exemple).

\textbf{Causes internes.} Il serait malhonnête de tout attribuer à l'extérieur. La \textbf{mauvaise gouvernance} et la \textbf{corruption} détournent les ressources publiques ; les \textbf{conflits armés} détruisent les infrastructures et font fuir les investisseurs ; l'instabilité politique empêche toute planification à long terme ; enfin, la \textbf{faible transformation locale des matières premières} maintient les économies africaines au bas de l'échelle de la valeur.

\begin{exoresolu}[Analyser une situation concrète : le cacao et le pétrole du Cameroun]
\textbf{Énoncé.} Le Cameroun exporte son cacao brut vers l'Europe, où il est transformé en chocolat, puis réimporte du chocolat fini vendu beaucoup plus cher. De même, il exporte son pétrole brut et importe de l'essence raffinée. À l'aide des notions de la leçon, expliquez en quoi cette situation illustre le sous-développement, et proposez une solution.
\tcblower
\textbf{Solution détaillée.}
\begin{enumerate}
\item \textbf{Identifier le mécanisme.} Le Cameroun vend des matières premières \emph{brutes} à bas prix et achète des produits \emph{transformés} à prix élevé : c'est la structure typique des \emph{échanges inégaux}.
\item \textbf{Appliquer la notion de valeur ajoutée.} La plus grande partie du profit (la \emph{valeur ajoutée}) se réalise au moment de la transformation : torréfaction, fabrication, emballage, marque. Cette étape se fait à l'étranger ; le pays producteur ne perçoit qu'une faible part du prix final. Un kilo de cacao brut rapporte infiniment moins au planteur camerounais qu'un kilo de chocolat ne rapporte au fabricant européen.
\item \textbf{Relier à la définition du sous-développement.} Cette situation manifeste la \emph{faible industrialisation} et la \emph{dépendance extérieure} : l'économie camerounaise subit des prix fixés ailleurs et ne contrôle pas la chaîne de valeur de ses propres richesses.
\item \textbf{Proposer une solution.} Développer la \emph{transformation locale} : construire des usines de chocolat et des raffineries au Cameroun, former les techniciens nécessaires, afin de vendre des produits finis et de conserver la valeur ajoutée sur le territoire national. C'est exactement la voie que préconise le panafricanisme économique, que nous étudierons dans la suite de la leçon.
\end{enumerate}
\end{exoresolu}

\courssub{Illustrer le sous-développement : flux et chiffres}

\begin{popfigure}
\begin{center}
\begin{tikzpicture}[scale=1]
% Europe (schématique)
\draw[fill=popblueL,draw=popblue] (-2.2,3.2) -- (0.4,3.6) -- (1.4,2.9) -- (0.6,2.3) -- (-1.4,2.4) -- cycle;
\node[font=\small\bfseries,color=popdark] at (-0.4,3.0) {Europe};
% Asie (schématique)
\draw[fill=popgoldL,draw=popgold] (4.6,3.4) -- (6.8,3.2) -- (7.2,1.6) -- (5.4,1.2) -- (4.4,2.2) -- cycle;
\node[font=\small\bfseries,color=popdark] at (5.8,2.4) {Asie};
% Afrique (contour simplifié)
\draw[fill=popgreenL,draw=popgreen!70!black,line width=0.8pt]
  (0.2,2.0) -- (1.6,2.2) -- (2.6,1.8) -- (3.4,1.0) -- (3.6,0.0)
  -- (3.0,-1.2) -- (2.2,-2.2) -- (1.6,-2.6) -- (1.0,-1.6)
  -- (0.4,-0.6) -- (-0.6,0.2) -- (-0.8,1.2) -- cycle;
\node[font=\small\bfseries,color=popdark] at (1.5,0.2) {Afrique};
% Points villes / ressources
\fill[popdark] (0.6,0.6) circle (1.5pt); \node[font=\scriptsize,anchor=east] at (0.6,0.6) {Cacao};
\fill[popdark] (2.6,0.6) circle (1.5pt); \node[font=\scriptsize,anchor=west] at (2.6,0.6) {Pétrole};
\fill[popdark] (1.8,-1.4) circle (1.5pt); \node[font=\scriptsize,anchor=west] at (1.8,-1.4) {Minerais};
% Flux sortants (épais)
\draw[-{Stealth[length=3mm]},line width=2.2pt,color=poporange] (0.4,1.6) .. controls (-0.8,2.6) .. (-0.6,2.7);
\draw[-{Stealth[length=3mm]},line width=2.2pt,color=poporange] (3.2,1.4) .. controls (4.2,2.2) .. (4.9,2.3);
% Flux entrants (fins)
\draw[-{Stealth[length=2.5mm]},line width=0.9pt,color=popblue,dashed] (0.2,2.5) .. controls (1.0,3.0) .. (1.8,2.3);
\draw[-{Stealth[length=2.5mm]},line width=0.9pt,color=popgold,dashed] (5.0,1.5) .. controls (4.4,0.8) .. (3.7,0.5);
% Nord
\draw[-{Stealth},color=popmuted] (-2.6,-1.6) -- (-2.6,-0.8); \node[font=\scriptsize,color=popmuted,below] at (-2.6,-1.6) {N};
\end{tikzpicture}
\end{center}
\legende{Carte schématique des échanges inégaux. Les flèches épaisses orange représentent l'exportation massive de matières premières brutes (cacao, pétrole, minerais) de l'Afrique vers l'Europe et l'Asie ; les flèches fines en pointillés représentent la réimportation, en plus faible volume mais à prix élevé, de produits manufacturés finis. Croquis simplifié, sans prétention à l'exactitude géographique.}
\end{popfigure}

\begin{popfigure}
\begin{center}
\begin{tikzpicture}
\begin{axis}[
  width=\linewidth, height=7cm,
  xlabel={Année}, ylabel={PIB par habitant (dollars constants)},
  xmin=1960, xmax=2020, ymin=0, ymax=11000,
  xtick={1960,1970,1980,1990,2000,2010,2020},
  legend style={at={(0.03,0.97)},anchor=north west,font=\small},
  grid=both, grid style={popline!40},
]
\addplot[popcurve,color=popblue,line width=1.4pt] coordinates {
(1960,900) (1970,1250) (1980,1500) (1990,1350) (2000,1250) (2010,1650) (2020,1600)};
\addlegendentry{Afrique subsaharienne}
\addplot[popcurve,color=poporange,line width=1.4pt] coordinates {
(1960,450) (1970,800) (1980,1400) (1990,2600) (2000,4500) (2010,7500) (2020,10500)};
\addlegendentry{Asie de l'Est et Pacifique}
\end{axis}
\end{tikzpicture}
\end{center}
\legende{Évolution comparée du PIB par habitant (ordres de grandeur, dollars constants) entre l'Afrique subsaharienne et l'Asie de l'Est, 1960-2020. En 1960, les deux régions étaient proches ; soixante ans plus tard, l'écart est immense. Ce graphique prouve que le sous-développement n'est pas une fatalité : des pays asiatiques aussi pauvres que l'Afrique en 1960 (Corée du Sud, Chine) ont réussi leur décollage industriel. Source : ordres de grandeur inspirés des données de la Banque mondiale.}
\end{popfigure}

Ce graphique est philosophiquement décisif : il réfute par les faits toute explication « naturelle » du sous-développement. Si le retard africain tenait à la nature des peuples ou des continents, l'Asie de l'Est, partie d'un niveau comparable voire inférieur en 1960, n'aurait pas pu le surmonter. La différence tient à des choix historiques et politiques : industrialisation, transformation locale, éducation massive, stabilité de l'État. Ce qui a été fait ailleurs peut être fait en Afrique.

\begin{aretenir}[L'essentiel de la section]
\begin{itemize}
\item Le \cle{sous-développement} est un état de retard économique, social et technique, marqué par la pauvreté et la dépendance extérieure ; le terme date du discours de Truman (1949).
\item Pour Samir Amin et A. G. Frank, il n'est pas un état naturel mais le produit historique de l'insertion inégale de l'Afrique dans le capitalisme mondial (« développement du sous-développement »).
\item Ses manifestations : pauvreté, analphabétisme, faible industrialisation, dépendances alimentaire et financière, fuite des cerveaux.
\item Ses causes sont à la fois \textbf{externes} (colonisation, échanges inégaux, dette) et \textbf{internes} (gouvernance, corruption, conflits, absence de transformation locale).
\item Le sous-développement n'est pas une fatalité : l'exemple asiatique le prouve.
\end{itemize}
\end{aretenir}

\begin{exercice}[Pour s'entraîner]
\textbf{Question de synthèse.} « Le sous-développement de l'Afrique est-il une fatalité ? » Rédigez un paragraphe argumenté (une quinzaine de lignes) en mobilisant : la définition du sous-développement, la thèse des théoriciens de la dépendance, et l'exemple comparatif de l'Asie de l'Est. Vous conclurez en montrant pourquoi cette question ouvre sur celle du panafricanisme comme philosophie du développement.
\end{exercice}

\begin{corrige}[Exercice]
Le sous-développement se définit comme un retard économique, social et technique d'un pays par rapport aux pays industrialisés, caractérisé par la pauvreté et la dépendance. Le présenter comme une fatalité reviendrait à en faire un état naturel, voire à accréditer le préjugé d'une infériorité des peuples africains. Or la thèse des théoriciens de la dépendance (S. Amin, A. G. Frank) montre qu'il s'agit d'un processus historique : c'est l'insertion inégale de l'Afrique dans le système capitaliste mondial — traite, colonisation, échanges inégaux — qui a produit ce retard. Ce qui est historique est réversible. L'exemple de l'Asie de l'Est le confirme : aussi pauvre que l'Afrique en 1960, elle a connu un décollage spectaculaire grâce à l'industrialisation et à l'éducation. Le sous-développement n'est donc pas une malédiction mais un problème à résoudre. Reste à savoir par quels moyens : c'est précisément la réponse que le panafricanisme, comme philosophie du développement de l'Afrique, se propose d'apporter, et que nous examinerons dans les sections suivantes.
\end{corrige}

\coursec{Historique du panafricanisme}

Nous venons de voir que le sous-développement de l'Afrique n'est pas une fatalité naturelle : il s'enracine dans une histoire faite de traite négrière, de colonisation et de dépendances économiques. Mais cette histoire n'a pas été subie passivement. Dès le début du XX\textsuperscript{e} siècle, des penseurs et des militants noirs, d'abord dans la diaspora puis sur le continent africain, ont élaboré une réponse : unir l'Afrique et ses fils dispersés pour la libérer et la développer. C'est ce mouvement, appelé \cle{panafricanisme}, dont nous allons maintenant retracer l'histoire. Comprendre cette histoire, c'est comprendre que le panafricanisme n'est pas une simple idée abstraite : c'est une philosophie née de la souffrance, nourrie par la résistance, et progressivement transformée en projet politique et économique.

\courssub{Définition du panafricanisme}

Commençons par une intuition simple. Imaginez une famille dont les membres ont été dispersés de force dans plusieurs maisons étrangères, humiliés, et dont la maison d'origine a été pillée. Chacun, isolément, est faible. Mais s'ils se retrouvent, se reconnaissent comme frères et s'unissent, ils deviennent capables de reconstruire la maison familiale et de se faire respecter. Le panafricanisme procède exactement de cette intuition : les Africains du continent et les Noirs de la diaspora (Amériques, Caraïbes, Europe) forment une seule famille que l'histoire a dispersée, et leur force réside dans leur unité.

\begin{definition}[Panafricanisme]
Le \cle{panafricanisme} est un mouvement de pensée et d'action politique qui promeut la \textbf{solidarité} et l'\textbf{unité} des peuples africains et d'ascendance africaine (la diaspora) en vue de leur \textbf{libération} et de leur \textbf{développement}. Il vise l'unité politique, économique et culturelle de l'Afrique et de sa diaspora.
\end{definition}

Ce mot se construit sur le grec \emph{pan} (« tout ») : le panafricanisme veut rassembler \emph{toute} l'Afrique, sans exclure les descendants d'Africains déportés par la traite négrière. Il repose sur trois convictions fondamentales :

\begin{enumerate}
\item Les Africains et les Afro-descendants partagent une \textbf{communauté de destin} : l'esclavage et la colonisation les ont frappés ensemble, ils doivent s'en libérer ensemble.
\item Aucun État africain, pris isolément, n'a la taille suffisante pour peser dans le monde : seule l'\textbf{unité} peut donner à l'Afrique une voix et une puissance.
\item Le développement de l'Afrique suppose d'abord sa \textbf{libération politique} et sa \textbf{réhabilitation culturelle} : un peuple aliéné ne peut pas se développer.
\end{enumerate}

\courssub{Les origines dans la diaspora : de la résistance à la négritude}

Fait remarquable : le panafricanisme n'est pas né en Afrique, mais dans la \textbf{diaspora}, c'est-à-dire chez les descendants d'esclaves africains déportés en Amérique et aux Caraïbes. Pourquoi ? Parce que c'est là que l'humiliation de la condition noire était la plus vive, et là aussi que les Noirs ont eu les premiers accès à l'éducation occidentale et à la liberté d'expression.

La première racine du panafricanisme est donc la \textbf{résistance à l'esclavage} : révoltes d'esclaves, abolitionnisme, et au début du XX\textsuperscript{e} siècle, premiers mouvements organisés de fierté noire.

\textbf{Marcus Garvey} (1887-1940), militant jamaïcain installé aux États-Unis, fonde en 1914 l'\emph{Universal Negro Improvement Association} (UNIA). Il prêche la fierté de la race noire, l'autosuffisance économique des Noirs et le « \emph{Back to Africa} » (le retour en Afrique) : les Noirs d'Amérique doivent regarder vers l'Afrique comme leur patrie et contribuer à sa libération. Son slogan, « \emph{Africa for the Africans} » (l'Afrique aux Africains), aura un retentissement immense.

\textbf{William Edward Burghardt Du Bois} (1868-1963), sociologue et historien américain, premier Noir docteur de l'université Harvard (1895), est le véritable \emph{organisateur} du panafricanisme naissant. Il est à l'origine des premiers congrès panafricains et finira sa vie au Ghana, invité par Nkrumah. Pour lui, le problème du XX\textsuperscript{e} siècle est « le problème de la ligne de couleur » : la domination blanche sur les peuples noirs ne tombera que par l'organisation solidaire des opprimés.

Sur le plan \textbf{culturel}, le panafricanisme s'incarne dans la \textbf{négritude}, mouvement né dans le Paris des années 1930 autour d'étudiants noirs francophones.

\begin{definition}[Négritude]
La \cle{négritude} est un mouvement littéraire et culturel fondé par \textbf{Aimé Césaire} (Martiniquais), \textbf{Léopold Sédar Senghor} (Sénégalais) et \textbf{Léon-Gontran Damas} (Guyanais), qui affirme la \textbf{valeur des cultures noires} face à l'aliénation coloniale. Il s'agit de réhabiliter l'identité noire, méprisée par la colonisation, et de proclamer la fierté d'être Noir.
\end{definition}

Césaire, dans son \emph{Cahier d'un retour au pays natal} (1939), invente le mot « négritude » et dénonce la colonisation comme machine à déshumaniser. Senghor, poète et futur président du Sénégal, théorise la négritude comme « l'ensemble des valeurs culturelles du monde noir ». La négritude est ainsi le \emph{versant culturel} du panafricanisme : avant de s'unir politiquement, les Noirs doivent cesser de se mépriser eux-mêmes et retrouver la fierté de leur identité.

\begin{attention}[Ne pas confondre négritude et panafricanisme]
La \textbf{négritude} est un mouvement \emph{culturel et littéraire} (réhabilitation de l'identité noire) ; le \textbf{panafricanisme} est un mouvement \emph{politique et économique} (unité et libération de l'Afrique). Les deux sont liés — Senghor est à la fois poète de la négritude et homme d'État panafricaniste — mais ils ne se recouvrent pas. Dans une dissertation, précise toujours de quel versant tu parles.
\end{attention}

\courssub{Les congrès panafricains : Paris (1919) et Manchester (1945)}

Le panafricanisme devient un mouvement organisé grâce à une série de \textbf{congrès} qui réunissent délégués africains et afro-descendants.

Le \textbf{premier congrès panafricain} se tient à \textbf{Paris en 1919}, en marge de la conférence de paix qui suit la Première Guerre mondiale. Organisé par Du Bois avec l'appui du député sénégalais \textbf{Blaise Diagne}, il rassemble une cinquantaine de délégués. Ses revendications restent modestes : il demande que les anciennes colonies allemandes (dont le Cameroun) soient administrées dans l'intérêt des populations, et réclame des droits pour les Noirs. Mais l'essentiel est ailleurs : pour la première fois, des Noirs des trois continents se parlent et se reconnaissent solidaires.

D'autres congrès suivent (Londres 1921 et 1923, Lisbonne 1923, New York 1927), mais le véritable tournant est le \textbf{cinquième congrès de Manchester, en octobre 1945}. Le contexte a changé : la Seconde Guerre mondiale vient de s'achever, des soldats africains ont combattu pour la liberté de l'Europe, et la Charte des Nations unies proclame le droit des peuples à disposer d'eux-mêmes. Manchester marque trois ruptures :

\begin{enumerate}
\item \textbf{Le passage de relais} : la direction du mouvement passe de la diaspora aux Africains eux-mêmes. Y participent de futurs chefs d'État : \textbf{Kwame Nkrumah} (futur président du Ghana), \textbf{Jomo Kenyatta} (futur président du Kenya).
\item \textbf{La radicalisation des revendications} : on ne demande plus des réformes, on exige l'\textbf{indépendance immédiate} des colonies.
\item \textbf{L'appel à l'action de masse} : le congrès recommande les grèves, les boycotts et la désobéissance civile, sur le modèle de Gandhi.
\end{enumerate}

Manchester 1945 transforme ainsi le panafricanisme : d'un cercle d'intellectuels, il devient un \textbf{programme de lutte pour l'indépendance}.

\courssub{L'indépendance du Ghana (1957) et Kwame Nkrumah}

La première victoire concrète du panafricanisme est l'\textbf{indépendance du Ghana, le 6 mars 1957}, première colonie d'Afrique noire (au sud du Sahara) à accéder à la souveraineté. Son artisan, \textbf{Kwame Nkrumah} (1909-1972), devient aussitôt la figure majeure du panafricanisme militant.

Pour Nkrumah, l'indépendance du Ghana n'a de sens que si elle sert toute l'Afrique : « Notre indépendance est vide de sens si elle n'est pas liée à la libération totale de l'Afrique. » Il formule sa doctrine dans un ouvrage célèbre, \emph{L'Afrique doit s'unir} (1963). Son raisonnement est le suivant : les États africains, artificiellement découpés par la colonisation, sont trop petits et trop pauvres pour être réellement indépendants ; divisés, ils resteront la proie du \textbf{néocolonialisme} ; seuls les « États-Unis d'Afrique » pourront négocier d'égal à égal avec les puissances mondiales. Sa formule restée célèbre : « Cherchez d'abord le royaume politique, et tout le reste vous sera donné en plus. »

Dès 1958, Nkrumah organise à \textbf{Accra} la première Conférence des États africains indépendants, puis la Conférence des peuples africains. Le Ghana devient la capitale morale du panafricanisme : Du Bois lui-même s'y installe et y meurt en 1963.

\begin{savaistu}[Le Cameroun et le combat panafricaniste]
Le Cameroun n'est pas en reste. \textbf{Ruben Um Nyobé} (1913-1958), secrétaire général de l'\textbf{UPC} (Union des populations du Cameroun), porte devant les Nations unies, dès 1952, la revendication de l'indépendance et de la réunification du Cameroun. Figure du nationalisme camerounais, il meurt au combat en 1958. Le Cameroun indépendant adhère ensuite à l'OUA dès 1963 et accueille à Yaoundé, en 1996, le sommet qui prépare la transition vers l'Union africaine.
\end{savaistu}

\courssub{La création de l'OUA : Addis-Abeba, 25 mai 1963}

Après les indépendances (année 1960 : dix-sept États africains deviennent souverains), l'unité africaine peut enfin s'organiser. Mais deux visions s'affrontent :

\begin{itemize}
\item Le \textbf{groupe de Casablanca} (Ghana, Guinée, Mali, Égypte, Maroc, Algérie), fidèle à Nkrumah, veut une \emph{fédération} immédiate, avec un gouvernement africain.
\item Le \textbf{groupe de Monrovia} (majoritaire, incluant la plupart des États francophones), plus prudent, veut une simple \emph{coopération} entre États souverains, sans renoncer à leurs frontières.
\end{itemize}

Le compromis est trouvé à \textbf{Addis-Abeba (Éthiopie) le 25 mai 1963} : trente-deux chefs d'État signent la Charte de l'\textbf{Organisation de l'Unité Africaine} (\cle{OUA}). C'est la vision prudente de Monrovia qui l'emporte. Les objectifs de l'OUA sont :

\begin{enumerate}
\item Promouvoir l'unité et la solidarité des États africains ;
\item Coordonner la coopération politique, économique et culturelle ;
\item Défendre la souveraineté et l'intégrité territoriale des États membres ;
\item Éradiquer toutes les formes de colonialisme (lutte pour les peuples encore colonisés et contre l'apartheid en Afrique du Sud).
\end{enumerate}

Mais l'OUA repose sur deux principes qui limiteront son efficacité : la \textbf{non-ingérence} dans les affaires intérieures des États membres, et l'\textbf{intangibilité des frontières} héritées de la colonisation. Ces principes protègent les jeunes États, mais ils empêchent l'organisation d'intervenir face aux dictatures et aux conflits internes. L'OUA réussit son combat anticolonial (dernières indépendances, fin de l'apartheid en 1994), mais échoue largement sur le terrain économique.

\begin{savaistu}[La Journée de l'Afrique]
Le \textbf{25 mai}, date de la signature de la Charte de l'OUA en 1963, est célébré chaque année comme la \textbf{Journée de l'Afrique} (\emph{Africa Day}). C'est un jour férié dans plusieurs pays africains, symbole de l'unité du continent.
\end{savaistu}

\courssub{De l'OUA à l'Union africaine (2002) : l'élargissement des ambitions}

À la fin du XX\textsuperscript{e} siècle, l'OUA apparaît dépassée : elle n'a pas empêché les guerres civiles ni sorti le continent de la pauvreté. Sous l'impulsion notamment du Libyen Mouammar Kadhafi et du Sud-Africain Thabo Mbeki, les chefs d'État décident sa transformation : le 9 juillet \textbf{2002}, à Durban (Afrique du Sud), l'OUA devient l'\textbf{Union africaine} (\cle{UA}), sur le modèle de l'Union européenne.

L'UA élargit considérablement les ambitions du panafricanisme :

\begin{itemize}
\item \textbf{Intégration politique} : Parlement panafricain, Cour africaine de justice, Conseil de paix et de sécurité ; l'UA abandonne la non-ingérence absolue et peut intervenir en cas de crimes de guerre ou de génocide.
\item \textbf{Intégration économique} : le programme \textbf{NEPAD} (Nouveau Partenariat pour le Développement de l'Afrique, 2001) vise à mobiliser les ressources africaines et les partenariats internationaux pour le développement.
\item \textbf{Marché commun} : la \textbf{ZLECAf} (Zone de libre-échange continentale africaine), entrée en vigueur le 1\textsuperscript{er} janvier \textbf{2021}, crée le plus grand marché de libre-échange du monde par le nombre de pays : plus d'un milliard trois cents millions de consommateurs. L'objectif est de développer les échanges \emph{entre} pays africains, aujourd'hui très faibles (environ 15 \% du commerce total du continent).
\end{itemize}

\begin{attention}[OUA et UA : ne pas confondre]
L'\textbf{OUA} (1963) est une organisation \emph{interétatique prudente} : coopération entre États souverains, non-ingérence stricte, priorité à la décolonisation. L'\textbf{UA} (2002) affiche une ambition d'\emph{intégration plus poussée} : institutions supranationales (Parlement, Cour), droit d'intervention humanitaire, intégration économique (NEPAD, ZLECAf). Confondre les deux dates ou les deux logiques est une erreur classique à l'examen.
\end{attention}

\courssub{Frise chronologique et repères géographiques}

\begin{popfigure}
\begin{center}
\begin{tikzpicture}[x=1.55cm]
\draw[-{Stealth},thick,popdark] (0,0) -- (6.6,0);
\foreach \x/\date/\lieu/\act in {0.3/{1919}/{Paris}/{1\textsuperscript{er} congrès (Du Bois)}, 1.6/{1945}/{Manchester}/{tournant (Nkrumah, Kenyatta)}, 2.9/{1957}/{Ghana}/{indépendance (Nkrumah)}, 4.2/{1963}/{Addis-Abeba}/{création de l'OUA}, 5.5/{2002}/{Durban}/{Union africaine}} {
  \fill[popblue] (\x,0) circle (2.5pt);
  \node[above,font=\bfseries\small,text=popblue] at (\x,0.05) {\date};
  \node[below,font=\footnotesize,text=popdark,align=center] at (\x,-0.08) {\lieu\\\act};
}
\fill[poporange] (6.3,0) circle (2.5pt);
\node[above,font=\bfseries\small,text=poporange] at (6.3,0.05) {2021};
\node[below,font=\footnotesize,text=popdark,align=center] at (6.3,-0.08) {ZLECAf\\libre-échange};
\end{tikzpicture}
\end{center}
\legende{Frise chronologique des grandes étapes du panafricanisme : de la diaspora intellectuelle (1919) à l'intégration économique continentale (2021).}
\end{popfigure}

\begin{popfigure}
\begin{center}
\begin{tikzpicture}[scale=0.62]
% Contour très simplifié de l'Afrique
\draw[thick,popdark,fill=popcream] 
  (-3.2,3.4) -- (-1.2,4.1) -- (0.8,4.0) -- (2.6,3.2) -- (3.4,2.2) -- (4.2,1.2) -- (4.6,0.2) -- (3.6,-0.6) -- (3.2,-1.8) -- (2.2,-3.2) -- (1.0,-4.4) -- (0.0,-4.7) -- (-1.0,-3.8) -- (-1.8,-2.4) -- (-2.2,-1.0) -- (-3.0,0.2) -- (-3.6,1.2) -- (-3.8,2.2) -- cycle;
% Villes
\fill[popblue] (-2.9,1.6) circle (3pt); \node[left,font=\footnotesize,text=popdark] at (-3.0,1.6) {\textbf{Accra} (1957)};
\fill[popblue] (3.3,0.4) circle (3pt); \node[right,font=\footnotesize,text=popdark] at (3.4,0.4) {\textbf{Addis-Abeba} (1963)};
\fill[popblue] (0.9,-4.2) circle (3pt); \node[right,font=\footnotesize,text=popdark] at (1.0,-4.2) {\textbf{Durban} (2002)};
\fill[poporange] (-1.6,0.7) circle (3pt); \node[below,font=\footnotesize,text=popdark] at (-1.6,0.55) {Yaoundé};
% Nord
\draw[-{Stealth},thick,popdark] (-4.6,3.2) -- (-4.6,4.2); \node[font=\footnotesize\bfseries] at (-4.6,4.5) {N};
\node[font=\small\itshape,text=popmuted] at (0.3,2.6) {AFRIQUE};
\end{tikzpicture}
\end{center}
\legende{Croquis très simplifié de l'Afrique localisant les lieux étudiés : Accra (indépendance du Ghana, 1957), Addis-Abeba (siège de l'OUA, 1963), Durban (naissance de l'UA, 2002). Manchester (1945) et Paris (1919), en Europe, rappellent que le mouvement est né dans la diaspora.}
\end{popfigure}

\begin{methode}[Mémoriser une chronologie en cinq dates-repères]
Pour retenir l'histoire du panafricanisme, fixe cinq dates et associe à chacune un \textbf{lieu} et un \textbf{acteur} :
\begin{enumerate}
\item \textbf{1919} — Paris — Du Bois : premier congrès panafricain ;
\item \textbf{1945} — Manchester — Nkrumah et Kenyatta : tournant vers l'indépendance ;
\item \textbf{1957} — Accra — Nkrumah : première indépendance en Afrique noire ;
\item \textbf{1963} — Addis-Abeba — les 32 chefs d'État : création de l'OUA ;
\item \textbf{2002} — Durban — les États membres : l'OUA devient l'Union africaine (ajoute 2021, la ZLECAf, pour l'étape économique).
\end{enumerate}
Récite ces cinq couples « date — lieu — acteur » à voix haute : c'est le squelette de tout paragraphe d'histoire du panafricanisme dans une dissertation.
\end{methode}

\begin{exoresolu}[Application : replacer un événement dans son contexte]
Énoncé. Un élève écrit : « La création de l'OUA en 1963 marque la naissance du panafricanisme. » Cette affirmation est-elle exacte ? Corrige-la et justifie ta réponse.
\tcblower
\textbf{Solution.} L'affirmation est \textbf{inexacte}. L'OUA est créée le 25 mai 1963 à Addis-Abeba, mais le panafricanisme est bien plus ancien : il naît dans la diaspora noire au début du XX\textsuperscript{e} siècle (Garvey, Du Bois), s'organise dès le premier congrès de Paris en 1919, se radicalise à Manchester en 1945 et remporte sa première victoire avec l'indépendance du Ghana en 1957. La création de l'OUA n'est donc pas la \emph{naissance} du panafricanisme, mais son \textbf{institutionnalisation} : le passage d'un mouvement d'idées et de luttes à une organisation permanente d'États indépendants. Formulation correcte : « La création de l'OUA en 1963 marque l'aboutissement institutionnel du panafricanisme, mouvement né un demi-siècle plus tôt dans la diaspora. »
\end{exoresolu}

\begin{aretenir}[L'essentiel de la section]
Le panafricanisme naît dans la \textbf{diaspora} (résistance à l'esclavage, Garvey, Du Bois) et trouve son expression culturelle dans la \textbf{négritude} (Césaire, Senghor, Damas). Les congrès de \textbf{Paris (1919)} puis de \textbf{Manchester (1945)} organisent le mouvement, qui triomphe avec l'indépendance du \textbf{Ghana (1957)} sous Nkrumah. L'\textbf{OUA (Addis-Abeba, 25 mai 1963)} institutionnalise l'unité africaine sur des bases prudentes (non-ingérence), avant que l'\textbf{Union africaine (2002)} n'élargisse l'ambition vers l'intégration politique et économique (NEPAD, ZLECAf en 2021). Le Cameroun participe à cette histoire, avec la figure nationaliste de Ruben Um Nyobé.
\end{aretenir}

Cette longue marche — de la diaspora humiliée aux institutions continentales — montre que le panafricanisme a toujours poursuivi un même dessein : faire de l'unité l'instrument de la libération puis du développement de l'Afrique. Reste à préciser quels sont exactement ces objectifs, et comment le panafricanisme prétend en faire une véritable philosophie du développement : c'est l'objet de la section suivante.

\coursec{Les objectifs du panafricanisme comme philosophie du développement}

Nous avons vu, dans la section précédente, que le panafricanisme est né de la rencontre entre la diaspora africaine et les élites du continent, des congrès de Manchester (1945) à la création de l'OUA en 1963, puis de l'Union africaine en 2002. Mais un mouvement ne se juge pas seulement à son histoire : il se juge à ce qu'il veut accomplir. Quels sont donc les \cle{objectifs} du panafricanisme ? Et pourquoi peut-on dire qu'il constitue une véritable \emph{philosophie du développement} de l'Afrique ? C'est la question qui guide cette section. Nous montrerons que le panafricanisme poursuit quatre objectifs complémentaires — politique, économique, culturel et social — qui reposent sur une même conviction : \textbf{l'unité est la condition du développement}.

\courssub{Une vision d'ensemble : les quatre objectifs du panafricanisme}

Commençons par une intuition simple. Un vendeur seul au marché de Mokolo négocie mal face à un gros acheteur ; dix vendeurs associés obtiennent de meilleurs prix. Le panafricanisme applique cette sagesse à l'échelle d'un continent : séparés, les 54 États africains pèsent peu ; unis, ils deviennent une force. De cette intuition découlent quatre grands objectifs.

\begin{popfigure}
\begin{center}
\begin{tikzpicture}[mindmap, grow cyclic, every node/.style={concept, text=white},
  root concept/.append style={concept color=popblue, font=\bfseries},
  level 1/.append style={level distance=3.4cm, sibling angle=90},
  level 2/.append style={level distance=2.2cm, sibling angle=45}]
\node[root concept]{Panafricanisme}
  child[concept color=poporange]{ node {Objectif politique}
    child { node[concept color=poporange!70] {Souveraineté} }
    child { node[concept color=poporange!70] {Une seule voix} }
  }
  child[concept color=popgreen]{ node {Objectif économique}
    child { node[concept color=popgreen!70] {ZLECAf} }
    child { node[concept color=popgreen!70] {Industrialisation} }
  }
  child[concept color=poppurple]{ node {Objectif culturel}
    child { node[concept color=poppurple!70] {Dignité} }
    child { node[concept color=poppurple!70] {Langues et histoire} }
  }
  child[concept color=poppink]{ node {Objectif social}
    child { node[concept color=poppink!70] {Paix} }
    child { node[concept color=poppink!70] {Éducation, bien-être} }
  };
\end{tikzpicture}
\end{center}
\legende{Carte mentale des quatre objectifs du panafricanisme : politique, économique, culturel et social. Chaque branche précise deux moyens privilégiés.}
\end{popfigure}

\begin{methode}[Analyser un objectif du panafricanisme]
Pour analyser correctement un objectif du panafricanisme (exercice fréquent au Baccalauréat technique), procéder en trois étapes :
\begin{enumerate}
\item \textbf{Identifier le domaine} concerné : politique, économique, culturel ou social.
\item \textbf{Dégager la finalité visée} : ce que l'Afrique veut obtenir (souveraineté, prospérité, dignité, paix).
\item \textbf{Repérer le moyen institutionnel mobilisé} : OUA, Union africaine, ZLECAf, Agenda 2063, organisations régionales (CEMAC, CEDEAO).
\end{enumerate}
Exemple d'application : « la création de la ZLECAf » = domaine \emph{économique} ; finalité \emph{l'intégration des marchés} ; moyen \emph{un accord continental de libre-échange}.
\end{methode}

\courssub{L'objectif politique : achever la décolonisation et parler d'une seule voix}

Le premier objectif du panafricanisme est \cle{politique}. Il faut se souvenir que, lors des premiers congrès panafricains, l'essentiel de l'Afrique était encore colonisé. L'objectif était alors clair : \textbf{libérer le continent}. Après les indépendances des années 1960, cet objectif s'est transformé sans disparaître : il s'agit désormais d'\emph{achever} la décolonisation (soutien aux peuples encore sous domination, lutte contre l'apartheid en Afrique du Sud jusqu'en 1994) et surtout de \textbf{défendre la souveraineté} des États africains face aux ingérences extérieures.

Mais l'ambition politique va plus loin : il s'agit de faire entendre \textbf{une seule voix africaine} sur la scène mondiale. Un pays comme le Cameroun, seul, pèse peu dans les négociations de l'ONU ou de l'Organisation mondiale du commerce ; un bloc de 54 États coordonnés peut faire pencher la balance. C'est le sens des positions communes africaines sur le climat, sur la réforme du Conseil de sécurité ou sur les accords commerciaux internationaux.

\begin{propriete}[Thèse de Nkrumah : le primat de l'unité politique]
Pour Kwame \cle{Nkrumah}, premier président du Ghana et théoricien majeur du panafricanisme, \textbf{l'unité politique de l'Afrique est la condition préalable de son indépendance économique réelle}. Sa formule célèbre — \emph{« Seek ye first the political kingdom »} (« Cherchez d'abord le royaume politique ») — signifie que sans puissance politique continentale (idéalement des « États-Unis d'Afrique »), les indépendances nationales restent formelles : les anciennes métropoles continuent de décider des prix, des investissements et des alliances. Autrement dit : pas de développement économique véritable sans souveraineté politique effective, et pas de souveraineté effective sans unité.
\end{propriete}

\courssub{L'objectif économique : intégrer, industrialiser, transformer}

Le deuxième objectif est \cle{économique}. L'observation de départ est accablante : l'Afrique exporte des matières premières brutes (cacao, pétrole, bois, café) et importe des produits manufacturés coûteux. Pire, les Africains commercent très peu \emph{entre eux}.

\begin{popfigure}
\begin{center}
\begin{tikzpicture}
\begin{axis}[
  ybar, bar width=1.6cm, width=0.85\linewidth, height=6cm,
  ymin=0, ymax=75,
  ylabel={Part du commerce intra-régional (en \%)},
  symbolic x coords={Afrique, Europe},
  xtick=data,
  nodes near coords, nodes near coords align={vertical},
  every node near coord/.append style={font=\bfseries},
  ymajorgrids=true, grid style={popline!40},
]
\addplot[fill=poporange, draw=popdark] coordinates {(Afrique,15)};
\addplot[fill=popblue, draw=popdark] coordinates {(Europe,60)};
\end{axis}
\end{tikzpicture}
\end{center}
\legende{Comparaison du commerce intra-régional : environ 15 \% des échanges africains se font entre pays africains, contre environ 60 \% en Europe (ordres de grandeur, sources UA / CNUCED).}
\end{popfigure}

Ce graphique illustre le problème : un continent fragmenté en petits marchés nationaux ne peut ni industrialiser ni négocier à armes égales. D'où la stratégie panafricaine de l'\cle{intégration régionale}.

\begin{definition}[Intégration régionale]
L'\cle{intégration régionale} est le \textbf{processus par lequel des États voisins mettent en commun leurs ressources et leurs marchés} pour renforcer leur développement. Elle passe par la suppression progressive des barrières douanières, la libre circulation des personnes et des biens, et des politiques communes (monnaie, infrastructures, énergie).
\end{definition}

L'aboutissement actuel de cette stratégie est la \cle{ZLECAf} (Zone de libre-échange continentale africaine), dont le démarrage des échanges a débuté en 2021.

\begin{exemplebox}[La ZLECAf en chiffres]
La ZLECAf réunit \textbf{54 pays signataires} (sur 55 membres de l'UA) et constitue un marché d'environ \textbf{1,4 milliard de consommateurs}, avec un PIB cumulé de plus de \num{3000} milliards de dollars. C'est, par le nombre de pays, la plus grande zone de libre-échange du monde. Son objectif : augmenter fortement le commerce intra-africain et encourager la \textbf{transformation locale} des matières premières.
\end{exemplebox}

\begin{exoresolu}[La ZLECAf, opportunité pour les entreprises camerounaises]
\textbf{Énoncé.} Une entreprise de Douala transforme du cacao en chocolat et du bois en meubles. Jusqu'ici, elle exporte surtout vers l'Europe et vend peu au Gabon, au Nigeria et au Tchad à cause des droits de douane. Expliquer, en vous appuyant sur les objectifs économiques du panafricanisme, ce que la ZLECAf change pour cette entreprise.
\tcblower
\textbf{Solution détaillée.} Appliquons la méthode en trois étapes.
\begin{enumerate}
\item \emph{Domaine} : économique (intégration des marchés).
\item \emph{Finalité} : élargir le marché africain et promouvoir la transformation locale des matières premières, conformément à l'objectif panafricain d'industrialisation.
\item \emph{Moyen institutionnel} : la ZLECAf, qui réduit progressivement les droits de douane entre pays africains.
\end{enumerate}
\emph{Analyse.} Grâce à la ZLECAf, les tablettes de chocolat et les meubles fabriqués à Douala peuvent être vendus au Nigeria (plus de 200 millions de consommateurs), au Gabon ou au Kenya avec des tarifs douaniers réduits, donc à des prix plus compétitifs. L'entreprise gagne des débouchés, peut augmenter sa production, embaucher davantage et investir. Surtout, le Cameroun cesse d'exporter uniquement des fèves de cacao et des grumes de bois bruts : il exporte des \textbf{produits transformés}, à plus forte valeur ajoutée. C'est exactement le passage du statut de « fournisseur de matières premières » à celui de « producteur industriel » que le panafricanisme appelle de ses vœux. \emph{Conclusion} : la ZLECAf concrétise l'objectif économique du panafricanisme en faisant de l'unité continentale un levier de développement endogène.
\end{exoresolu}

\courssub{L'objectif culturel : restaurer la dignité africaine}

Le troisième objectif est \cle{culturel}. La colonisation n'a pas seulement exploité les richesses : elle a dévalorisé les cultures africaines, présentées comme « primitives », imposé les langues coloniales et nié l'histoire africaine (Hegel prétendait même que l'Afrique n'était « pas entrée dans l'Histoire »). Le panafricanisme, héritier de la \emph{Négritude} de Senghor et Césaire, veut \textbf{restaurer la dignité africaine} : revaloriser les langues, les arts, les philosophies et l'histoire du continent, afin que l'Africain cesse de se percevoir à travers le regard de l'autre.

Cet objectif n'est pas un luxe : un peuple qui doute de lui-même ne peut pas se développer. La confiance culturelle est le socle psychologique de l'initiative économique et politique. Enseigner l'histoire des royaumes du Sahel, des civilisations du Grand Zimbabwe ou de l'Égypte antique, promouvoir le bamiléké, l'ewondo ou le fulfuldé à côté du français et de l'anglais, c'est reconstruire l'estime de soi collective.

\courssub{L'objectif social : paix, éducation et bien-être}

Enfin, le panafricanisme vise un objectif \cle{social} : le développement n'a de sens que s'il améliore la vie concrète des populations. Cela suppose d'abord la \textbf{paix et la sécurité} — on ne construit ni école ni usine en zone de guerre — puis l'\textbf{éducation}, la santé et le bien-être. L'Union africaine s'est dotée d'une feuille de route ambitieuse pour concrétiser cet objectif.

\begin{savaistu}[L'Agenda 2063 : « l'Afrique que nous voulons »]
Adopté en 2013 pour le cinquantenaire de l'OUA, l'\cle{Agenda 2063} de l'Union africaine ambitionne de bâtir « l'Afrique que nous voulons » : une Afrique \textbf{prospère} (croissance inclusive), \textbf{intégrée} (unie politiquement et économiquement), \textbf{pacifique} (objectif « faire taire les armes »), \textbf{culturellement affirmée} et \textbf{gouvernée démocratiquement}. Il prévoit notamment un réseau de trains à grande vitesse reliant les capitales africaines, un passeport africain (en cours de déploiement) et un marché unique. C'est la traduction contemporaine, en programmes concrets, des quatre objectifs du panafricanisme.
\end{savaistu}

\courssub{Le panafricanisme comme philosophie : l'unité, condition du développement endogène}

Pourquoi parler de \emph{philosophie} et pas simplement de « programme politique » ? Parce que le panafricanisme propose une \textbf{vision du monde} cohérente, fondée sur trois principes :
\begin{itemize}
\item l'\textbf{unité} : les Africains du continent et de la diaspora forment une seule communauté de destin ;
\item la \textbf{solidarité} : le développement de chacun dépend du développement de tous ;
\item la \textbf{prise en charge de son propre destin} : l'Afrique ne doit pas attendre son salut de l'aide extérieure, mais compter d'abord sur ses propres forces.
\end{itemize}

Ce troisième principe relie directement le panafricanisme à la notion de \cle{développement endogène} : un développement qui part des ressources, des besoins et des décisions \emph{internes} au continent. L'unité en est la condition : un marché continental de 1,4 milliard de consommateurs rend l'industrialisation rentable ; un bloc uni pèse dans les négociations internationales (climat, prix des matières premières, dette). La division, au contraire, maintient chaque État dans la dépendance — c'est le mécanisme du \emph{sous-développement} étudié dans la section 2.

\begin{attention}[Pièges à éviter dans une copie]
\begin{itemize}
\item Ne pas réduire le panafricanisme à un \emph{slogan sentimental} sur la « fraternité noire » : c'est un projet politique, économique et culturel structuré, avec des institutions (OUA, UA, ZLECAf).
\item Ne pas confondre \emph{unité} et \emph{uniformité} : le panafricanisme ne veut pas effacer les diversités nationales et culturelles, mais les articuler dans une solidarité continentale.
\item Ne pas présenter les objectifs comme déjà atteints : la ZLECAf est en construction, le commerce intra-africain reste faible (environ 15 \%). Distinguer toujours l'\emph{objectif visé} et la \emph{réalité actuelle}.
\end{itemize}
\end{attention}

\begin{aretenir}[L'essentiel de la section]
\begin{itemize}
\item Le panafricanisme poursuit quatre objectifs : \textbf{politique} (souveraineté, une seule voix), \textbf{économique} (intégration, industrialisation, ZLECAf), \textbf{culturel} (dignité, revalorisation des cultures) et \textbf{social} (paix, éducation, Agenda 2063).
\item Pour Nkrumah, l'unité politique précède et conditionne l'indépendance économique réelle.
\item Le panafricanisme est une \textbf{philosophie du développement endogène} : l'unité et la solidarité sont les conditions pour que l'Afrique prenne en charge son propre destin.
\end{itemize}
\end{aretenir}

\begin{exercice}[Pour s'entraîner]
En vous appuyant sur la méthode d'analyse d'un objectif, identifier le domaine, la finalité et le moyen institutionnel des actions suivantes : (a) la position commune africaine sur la réforme du Conseil de sécurité de l'ONU ; (b) la création d'une usine de transformation de chocolat à Douala destinée au marché de la CEMAC ; (c) l'introduction de l'histoire des royaumes africains dans les programmes scolaires ; (d) l'objectif « faire taire les armes » de l'Agenda 2063.
\end{exercice}

\begin{corrige}[Exercice]
(a) Domaine \emph{politique} ; finalité : faire entendre une seule voix africaine et peser dans les décisions mondiales ; moyen : la coordination diplomatique au sein de l'Union africaine.
(b) Domaine \emph{économique} ; finalité : transformation locale des matières premières et intégration des marchés ; moyen : le marché régional (CEMAC), prolongé par la ZLECAf.
(c) Domaine \emph{culturel} ; finalité : restaurer la dignité et la confiance en soi africaines ; moyen : les politiques éducatives nationales inspirées des objectifs panafricains.
(d) Domaine \emph{social} ; finalité : la paix comme condition du bien-être et du développement ; moyen : l'Agenda 2063 de l'Union africaine et son Conseil de paix et de sécurité.
\end{corrige}

Ainsi, le panafricanisme apparaît comme un projet total : politique, économique, culturel et social. Mais entre les ambitions affichées et les réalisations concrètes, l'écart demeure grand. Reste alors la question décisive, qui fera l'objet de notre synthèse : face à ses difficultés, le panafricanisme peut-il réellement développer l'Afrique ?

\coursec{Les difficultés du panafricanisme}

Après avoir défini les objectifs ambitieux du panafricanisme — unité politique, indépendance économique, développement solidaire du continent —, une question s'impose avec force : si ce projet est si beau, pourquoi l'Afrique reste-t-elle, plus de soixante ans après les indépendances, le continent le plus pauvre et le plus fragmenté du monde ? C'est ici que la philosophie doit quitter le registre du rêve pour celui de l'analyse critique. Étudier les \cle{difficultés} du panafricanisme, ce n'est pas le condamner : c'est comprendre les obstacles réels qui se dressent entre l'idéal et sa réalisation, afin de juger si le projet est encore tenable. Nous distinguerons les difficultés \emph{internes} (politiques, économiques, culturelles, institutionnelles) des difficultés \emph{externes} (ingérences étrangères, néocolonialisme, mondialisation inégale), avant de proposer un bilan nuancé.

\courssub{Les difficultés politiques : un continent divisé}

\textbf{L'intuition.} Imaginez une grande famille qui veut construire une maison commune, mais dont chaque membre refuse de partager la clé de sa propre chambre. Telle est la situation africaine : chacun veut l'unité, mais personne ne veut renoncer à sa souveraineté.

\textbf{L'analyse.} Dès l'indépendance, les États africains se sont divisés sur le sens à donner à l'unité. En 1961, deux blocs s'affrontent :

\begin{itemize}
\item le \cle{groupe de Casablanca} (Ghana, Guinée, Mali, Maroc, Égypte, GPRA d'Algérie), conduit par Kwame \cle{Nkrumah}, qui réclame une unité politique immédiate, voire des « États-Unis d'Afrique » dotés d'un gouvernement continental ;
\item le \cle{groupe de Monrovia} (Nigeria, Libéria, Sierra Leone, Éthiopie, Togo, Cameroun, République centrafricaine, Congo-Brazzaville, Gabon, Madagascar, Mauritanie, Niger, Haute-Volta, Somalie, et ultérieurement Côte d'Ivoire, Sénégal), favorable à une simple coopération entre États souverains, sans abandon de pouvoir.
\end{itemize}

Le compromis de 1963 — la création de l'OUA — consacre la victoire du principe de Monrovia : respect des souverainetés, non-ingérence dans les affaires intérieures, intangibilité des frontières héritées de la colonisation. L'unité africaine naît donc \emph{limitée dès l'origine} par la jalousie des souverainetés nationales.

À cela s'ajoutent l'instabilité politique chronique : coups d'État à répétition (plus d'une centaine réussis ou tentés depuis 1960), guerres civiles (Nigeria-Biafra 1967-1970, Rwanda 1994), conflits interétatiques (Éthiopie-Érythrée, tensions RDC-Rwanda). Un continent en proie à ses propres conflits peine à construire un front commun.

\begin{savaistu}[1961 : l'Afrique coupée en deux]
En janvier 1961, les chefs d'État du groupe de Casablanca se réunissent au Maroc et proclament la nécessité d'une unité politique rapide du continent, selon la vision de Nkrumah : « L'Afrique doit s'unir maintenant ou périr. » En mai de la même année, le groupe de Monrovia répond en défendant une coopération prudente entre États souverains. Cette fracture fondatrice explique pourquoi l'OUA, créée en 1963 à Addis-Abeba, sera une organisation faible, fondée sur le consensus et la non-ingérence. Notez que l'Algérie, pas encore indépendante, y était représentée par le GPRA (Gouvernement provisoire de la République algérienne).
\end{savaistu}

\courssub{Les difficultés économiques : des économies rivales et dépendantes}

\textbf{L'intuition.} Deux vendeuses du marché Mokolo à Yaoundé qui vendent exactement le même manioc n'ont aucune raison de commercer entre elles : elles se font concurrence. Les économies africaines se ressemblent trop pour commercer entre elles.

\textbf{L'analyse.} Trois obstacles économiques majeurs freinent l'intégration :

\begin{enumerate}
\item \textbf{Des économies concurrentes et extraverties.} Héritières de la colonisation, les économies africaines exportent toutes des matières premières brutes (cacao, café, pétrole, cuivre) vers l'Europe, l'Amérique ou la Chine, et importent des produits manufacturés. Elles sont donc \emph{concurrentes} entre elles et \emph{extraverties}, c'est-à-dire tournées vers l'extérieur du continent.
\item \textbf{La faiblesse du commerce intra-africain.} Les routes, chemins de fer et ports ont été construits pour relier l'intérieur aux côtes (vers l'ancienne métropole), non les pays africains entre eux. Résultat : les échanges entre pays africains restent marginaux.
\item \textbf{La dépendance financière.} Quatorze pays d'Afrique centrale et de l'Ouest utilisent le \cle{Franc CFA}, monnaie arrimée à l'euro et garantie par le Trésor public français, ce qui limite leur souveraineté monétaire. S'ajoutent l'aide publique au développement, les prêts du FMI et de la Banque mondiale, souvent assortis de conditions (plans d'ajustement structurel) qui privent les États de leur liberté de décision.
\end{enumerate}

\begin{exemplebox}[Un chiffre qui parle]
Le commerce intra-africain ne représente qu'environ \textbf{15 \%} du commerce total du continent, contre plus de \textbf{60 \%} en Europe et près de \textbf{60 \%} en Asie. Autrement dit : les Africains commercent environ quatre fois plus avec le reste du monde qu'entre eux. Comment construire un marché commun quand les partenaires naturels s'ignorent ?
\end{exemplebox}

\begin{popfigure}
\begin{center}
\begin{tikzpicture}
\begin{axis}[
ybar,
bar width=1.1cm,
width=0.85\linewidth,
height=6cm,
ymin=0, ymax=75,
ylabel={Part du commerce intra-régional (en \%)},
symbolic x coords={Afrique, Asie, Europe},
xtick=data,
nodes near coords,
nodes near coords align={vertical},
every node near coord/.append style={font=\small\bfseries},
ymajorgrids=true,
grid style={popline!40},
axis line style={popdark},
]
\addplot[fill=poporange, draw=popdark] coordinates {(Afrique,15) (Asie,59) (Europe,63)};
\end{axis}
\end{tikzpicture}
\end{center}
\legende{Part approximative du commerce intra-régional dans le commerce total de chaque continent (valeurs indicatives, sources CNUCED/Commission économique pour l'Afrique). L'Afrique échange très peu entre ses propres pays, contrairement à l'Europe et à l'Asie.}
\end{popfigure}

\courssub{Les difficultés culturelles et linguistiques}

\textbf{L'intuition.} À une réunion panafricaine, le délégué camerounais de Douala parle français, son voisin nigérian parle anglais, l'Angolais parle portugais : pour se comprendre, ils doivent utiliser... les langues des anciens colonisateurs. Voilà un paradoxe douloureux.

\textbf{L'analyse.} La colonisation a légué au continent un morcellement linguistique durable : Afrique \cle{francophone}, \cle{anglophone}, \cle{lusophone}. Ces blocs linguistiques entretiennent des réseaux économiques, juridiques et diplomatiques distincts, tournés vers Paris, Londres ou Lisbonne plutôt que vers Addis-Abeba. Les sommets de l'Union africaine fonctionnent avec de coûteuses traductions simultanées, et les coopérations régionales suivent souvent les lignes linguistiques coloniales.

S'y ajoute la \cle{diversité ethnique} : l'Afrique compte plusieurs milliers de groupes ethniques. Cette richesse devient un obstacle lorsqu'elle est \emph{instrumentalisée} par des politiciens en quête de pouvoir (ethnicisme, clientélisme), comme l'ont tragiquement montré le génocide rwandais de 1994 ou les crises post-électorales en Côte d'Ivoire et au Kenya. L'identité nationale et continentale peine à primer sur les appartenances communautaires.

\courssub{Les difficultés externes : ingérences et néocolonialisme}

\textbf{L'intuition.} Un joueur de football ne peut gagner le match si l'entraîneur adverse choisit sa composition d'équipe à sa place. L'Afrique joue sur un terrain où d'autres puissances écrivent une partie des règles.

\begin{definition}[Néocolonialisme]
Le \cle{néocolonialisme} désigne la domination économique et politique \emph{indirecte} exercée par d'anciennes puissances coloniales ou de grandes puissances sur des États formellement indépendants. Contrairement au colonialisme classique, il n'a plus besoin d'occupation militaire ni d'administration directe : il passe par le contrôle des marchés, des monnaies, des dettes, des entreprises multinationales et des appuis politiques à des régimes complaisants.
\end{definition}

\textbf{L'analyse.} Les ingérences prennent plusieurs formes : soutien ou renversement de chefs d'État pendant la guerre froide (l'assassinat de Patrice Lumumba au Congo en 1961, avec l'implication de puissances étrangères, reste le symbole de la destruction d'un nationalisme africain) ; accords de défense et présence de bases militaires étrangères ; prédation sur les matières premières par les multinationales ; concurrence nouvelle de la Chine, de la Turquie ou des pays du Golfe. Enfin, la \cle{mondialisation inégale} place l'Afrique en position de fournisseur de matières brutes dont elle ne fixe pas les prix : le continent subit des règles commerciales (OMC, accords de partenariat) écrites ailleurs.

\begin{attention}[Ne pas tout imputer à l'étranger]
Une erreur fréquente en dissertation consiste à attribuer \emph{toutes} les difficultés du panafricanisme aux seules puissances étrangères, en occultant les responsabilités internes : mauvaise gouvernance, corruption, détournement des deniers publics, confiscation du pouvoir, choix économiques erronés. Un bon devoir distingue rigoureusement les causes externes (réelles) des causes internes (tout aussi réelles) et montre comment elles s'articulent : le néocolonialisme prospère souvent grâce à des complicités locales.
\end{attention}

\courssub{Les difficultés institutionnelles : une Union africaine faible}

\textbf{L'intuition.} Un arbitre sans sifflet ni carton rouge ne peut pas faire respecter les règles du jeu. L'Union africaine est souvent cet arbitre désarmé.

\textbf{L'analyse.} L'UA, succédant à l'OUA en 2002, souffre de maux structurels :

\begin{itemize}
\item \textbf{Faiblesse des moyens financiers} : une large part du budget de l'organisation (notamment pour les opérations de paix) provient de bailleurs extérieurs (Union européenne, Nations unies, États partenaires), ce qui compromet son autonomie de décision ;
\item \textbf{Non-application des décisions} : les sommets adoptent des résolutions ambitieuses, mais les États les ratifient lentement ou ne les transposent pas dans leur droit national ;
\item \textbf{Absence de contrainte} : l'UA ne dispose d'aucun pouvoir coercitif réel pour sanctionner un État qui viole ses engagements, sauf cas extrêmes (suspension après un coup d'État, appliquée de façon inconstante).
\end{itemize}

\courssub{Bilan nuancé : échecs réels, avancées réelles}

Il serait aussi faux de décréter le panafricanisme mort que de le proclamer triomphant. La lucidité philosophique exige un jugement équilibré.

\begin{methode}[Construire un bilan nuancé en dissertation]
\begin{enumerate}
\item \textbf{Lister d'abord les réussites} : cela montre que vous ne jugez pas avec partialité.
\item \textbf{Lister ensuite les échecs} : cela montre votre lucidité critique.
\item \textbf{Formuler enfin un jugement argumenté} : peser les deux plateaux de la balance, hiérarchiser, et conclure par une thèse personnelle justifiée (et non par un simple « il y a du bon et du mauvais »).
\end{enumerate}
\end{methode}

\textbf{Les avancées.} Malgré tout, le panafricanisme a enregistré des progrès notables :

\begin{itemize}
\item la \cle{ZLECAf} (Zone de libre-échange continentale africaine), entrée en vigueur en 2019 (début des échanges commerciaux en 2021), qui vise à créer le plus grand marché commun du monde par le nombre de pays et à faire passer le commerce intra-africain de 15 \% à des niveaux bien supérieurs ;
\item les \cle{forces de paix} africaines : missions de l'UA en Somalie (AMISOM/ATMIS/AUSSOM), forces régionales au Sahel (G5 Sahel, Force conjointe MNJTF au bassin du lac Tchad), qui montrent une volonté de « solutions africaines aux problèmes africains » ;
\item l'\cle{Agenda 2063}, feuille de route continentale pour une « Afrique intégrée, prospère et pacifique », qui fixe des objectifs concrets d'industrialisation, d'infrastructures et de gouvernance ;
\item la condamnation (imparfaite mais réelle) des coups d'État et la suspension des États putschistes.
\end{itemize}

\textbf{Les échecs persistants.} Pauvreté massive, dépendance alimentaire et financière, conflits non résolus, fuite des cerveaux, migrations dramatiques : les objectifs proclamés en 1963 restent largement inatteints.

\courssub{Débat philosophique : utopie ou projet réaliste ?}

\begin{popfigure}
\begin{center}
\begin{tikzpicture}[
box/.style={draw=popdark, rounded corners=2pt, align=left, font=\small, inner sep=6pt},
header/.style={draw=popdark, rounded corners=2pt, align=center, font=\small\bfseries, inner sep=6pt}
]
\node[header, fill=popblueL, minimum width=7.2cm] (h1) at (0,0) {Difficultés internes};
\node[header, fill=poporangeL, minimum width=7.2cm] (h2) at (8.2,0) {Difficultés externes};
\node[box, fill=popblueL!40, minimum width=7.2cm, anchor=north west, text width=6.8cm] at (-3.6,-0.5)
{\textbullet\ Divisions politiques (Casablanca / Monrovia, 1961)\\
\textbullet\ Souverainetés jalouses, non-ingérence\\
\textbullet\ Coups d'État, guerres civiles, conflits\\
\textbullet\ Économies concurrentes et extraverties\\
\textbullet\ Commerce intra-africain faible ($\approx$ 15 \%)\\
\textbullet\ Morcellement linguistique (franco-/anglo-/lusophone)\\
\textbullet\ Ethnicisme instrumentalisé\\
\textbullet\ Gouvernance défaillante, corruption\\
\textbullet\ UA sans moyens ni pouvoir de contrainte};
\node[box, fill=poporangeL!40, minimum width=7.2cm, anchor=north west, text width=6.8cm] at (4.6,-0.5)
{\textbullet\ Néocolonialisme (contrôle économique indirect)\\
\textbullet\ Ingérences des puissances étrangères\\
\textbullet\ Dépendance monétaire (Franc CFA)\\
\textbullet\ Dette et plans d'ajustement (FMI, BM)\\
\textbullet\ Prédation des multinationales\\
\textbullet\ Règles commerciales mondiales défavorables\\
\textbullet\ Prix des matières premières fixés ailleurs\\
\textbullet\ Concurrence des nouvelles puissances (Chine, etc.)};
\end{tikzpicture}
\end{center}
\legende{Classification des difficultés du panafricanisme : les obstacles internes relèvent des choix et de l'organisation des Africains eux-mêmes ; les obstacles externes relèvent de l'environnement international. Les deux registres interagissent.}
\end{popfigure}

Le débat philosophique peut maintenant être posé en termes rigoureux : l'unité africaine est-elle un \emph{rêve utopique} ou un \emph{projet réaliste} ?

\textbf{Thèse de l'utopie.} Les sceptiques font valoir que les obstacles sont structurels : 55 États souverains jaloux de leurs prérogatives, des milliers d'ethnies, des économies rivales, des dépendances profondes. L'unité continentale serait un mythe mobilisateur, utile pour légitimer les discours, mais irréalisable — au sens où Kant distinguait les idées régulatrices de la raison, qui orientent l'action sans jamais pouvoir être pleinement réalisées.

\textbf{Thèse du projet réaliste.} Les optimistes répondent que l'histoire dément le fatalisme : l'Union européenne est née de siècles de guerres fratricides bien plus sanglantes ; la ZLECAf prouve qu'une volonté collective peut aboutir ; la démographie (l'Afrique comptera 2,5 milliards d'habitants en 2050) fait de l'unité non un luxe mais une nécessité de survie économique. Pour Nkrumah, l'unité n'était pas un rêve mais la condition rationnelle de la liberté réelle : « Nous devons nous unir maintenant ou périr. »

\textbf{Position synthétique.} On peut soutenir, avec Sartre, que l'utopie n'est pas le contraire du réalisme mais son moteur : un projet n'est condamné que si l'on renonce à le vouloir. L'unité africaine est un \emph{projet difficile mais possible}, à condition d'en traiter les causes internes (gouvernance, intégration économique effective) autant que les causes externes (dépendances). Le panafricanisme est moins un état à atteindre qu'un \emph{processus} en marche.

\begin{aretenir}[L'essentiel de la section]
\begin{itemize}
\item Les difficultés du panafricanisme sont \textbf{internes} (divisions politiques depuis Casablanca/Monrovia en 1961, souverainetés jalouses, conflits, économies concurrentes, faible commerce intra-africain d'environ 15 \%, morcellement linguistique, faiblesse institutionnelle de l'UA) et \textbf{externes} (néocolonialisme, ingérences, mondialisation inégale).
\item Le \textbf{néocolonialisme} est une domination indirecte sur des États formellement indépendants ; il ne dispense pas d'analyser les responsabilités internes (gouvernance, corruption).
\item Le bilan est \textbf{nuancé} : échecs réels, mais avancées concrètes (ZLECAf, forces de paix, Agenda 2063).
\item Philosophiquement, l'unité africaine oscille entre \textbf{utopie} et \textbf{projet réaliste} : elle peut être pensée comme un processus exigeant plutôt qu'un rêve impossible.
\end{itemize}
\end{aretenir}

\begin{exoresolu}[Dissertation : « L'unité africaine est-elle un rêve impossible ? »]
\textbf{Énoncé.} Un élève affirme : « Puisque le panafricanisme échoue depuis soixante ans, l'unité africaine est un rêve impossible. » Construisez en quelques étapes la réfutation argumentée de cette affirmation.

\tcblower
\textbf{Solution détaillée.}
\begin{enumerate}
\item \textbf{Analyser l'argument.} L'élève commet un raisonnement par induction abusive : il conclut d'un passé d'échecs à une impossibilité définitive. Or l'échec passé ne prouve pas l'impossibilité future (sinon aucun progrès humain n'aurait jamais eu lieu).
\item \textbf{Objecter par l'histoire comparée.} L'Europe, déchirée par deux guerres mondiales, a bâti en quelques décennies une union politique et un marché unique où le commerce intra-européen dépasse 60 \%. Ce qui semblait impossible en 1945 est devenu réalité : la difficulté n'est pas l'impossibilité.
\item \textbf{Objecter par les faits africains récents.} La ZLECAf (2019/2021), les forces de paix africaines et l'Agenda 2063 montrent que le processus avance, même lentement.
\item \textbf{Nuancer.} L'unité n'est possible qu'à des conditions : bonne gouvernance interne, développement du commerce intra-africain, réduction des dépendances externes. Elle n'est ni impossible ni automatique : elle dépend de la volonté et de l'action.
\item \textbf{Conclure.} L'unité africaine n'est pas un rêve impossible, mais un projet difficile dont la réalisation est suspendue aux choix des Africains eux-mêmes. Comme le soutenait Nkrumah, l'unité est moins un rêve qu'une nécessité.
\end{enumerate}
\end{exoresolu}

\begin{exercice}[Pour s'entraîner]
\begin{enumerate}
\item Définissez le néocolonialisme et donnez deux exemples concrets de son fonctionnement en Afrique.
\item Pourquoi le commerce intra-africain est-il si faible ? Identifiez deux causes économiques et une cause historique.
\item Rédigez un paragraphe argumenté (10 à 15 lignes) répondant à la question : « Les difficultés du panafricanisme viennent-elles surtout de l'extérieur ? » Veillez à ne pas occulter les responsabilités internes.
\end{enumerate}
\end{exercice}

\begin{corrige}[Exercice]
\textbf{1.} Le néocolonialisme est la domination économique et politique indirecte exercée par d'anciennes puissances coloniales ou de grandes puissances sur des États formellement indépendants. Exemples : le contrôle monétaire via le Franc CFA garanti par le Trésor public français ; l'exploitation des matières premières par des multinationales qui rapatrient les profits ; l'appui diplomatique et militaire à des régimes complaisants.

\textbf{2.} Causes économiques : les économies africaines produisent les mêmes matières premières (elles sont concurrentes et non complémentaires) et sont extraverties, tournées vers l'exportation hors du continent. Cause historique : les infrastructures coloniales (routes, voies ferrées, ports) ont été construites pour relier l'intérieur aux côtes vers les métropoles, et non les pays africains entre eux.

\textbf{3.} Le paragraphe doit suivre la structure : thèse adverse (les causes externes sont réelles : néocolonialisme, ingérences, règles commerciales inégales) ; objection (les causes internes sont tout aussi déterminantes : divisions politiques depuis 1961, gouvernance défaillante, corruption, non-application des décisions de l'UA) ; synthèse (les deux registres s'articulent : les ingérences extérieures prospèrent souvent sur des complicités et des faiblesses internes ; traiter les unes sans les autres est illusoire).
\end{corrige}

\coursec{Synthèse : le panafricanisme peut-il développer l'Afrique ?}

Nous avons vu, dans les sections précédentes, que le panafricanisme se heurte à de sérieuses difficultés : divisions entre États, héritages coloniaux, dépendance économique, faiblesse des institutions continentales, mauvaise gouvernance. Faut-il en conclure que le projet panafricain est un échec ou une utopie ? Ce serait aller trop vite. Une difficulté n'est pas une impossibilité : elle appelle une analyse dialectique. C'est précisément l'objet de cette synthèse : confronter la thèse optimiste et l'antithèse sceptique, puis formuler une position raisonnée, avant de tirer les leçons pratiques pour la jeunesse camerounaise et pour la rédaction d'une dissertation au Baccalauréat technique.

\courssub{Rappel de la problématique}

Toute la leçon tourne autour d'une question simple mais fondamentale : \emph{l'unité africaine est-elle la voie du développement du continent ?} Autrement dit, le panafricanisme, né comme mouvement de libération politique et culturelle, peut-il devenir une véritable \cle{philosophie du développement}, capable de sortir l'Afrique du \cle{sous-développement} ?

Rappelons les termes. Le \cle{développement} est le passage d'une société d'un état de pauvreté et de dépendance à un état de progrès économique, social et humain durable ; il ne se réduit pas à la croissance du PIB, mais inclut l'éducation, la santé, la liberté et la dignité (Amartya Sen : « le développement comme liberté »). Le \cle{sous-développement} désigne l'état d'un pays dont les richesses potentielles sont réelles mais mal exploitées, en raison de blocages internes (gouvernance, infrastructures) et externes (échanges inégaux, domination néocoloniale). Le \cle{panafricanisme} est le mouvement qui affirme la solidarité de tous les Africains et des descendants d'Africains, et qui fait de l'unité du continent la condition de sa libération et de son progrès.

La problématique peut donc se formuler ainsi : \textbf{l'unité africaine est-elle une condition du développement du continent ?}

\courssub{Thèse : oui, l'unité est la voie du développement}

\begin{exemplebox}[Une intuition pour commencer]
Un seul doigt ne saisit rien ; cinq doigts unis forment un poing. De même, un seul État africain, même riche, pèse peu dans les négociations mondiales ; cinquante-cinq États unis représentent un marché de plus d'un milliard d'habitants et une voix qui compte aux Nations unies. C'est l'intuition fondatrice du panafricanisme économique.
\end{exemplebox}

La thèse panafricaine s'appuie sur trois arguments principaux.

\begin{enumerate}
\item \textbf{L'argument du marché.} Kwame Nkrumah l'affirmait dès \emph{Africa Must Unite} (1963) : les frontières coloniales ont morcelé l'Afrique en économies trop petites pour industrialiser. « L'Afrique doit s'unir » car seule l'unité crée un marché continental capable de soutenir de grandes industries. La \cle{ZLECAf} (Zone de libre-échange continentale africaine), dont l'accord est entré en vigueur en 2019 et dont les échanges ont démarré en 2021, réalise concrètement cette idée : en supprimant progressivement les droits de douane entre pays africains, elle vise à augmenter les échanges intra-africains, qui restent faibles (environ 15 \% du commerce africain, contre plus de 60 \% en Europe).
\item \textbf{L'argument du poids diplomatique.} Divisés, les États africains signent séparément des accords défavorables avec les grandes puissances et les firmes multinationales. Unis, ils négocient en position de force : prix des matières premières, siège permanent au Conseil de sécurité, conditions des accords de partenariat (ex. : position commune africaine sur le climat à la COP).
\item \textbf{L'argument de la solidarité et de la paix.} Le développement suppose la stabilité. L'Union africaine, héritière de l'OUA (1963), se donne pour mission de prévenir et de régler les conflits : « solutions africaines aux problèmes africains ». Sans paix, pas d'investissement ; sans unité, pas de paix durable.
\end{enumerate}

\begin{aretenir}[La thèse en une phrase]
Selon Nkrumah et les panafricanistes, l'unité politique et économique de l'Afrique est la \emph{condition nécessaire} du développement : elle donne au continent un marché, un poids diplomatique et une force de négociation qu'aucun État isolé ne possède.
\end{aretenir}

\courssub{Antithèse : non, ou pas encore}

\begin{exemplebox}[Une contre-observation]
Pourtant, soixante ans après les indépendances et la création de l'OUA (1963), l'Afrique reste le continent le plus pauvre : les échanges intra-africains stagnent, les conflits persistent, et les matières premières continuent de partir brutes vers l'étranger. Si l'unité était la solution, pourquoi le résultat tarde-t-il ?
\end{exemplebox}

L'antithèse sceptique avance trois contre-arguments.

\begin{enumerate}
\item \textbf{Les divisions réelles.} Le continent est traversé par des rivalités linguistiques (pays anglophones, francophones, arabophones, lusophones), des conflits frontaliers hérités de la colonisation et des nationalismes étroits. Les chefs d'État défendent souvent leur souveraineté nationale contre les institutions continentales.
\item \textbf{La dépendance extérieure.} L'économie africaine reste tournée vers l'extérieur : exportation de matières brutes, importation de produits manufacturés, endettement, aide conditionnée. Tant que cette structure extravertie n'est pas transformée, l'unité politique reste un discours sans prise sur la réalité économique.
\item \textbf{La mauvaise gouvernance.} Corruption, détournements, absence d'État de droit : beaucoup de maux africains sont d'abord internes. Un sceptique objecte que l'unité de gouvernements mal gérés ne produirait qu'une mauvaise gouvernance à plus grande échelle.
\end{enumerate}

\begin{attention}[Piège de raisonnement]
Ne confonds pas « le panafricanisme a échoué » avec « le panafricanisme est impossible ». Les difficultés constatées montrent que le projet est \emph{inachevé}, non qu'il est \emph{absurde }. Dans une dissertation, évite aussi le manichéisme : présenter la thèse comme toute blanche et l'antithèse comme toute noire empêche toute synthèse.
\end{attention}

\courssub{Synthèse : un horizon nécessaire, à condition de passer aux actes}

La démarche dialectique consiste à dépasser l'opposition : la thèse a raison sur le \emph{principe} (l'unité est nécessaire), l'antithèse a raison sur le \emph{constat} (l'unité n'est pas encore effective). La synthèse s'énonce ainsi : \textbf{le panafricanisme reste un horizon nécessaire du développement africain, à condition de réformer la gouvernance et de passer des discours aux actes.}

Concrètement, cela signifie :
\begin{itemize}
\item que l'unité n'est pas un événement mais un \emph{processus} : elle se construit par étapes (intégration régionale d'abord, continentale ensuite), comme le montre l'exemple de la construction européenne ;
\item que l'unité politique sans transformation économique est vide : il faut transformer sur place les matières premières (cacao, pétrole, bois, coton) au lieu de les exporter brutes ;
\item que l'unité exige des États démocratiques et bien gouvernés : la lutte contre la corruption est une condition du panafricanisme, non son ennemie.
\end{itemize}

\begin{popfigure}
\begin{center}
\begin{tikzpicture}[
box/.style={draw=popblue, thick, rounded corners, fill=popblueL, text width=4.6cm, align=center, minimum height=2.6cm, font=\small},
arr/.style={-{Stealth[length=3mm]}, very thick, poporange}
]
\node[box, align=center] (T) at (0,0){\textbf{THÈSE}\\ Oui : l'unité donne un marché, un poids diplomatique, une force de négociation (Nkrumah, ZLECAf)};
\node[box, fill=poppinkL, draw=poppink, align=center] (A) at (7.6,0){\textbf{ANTITHÈSE}\\ Non, pas encore : divisions, dépendance extérieure, mauvaise gouvernance};
\node[box, fill=popgreenL, draw=popgreen, text width=6.2cm] (S) at (3.8,-3.6) {\textbf{SYNTHÈSE}\\ Le panafricanisme est un horizon nécessaire, à condition de réformer la gouvernance et de passer des discours aux actes};
\draw[arr] (T.south) -- (S.north west);
\draw[arr] (A.south) -- (S.north east);
\end{tikzpicture}
\end{center}
\legende{Schéma de la démarche dialectique appliquée au sujet : la synthèse conserve la vérité de la thèse (l'unité est nécessaire) et la vérité de l'antithèse (elle n'est pas encore réalisée).}
\end{popfigure}

\courssub{Le rôle de la jeunesse et de la diaspora : l'afro-optimisme}

Le panafricanisme du XXI\textsuperscript{e} siècle ne se fait plus seulement dans les sommets des chefs d'État : il se fait aussi par la \cle{jeunesse} et la \cle{diaspora}. C'est ce qu'on appelle parfois l'\cle{afro-optimisme} : la conviction, portée par les nouvelles générations, que l'Afrique peut se développer par ses propres forces.

\begin{itemize}
\item \textbf{L'entrepreneuriat.} De jeunes Africains créent des entreprises qui répondent aux besoins locaux : paiement mobile, agriculture transformée, énergie solaire. Au Cameroun, des start-up de Douala et de Yaoundé développent des applications de livraison, de santé ou de services agricoles.
\item \textbf{La culture.} La musique (afrobeat, makossa, bikutsi), le cinéma (Nollywood et ses équivalents), la mode et la littérature créent une fierté africaine partagée et des industries culturelles créatrices d'emplois. La culture est aussi une économie.
\item \textbf{Le numérique.} Internet abolit les distances : un jeune développeur de Bafoussam peut travailler avec un client de Nairobi ou de la diaspora de Paris. La diaspora, par ses transferts d'argent (qui dépassent souvent l'aide publique au développement), ses compétences et ses réseaux, est un acteur majeur du développement.
\end{itemize}

\begin{savaistu}[Le savais-tu ?]
Les sujets de philosophie au Baccalauréat technique camerounais portent régulièrement sur le développement de l'Afrique et l'unité africaine : « L'unité africaine est-elle possible ? », « Le développement de l'Afrique passe-t-il par l'unité ? », « Le panafricanisme est-il une utopie ? ». Maîtriser cette leçon, c'est donc se préparer directement à l'épreuve.
\end{savaistu}

\courssub{Application au Cameroun}

Que signifie concrètement le panafricanisme pour un élève camerounais de Terminale technique ?

\begin{enumerate}
\item \textbf{La transformation locale.} Le Cameroun exporte son cacao brut et importe du chocolat. Transformer sur place (chocolateries, huileries, scieries) crée des emplois, conserve la valeur ajoutée au pays et réduit la dépendance : c'est le panafricanisme économique appliqué à l'échelle nationale.
\item \textbf{L'intégration régionale.} Le Cameroun est membre de la \cle{CEMAC} (Communauté économique et monétaire de l'Afrique centrale) et de la \cle{CEEAC} (Communauté économique des États de l'Afrique centrale). Ces organisations préparent l'unité continentale : libre circulation, monnaie commune (le franc CFA de la zone CEMAC émis par la BEAC), projets d'infrastructures transfrontalières (routes, corridors Douala--N'Djamena).
\item \textbf{L'éducation aux valeurs panafricaines.} L'école forme des citoyens ouverts sur le continent : apprendre l'histoire africaine, respecter les autres cultures africaines, refuser les préjugés et le tribalisme. Le panafricanisme commence dans les mentalités.
\end{enumerate}

\begin{popfigure}
\begin{center}
\begin{tikzpicture}
\draw[popline, thick] (0,0) rectangle (15.4,5.6);
\draw[fill=popblue, draw=popblue] (0,4.7) rectangle (15.4,5.6);
\node[white, font=\bfseries] at (7.7,5.15) {Concepts clés de la leçon};
\draw[fill=popblueL, draw=popblue] (0,3.8) rectangle (3.2,4.7); \node[font=\bfseries\small] at (1.6,4.25) {Concept};
\draw[fill=popblueL, draw=popblue] (3.2,3.8) rectangle (9.4,4.7); \node[font=\bfseries\small] at (6.3,4.25) {Définition};
\draw[fill=popblueL, draw=popblue] (9.4,3.8) rectangle (15.4,4.7); \node[font=\bfseries\small] at (12.4,4.25) {Exemple};
\draw[fill=popcream, draw=popline] (0,1.9) rectangle (3.2,3.8); \node[text width=2.8cm, align=center, font=\small] at (1.6,2.85) {\textbf{Développement}};
\draw[fill=popcream, draw=popline] (3.2,1.9) rectangle (9.4,3.8); \node[text width=5.8cm, align=center, font=\small] at (6.3,2.85) {Progrès économique, social et humain durable ; passage de la pauvreté et de la dépendance à l'autonomie et à la dignité};
\draw[fill=popcream, draw=popline] (9.4,1.9) rectangle (15.4,3.8); \node[text width=5.6cm, align=center, font=\small] at (12.4,2.85) {Émergence visée par le Cameroun à l'horizon 2035};
\draw[fill=white, draw=popline] (0,0.95) rectangle (3.2,1.9); \node[text width=2.8cm, align=center, font=\small] at (1.6,1.42) {\textbf{Sous-développement}};
\draw[fill=white, draw=popline] (3.2,0.95) rectangle (9.4,1.9); \node[text width=5.8cm, align=center, font=\small] at (6.3,1.42) {Richesses potentielles réelles mais mal exploitées, à cause de blocages internes et de la domination extérieure (échanges inégaux)};
\draw[fill=white, draw=popline] (9.4,0.95) rectangle (15.4,1.9); \node[text width=5.6cm, align=center, font=\small] at (12.4,1.42) {Exportation du cacao brut et importation du chocolat};
\draw[fill=popcream, draw=popline] (0,0) rectangle (3.2,0.95); \node[text width=2.8cm, align=center, font=\small] at (1.6,0.47) {\textbf{Panafricanisme}};
\draw[fill=popcream, draw=popline] (3.2,0) rectangle (9.4,0.95); \node[text width=5.8cm, align=center, font=\small] at (6.3,0.47) {Mouvement d'unité et de solidarité des Africains et de la diaspora, philosophie du développement du continent};
\draw[fill=popcream, draw=popline] (9.4,0) rectangle (15.4,0.95); \node[text width=5.6cm, align=center, font=\small] at (12.4,0.47) {OUA (1963), Union africaine (2002), ZLECAf (2021)};
\end{tikzpicture}
\end{center}
\legende{Tableau récapitulatif des trois concepts clés de la leçon, avec définitions et exemples camerounais et africains.}
\end{popfigure}

\courssub{Méthode : rédiger la dissertation sur ce sujet}

\begin{methode}[Rédiger une dissertation de philosophie]
\begin{enumerate}
\item \textbf{Introduction.} Elle comprend quatre moments : (a) l'\emph{amorce} ou accroche, qui part d'un fait général ou d'une observation concrète ; (b) la \emph{définition des termes} du sujet (ici : unité africaine, développement) ; (c) la \emph{problématique}, formulée par une question claire ; (d) l'\emph{annonce du plan}, qui présente les deux ou trois parties.
\item \textbf{Développement.} Deux ou trois parties argumentées : chaque partie défend une thèse avec des arguments, des auteurs cités brièvement et des exemples ; les parties sont reliées par des transitions (« Mais... », « Cependant... », « Ainsi... »).
\item \textbf{Conclusion.} Elle répond à la problématique par une position justifiée, et peut s'ouvrir sur une question nouvelle.
\end{enumerate}
\end{methode}

\begin{methode}[Justifier une conclusion]
\begin{enumerate}
\item \textbf{Reprendre la problématique} : rappeler la question posée dans l'introduction.
\item \textbf{Évaluer les arguments des deux parties} : dire ce que chaque thèse a de vrai et de limité (« la thèse a raison de... mais elle néglige... »).
\item \textbf{Formuler une position personnelle argumentée} : ne pas rester neutre ; trancher en donnant une raison décisive, éventuellement en dépassant l'opposition (synthèse).
\end{enumerate}
\end{methode}

\begin{attention}[La dissertation « catalogue »]
L'erreur la plus fréquente est la dissertation \emph{catalogue} : l'élève juxtapose des idées, des noms d'auteurs et des exemples sans problématique ni liens logiques. Le correcteur ne demande pas une liste, mais un \emph{raisonnement} : chaque partie doit répondre à la question posée, et chaque paragraphe doit enchaîner avec le précédent. Avant d'écrire, demande-toi toujours : « En quoi ce que je viens de dire répond-il à la problématique ? »
\end{attention}

\begin{exemplebox}[Problématique type Bac]
Sujet : \emph{« L'unité africaine est-elle une condition du développement du continent ? »}\\
Introduction possible (schéma) : \textbf{Amorce} — Soixante ans après les indépendances, l'Afrique, riche de ses matières premières, demeure le continent le plus pauvre. \textbf{Définitions} — L'unité africaine désigne la solidarité politique et économique des États du continent (panafricanisme) ; le développement est le progrès économique, social et humain durable. \textbf{Problématique} — L'unité africaine est-elle une condition du développement du continent ? \textbf{Annonce du plan} — Nous verrons d'abord que l'unité est nécessaire au développement, puis qu'elle se heurte à de réels obstacles, enfin qu'elle demeure un horizon à construire par des actes concrets.
\end{exemplebox}

\begin{exoresolu}[Construire le plan détaillé d'une dissertation]
Énoncé : À partir du sujet « Le panafricanisme peut-il développer l'Afrique ? », construis un plan dialectique en trois parties, avec pour chaque partie une thèse, deux arguments et un exemple.
\tcblower
\textbf{Solution détaillée.}
\textbf{I. Thèse : le panafricanisme peut développer l'Afrique.} Argument 1 : l'unité crée un grand marché continental (Nkrumah, \emph{Africa Must Unite} ; exemple : la ZLECAf). Argument 2 : l'unité donne un poids diplomatique dans les négociations mondiales (exemple : position commune africaine sur le climat ou les matières premières).
\textbf{II. Antithèse : le panafricanisme ne peut pas (encore) développer l'Afrique.} Argument 1 : les divisions linguistiques, frontalières et nationalistes affaiblissent l'unité (exemple : conflits frontaliers hérités de la colonisation). Argument 2 : la dépendance économique et la mauvaise gouvernance bloquent le développement, avec ou sans unité (exemple : exportation de matières brutes, corruption).
\textbf{III. Synthèse : le panafricanisme est un horizon nécessaire à condition de passer aux actes.} Argument 1 : l'unité se construit par étapes, régionales puis continentales (exemple : CEMAC, CEEAC). Argument 2 : elle exige la transformation locale des richesses, la bonne gouvernance et l'engagement de la jeunesse et de la diaspora (exemple : entrepreneuriat numérique au Cameroun).
\textbf{Conclusion} : le panafricanisme ne développe pas l'Afrique automatiquement, mais sans lui le développement du continent est impossible ; il reste donc une tâche à accomplir, non un rêve à abandonner.
\end{exoresolu}

\begin{exercice}[Entraînement au sujet de dissertation]
Rédige l'introduction complète (amorce, définitions, problématique, annonce du plan) du sujet suivant : « Le développement de l'Afrique passe-t-il nécessairement par l'unité africaine ? » Puis rédige une conclusion justifiée qui évalue les arguments des deux parties et prend position.
\end{exercice}

\begin{corrige}[Exercice : éléments de correction]
\textbf{Introduction attendue.} Amorce : constat du paradoxe africain (continent riche en ressources, population pauvre). Définitions : développement (progrès économique, social et humain durable) ; unité africaine (solidarité politique et économique des États du continent, panafricanisme). Problématique : le développement de l'Afrique passe-t-il nécessairement par l'unité africaine ? Annonce du plan : I. Oui, l'unité est nécessaire (marché, poids diplomatique) ; II. Mais des obstacles internes et externes la rendent difficile ; III. Elle demeure un horizon à construire par la réforme de la gouvernance et l'action concrète.
\textbf{Conclusion attendue.} La thèse a raison d'affirmer qu'aucun État africain isolé ne peut peser seul ; l'antithèse a raison de rappeler que l'unité reste inachevée et que la dépendance et la mauvaise gouvernance sont des freins réels. Position justifiée : l'unité n'est pas une condition suffisante, mais elle est une condition nécessaire du développement ; elle exige de passer des discours aux actes (transformation locale, intégration régionale, engagement de la jeunesse). Ouverture possible : la jeunesse africaine saura-t-elle réaliser l'unité que les générations précédentes ont rêvée ?
\end{corrige}

\begin{aretenir}[Conclusion de la leçon]
Le panafricanisme n'est ni une baguette magique ni une utopie à abandonner. La thèse (Nkrumah, ZLECAf) montre que l'unité donne à l'Afrique un marché, un poids diplomatique et une force de négociation ; l'antithèse rappelle les divisions, la dépendance et la mauvaise gouvernance. La synthèse dépasse l'opposition : le panafricanisme est un \emph{horizon nécessaire}, qui ne deviendra réalité que par la réforme de la gouvernance, la transformation locale des richesses et l'engagement de la jeunesse et de la diaspora. Au Cameroun, il se traduit par la transformation locale, l'intégration dans la CEMAC et la CEEAC, et l'éducation aux valeurs panafricaines.
\end{aretenir}

\newpage

\coursec{Activité d'intégration}

\coursec{Panafricanisme et développement de l’Afrique}

\courssub{Objectifs de l'activité}

Cette activité d'intégration mobilise l'ensemble des savoirs et savoir-faire de la leçon sur le \cle{panafricanisme} et le \cle{développement} de l'Afrique. À l'issue du travail, l'élève doit être capable de :

\begin{itemize}
\item définir et distinguer les concepts de \cle{développement} et de \cle{sous-développement} à partir d'une situation économique concrète ;
\item identifier, dans des documents chiffrés, les \cle{manifestations du sous-développement} (exportation de matières premières brutes, dépendance extérieure, faible industrialisation, échange inégal) ;
\item appliquer les notions de la leçon (\cle{panafricanisme}, \cle{intégration africaine}, \cle{transformation locale}, \cle{ZLECAf}) à une situation camerounaise authentique ;
\item construire et interpréter un diagramme en barres et un croquis de flux commerciaux ;
\item rédiger et justifier un \cle{paragraphe argumenté} selon la démarche de la dissertation philosophique (thèse, argument, exemple, objection, réponse).
\end{itemize}

\courssub{Situation}

\textbf{Document 1. La filière cacao au Cameroun.}

Le Cameroun est le quatrième producteur mondial de \cle{cacao}. La campagne 2022-2023 a produit environ \num{300000} tonnes de fèves, cultivées par près de \num{600000} petits planteurs, principalement dans les régions du Centre, du Sud, du Sud-Ouest et du Littoral. Mais environ \textbf{80 \%} de cette production est exportée \emph{brute} (fèves séchées, non transformées), principalement vers les Pays-Bas, la Belgique et l'Allemagne. Le chocolat consommé au Cameroun est, lui, en grande partie \emph{importé}.

\begin{center}
\begin{tabularx}{\linewidth}{|l|X|X|}\hline
\rowcolor{popblueL} \textbf{Produit} & \textbf{Prix indicatif (FCFA/kg)} & \textbf{Observations} \\ \hline
Cacao brut exporté (fèves) & \num{1500} & Prix payé au producteur souvent inférieur (\num{1000} à \num{1200}) \\ \hline
Chocolat importé (tablette) & \num{7500} & Soit 5 fois le prix du kilogramme de fèves \\ \hline
\end{tabularx}
\end{center}

\textbf{Document 2. Extrait du discours de Kwame Nkrumah (Accra, 1963).}

\begin{quote}
\og L'indépendance politique n'a de sens que si elle s'accompagne de l'indépendance économique. Aucune nation africaine, prise isolément, ne peut se développer seule. C'est dans l'unité que réside notre force : unissons nos marchés, nos industries et nos ressources, ou nous resterons les fournisseurs de matières premières des autres. \fg
\end{quote}

\textbf{Document 3. La Zone de libre-échange continentale africaine (ZLECAf).}

Entrée en vigueur en 2021, la \cle{ZLECAf} regroupe 54 États africains et vise à créer un marché unique de 1,4 milliard de consommateurs. Le commerce \emph{intra-africain} ne représente encore qu'environ \textbf{15 \%} du commerce total du continent (contre environ 60 \% en Europe). Les obstacles identifiés : frontières héritées de la colonisation, monnaies différentes, routes et chemins de fer insuffisants, concurrence entre économies semblables.

\textbf{Problème posé à la classe :} la coopérative des planteurs de cacao de la commune d'Obala (région du Centre) hésite entre deux stratégies : continuer à vendre les fèves brutes aux exportateurs européens, ou investir dans une unité locale de transformation (beurre de cacao, poudre, chocolat) destinée au marché camerounais et aux marchés africains voisins grâce à la ZLECAf. Les élèves, réunis en groupes de quatre, sont les conseillers de la coopérative : ils doivent analyser les documents et rédiger un avis argumenté.

\courssub{Consignes}

\begin{enumerate}
\item À partir du document 1, construis un \textbf{diagramme en barres} comparant le prix du kilogramme de cacao brut exporté et celui du kilogramme de chocolat importé. Calcule le rapport entre les deux prix et commente : qui capture la plus grande partie de la valeur ?
\item Réalise un \textbf{croquis schématique} montrant les flux : fèves brutes du Cameroun vers l'Europe, chocolat fini de l'Europe vers le Cameroun. Fais apparaître la direction des flux et la valeur ajoutée.
\item Relève dans les trois documents \textbf{trois manifestations du sous-développement} du Cameroun et justifie chaque choix à l'aide de la définition du sous-développement vue en cours.
\item À partir des documents 2 et 3, explique en quoi le \textbf{panafricanisme} (unité, intégration, marché commun) propose une réponse à ces manifestations, puis identifie \textbf{deux difficultés} concrètes qui freinent cette réponse.
\item Rédige collectivement un \textbf{paragraphe argumenté} (15 à 20 lignes) répondant à la question : \og La transformation locale du cacao et l'intégration africaine peuvent-elles sortir le Cameroun du sous-développement ? \fg{} Ton paragraphe devra contenir une thèse, deux arguments étayés par les données du dossier, une objection et une réponse à l'objection.
\item Restitution : chaque groupe présente son diagramme, son croquis et son paragraphe en 5 minutes. La classe corrige collectivement à l'aide de la grille de la dissertation (pertinence de la thèse, solidité des arguments, exactitude des exemples, clarté de la conclusion).
\end{enumerate}

\courssub{Ressources à mobiliser}

\begin{aretenir}[Notions utiles de la leçon]
\begin{itemize}
\item \textbf{Développement} : processus global de transformation d'une société (économique, social, culturel, humain) qui améliore durablement les conditions de vie des populations. Il ne se réduit pas à la croissance du PIB.
\item \textbf{Sous-développement} : état d'une économie dépendante, peu industrialisée, qui exporte des matières premières brutes à bas prix et importe des produits finis à prix élevé (\emph{échange inégal}).
\item \textbf{Panafricanisme} : philosophie et mouvement visant l'unité politique et économique de l'Afrique (Nkrumah, Senghor, Nyerere), prolongé aujourd'hui par l'Union africaine et la ZLECAf.
\item \textbf{Difficultés} : héritages coloniaux (frontières, langues, monnaies), faiblesse des infrastructures, rivalités entre États, dépendance financière extérieure.
\end{itemize}
\end{aretenir}

\begin{methode}[Rappel : la structure du paragraphe argumenté]
\begin{enumerate}
\item Annoncer la \textbf{thèse} (la réponse à la question).
\item Développer les \textbf{arguments}, chacun étayé par un exemple ou une donnée chiffrée du dossier.
\item Formuler une \textbf{objection} honnête (ce que dirait un adversaire de la thèse).
\item Répondre à l'objection et conclure en nuançant si nécessaire.
\end{enumerate}
\end{methode}

\courssub{Démarche et corrigé commenté}

\begin{exoresolu}[Corrigé guidé de l'activité]
\textbf{Énoncé :} À partir du dossier (filière cacao, discours de Nkrumah, ZLECAf), analyser les mécanismes du sous-développement, proposer une réponse panafricaniste et rédiger un paragraphe argumenté structuré.

\tcblower

\textbf{Étape 1. Le diagramme en barres.} On représente les deux prix du document 1 (échelle : 1 cm pour \num{1000} FCFA).

\begin{popfigure}
\begin{tikzpicture}
\begin{axis}[
ybar, bar width=1.6cm, width=0.75\linewidth, height=6cm,
ymin=0, ymax=9000,
ylabel={Prix (FCFA/kg)},
symbolic x coords={Cacao brut exporte, Chocolat importe},
xtick=data,
xticklabels={Cacao brut export\'e, Chocolat import\'e},
nodes near coords,
every node near coord/.append style={font=\small},
ymajorgrids=true, grid style={popline!40},
]
\addplot[fill=poporange!70, draw=poporange] coordinates {(Cacao brut exporte,1500) (Chocolat importe,7500)};
\end{axis}
\end{tikzpicture}
\legende{Comparaison des prix : cacao brut exporté vs chocolat importé (données Document 1).}
\end{popfigure}

Rapport : $7500 \div 1500 = 5$. Le kilogramme de chocolat importé coûte \textbf{5 fois} le kilogramme de fèves exportées. \emph{Interprétation} : la \cle{valeur ajoutée} (transformation, emballage, marque, logistique, marketing) est captée par les industries européennes, non par le planteur camerounais. C'est l'illustration chiffrée de l'\emph{échange inégal}.

\textbf{Étape 2. Le croquis des flux.}

\begin{popfigure}
\begin{tikzpicture}[scale=1, >=Stealth]
\node[draw=popgreen, fill=popgreenL, rounded corners, minimum width=3.2cm, minimum height=1.4cm, align=center] (cam) at (0,0) {\textbf{Cameroun}\\ f\`eves brutes};
\node[draw=popblue, fill=popblueL, rounded corners, minimum width=3.2cm, minimum height=1.4cm, align=center] (eur) at (8,0) {\textbf{Europe}\\ usines de chocolat};
\draw[->, line width=2.5pt, poporange] (cam.north) to[bend left=25] node[above, align=center, font=\small]{F\`eves brutes (\num{1500} FCFA/kg)} (eur.north);
\draw[->, line width=2.5pt, poppurple] (eur.south) to[bend left=25] node[below, align=center, font=\small]{Chocolat fini (\num{7500} FCFA/kg)} (cam.south);
\node[align=center, font=\small\itshape, popmuted] at (4,-2.2) {La valeur ajout\'ee ($\times 5$) reste en Europe : \'echange in\'egal};
\end{tikzpicture}
\legende{Schéma des flux commerciaux Cameroun--Europe : exportation de matière première brute (flèche orange) et importation de produit fini (flèche violette).}
\end{popfigure}

\textbf{Étape 3. Manifestations du sous-développement.}
\begin{enumerate}
\item \emph{Exportation de matières premières brutes / faible industrialisation} (Doc 1) : 80 \% des fèves partent sans transformation locale, signe d'un appareil productif incomplet qui ne maîtrise pas la transformation de ses propres ressources.
\item \emph{Dépendance extérieure / absence de maîtrise des prix} (Doc 1) : le prix du cacao est fixé sur les marchés étrangers (bourse de Londres/New York), le planteur ne maîtrise ni le prix ni le débouché ; le chocolat consommé localement est importé.
\item \emph{Échange inégal / appauvrissement de la balance commerciale} (Doc 1) : le Cameroun vend à \num{1500} FCFA ce qu'il rachète transformé à \num{7500} FCFA (rapport 1 à 5), ce qui structure un déficit commercial chronique et un transfert de richesses vers l'extérieur.
\end{enumerate}
Chaque manifestation correspond exactement à la définition du sous-développement comme économie dépendante, désarticulée et soumise à la division internationale du travail héritée de la colonisation.

\textbf{Étape 4. Le panafricanisme comme réponse et ses difficultés.}
\emph{Réponse panafricaniste} (Doc 2 et 3) : Nkrumah affirme qu'aucun État africain ne peut se développer seul ; l'\cle{unité} crée un marché assez vaste pour justifier l'industrialisation (économies d'échelle). La ZLECAf concrétise cette visée : un marché unique de 1,4 milliard de consommateurs permettrait au chocolat camerounais d'être vendu au Nigeria, au Ghana, en Côte d'Ivoire, rendant l'investissement industriel rentable. L'intégration permet aussi de mutualiser les infrastructures (énergie, transport) et de négocier collectivement les cours des matières premières.
\emph{Difficultés concrètes} (Doc 3 et cours) :
\begin{enumerate}
\item Faiblesse du commerce intra-africain (15 \% seulement) due aux frontières coloniales, à la multiplicité des monnaies, au déficit d'infrastructures routières/ferroviaires (ex. : pas de liaison ferroviaire Yaoundé-Douala performante, routes Obala-Douala dégradées).
\item Concurrence entre économies structurellement semblables (productrices de cacao, café, coton) qui peinent à se compléter, et dépendance financière extérieure (financement des usines, technologie) qui maintient la tutelle économique.
\end{enumerate}

\textbf{Étape 5. Paragraphe argumenté (modèle de rédaction).}

\emph{La transformation locale du cacao et l'intégration africaine constituent une voie nécessaire, quoique insuffisante à elle seule, pour sortir le Cameroun du sous-développement.} D'abord, transformer sur place permet de capter la valeur ajoutée : le dossier montre que le chocolat importé vaut cinq fois le prix des fèves exportées ; chaque tonne transformée à Obala enrichit donc le pays au lieu d'enrichir l'industriel étranger. Ensuite, l'intégration africaine, conformément à l'intuition de Nkrumah, offre le marché que le seul Cameroun ne peut offrir : grâce à la ZLECAf, le chocolat camerounais peut viser 1,4 milliard de consommateurs, ce qui rend l'investissement industriel rentable. \emph{Objection} : on dira que les obstacles (routes défaillantes, monnaies différentes, commerce intra-africain limité à 15 \%) rendent ce projet illusoire à court terme. \emph{Réponse} : ces difficultés sont réelles mais elles désignent précisément le travail politique à accomplir ; le sous-développement n'est pas une fatalité naturelle mais un état historique que des décisions souveraines — infrastructures, éducation technique, solidarité continentale, politique industrielle — peuvent transformer. Ainsi, la transformation locale unie à l'intégration africaine n'est pas une baguette magique, mais elle est la condition \emph{sine qua non} sans laquelle le Cameroun restera éternellement le fournisseur de fèves des autres.
\end{exoresolu}

\begin{attention}[Erreurs fréquentes à signaler lors de la correction]
\begin{itemize}
\item Confondre \textbf{croissance} (augmentation quantitative du PIB) et \textbf{développement} (transformation qualitative globale : santé, éducation, autonomie, structure productive).
\item Affirmer que le sous-développement est une \og malédiction \fg{} naturelle ou culturelle : le cours montre qu'il a des causes historiques structurelles (colonisation, division internationale du travail, échange inégal) et qu'il peut être surmonté par l'action politique.
\item Citer des chiffres sans les interpréter philosophiquement : une donnée n'est un argument que si on la relie à une notion du cours (ex. : le rapport 1/5 illustre l'échange inégal).
\item Oublier l'objection dans le paragraphe argumenté : une argumentation philosophique rigoureuse doit affronter le point de vue contraire pour le dépasser.
\item Ne pas distinguer \emph{indépendance politique} (1960) et \emph{indépendance économique} (encore inachevée), distinction centrale chez Nkrumah.
\end{itemize}
\end{attention}

\courssub{Bilan de l'activité}

Cette activité a permis de rendre \emph{visibles} les notions abstraites de la leçon. Le diagramme en barres a transformé la définition de l'\cle{échange inégal} en un écart chiffré (rapport de 1 à 5) que chaque élève peut désormais mobiliser comme exemple précis dans une dissertation. Le croquis des flux a montré que le sous-développement n'est pas d'abord une pénurie de richesses naturelles — le Cameroun produit le cacao — mais une \emph{mauvaise position} dans la chaîne de valeur mondiale : le pays fournit la matière première et laisse aux autres le profit de la transformation, de la technologie et du marketing. Les documents de Nkrumah et de la ZLECAf ont enfin établi le pont entre l'histoire du \cle{panafricanisme} et l'actualité brûlante : l'unité africaine n'est pas un rêve sentimental, mais une stratégie économique rationnelle de rupture avec la division internationale du travail, dont les difficultés (frontières, infrastructures, financement, concurrence) sont elles-mêmes des problèmes philosophiques et politiques, non des fatalités. En rédigeant un paragraphe avec thèse, arguments, objection et réponse, les élèves ont enfin pratiqué la démarche exacte exigée à l'épreuve de dissertation du Baccalauréat technique, sur un sujet qui touche directement leur environnement national et leur avenir professionnel.

\newpage

\coursec{Méthodes et exercices résolus}

\coursec{Leçon N20 : Panafricanisme et développement de l’Afrique}

\courssub{Le développement : définition et dimensions}

\begin{definition}[Développement]
Le \cle{développement} est un processus global, continu et endogène par lequel une société améliore durablement le bien-être de l'ensemble de sa population. Il ne se réduit pas à la croissance économique (augmentation du PIB) ; il intègre trois dimensions indissociables :
\begin{itemize}
\item \textbf{Économique :} diversification de l'appareil productif, industrialisation, maîtrise de la transformation des matières premières, création d'emplois décents.
\item \textbf{Social :} accès universel à la santé, à l'éducation, à l'eau potable, au logement, protection sociale, réduction des inégalités.
\item \textbf{Humain et culturel :} épanouissement de la personne (libertés, participation citoyenne, dignité), valorisation des cultures endogènes, égalité homme-femme.
\end{itemize}
\end{definition}

\begin{propriete}[Indicateurs de développement (repères pour le Bac)]
\begin{itemize}
\item \textbf{PIB/habitant :} mesure la richesse moyenne produite (insuffisant seul).
\item \textbf{IDH (Indice de Développement Humain, PNUD) :} combine PIB/hab, espérance de vie, scolarisation (0 à 1).
\item \textbf{IDHI (IDH inégalités) :} corrige l'IDH selon les inégalités internes.
\item \textbf{ODD (Objectifs de Développement Durable, ONU 2015-2030) :} 17 objectifs de référence mondiale.
\end{itemize}
\textbf{Exigence Bac :} savoir distinguer \cle{croissance} (quantitatif) et \cle{développement} (qualitatif + structurel). Citer Amartya \textbf{Sen} : « Le développement, c'est l'expansion des libertés réelles dont jouissent les hommes. » (\emph{Development as Freedom}, 1999).
\end{propriete}

\begin{exemplebox}[Le paradoxe camerounais]
Le Cameroun affiche une croissance moyenne de 3 à 4 \% par an (PIB), mais son IDH reste faible (0,576 en 2022, rang 151/193). L'exploitation du pétrole, du bois, du cacao, du coton n'a pas entraîné l'industrialisation ni l'éradication de la pauvreté rurale (ex. : zones de production cacaoyère d'Obala ou Sangmélima sans routes ni eau potable). Cela illustre \textbf{croissance sans développement}.
\end{exemplebox}

\courssub{Le sous-développement : concept et manifestations en Afrique}

\begin{definition}[Sous-développement]
Le \cle{sous-développement} n'est pas un « stade » naturel ni l'absence de richesses, mais un \textbf{état structurel de retard} produit historiquement (traite, colonisation, néocolonialisme) et entretenu par des mécanismes de dépendance. Il se caractérise par :
\begin{itemize}
\item \textbf{Extraversion de l'économie :} spécialisation dans l'exportation de matières premières brutes (cacao, café, pétrole, minerais) et importation de produits manufacturés.
\item \textbf{Dépendance extérieure :} technologique, financière (dette, FMI/Banque Mondiale), alimentaire, idéologique.
\item \textbf{Pauvreté de masse :} monétaire (seuil 2,15 \$/jour) et multidimensionnelle (santé, éducation, énergie).
\item \textbf{Faiblesse des institutions :} instabilité politique, corruption, prédation des ressources par les élites.
\end{itemize}
\end{definition}

\begin{aretenir}[Théories explicatives (culture générale Bac)]
\begin{itemize}
\item \textbf{Théorie de la dépendance (A. G. Frank, S. Amin) :} le sous-développement du « Sud » est la contrepartie nécessaire du développement du « Nord » (échange inégal, transfert de valeur).
\item \textbf{Walter Rodney (\emph{Comment l'Europe a sous-développé l'Afrique}, 1972) :} la colonisation a détruit les structures autonomes africaines pour les intégrer de force au capitalisme mondial comme fournisseurs de matières premières.
\item \textbf{Samir Amin :} l'accumulation à l'échelle mondiale polarise les richesses vers les centres (Europe, USA, Chine) au détriment des périphéries (Afrique).
\end{itemize}
\end{aretenir}

\begin{exemplebox}[Manifestations concrètes au Cameroun]
\begin{itemize}
\item \textbf{Économie :} Exportation de 300 000 t de cacao brut/an (2\textsuperscript{e} exportateur africain) ; transformation locale < 30 \%. Importation de chocolat, huiles, farines.
\item \textbf{Social :} Taux de pauvreté \textasciitilde 37 \% (Banque Mondiale 2022) ; mortalité maternelle 438/100 000 naissances ; 1,5 million d'enfants hors école (UNICEF).
\item \textbf{Infrastructures :} Taux d'électrification rural < 25 \% ; réseau routier bitumé \textasciitilde 15 \%.
\end{itemize}
\end{exemplebox}

\courssub{Historique du panafricanisme}

\begin{methode}[Repères chronologiques obligatoires (à mémoriser pour le devoir)]
\begin{enumerate}
\item \textbf{Origines (fin XIX\textsuperscript{e} - 1945) :} Naissance dans la \textbf{diaspora} (Amériques, Caraïbes). Figures : \textbf{Edward Blyden} (retour à l'Afrique), \textbf{W.E.B. Du Bois} (premiers congrès panafricains : 1900 Londres, 1919 Paris, 1921 Londres/Bruxelles/Paris, 1923 Londres/Lisbonne, 1927 New York), \textbf{Marcus Garvey} (UNIA, « Africa for Africans », retour physique), \textbf{George Padmore} (liaison diaspora/continent).
\item \textbf{Tournant de Manchester (1945) :} 5\textsuperscript{e} Congrès panafricain. Passage du culturel au \textbf{politique}. Revendication explicite de l'indépendance. Présence de \textbf{Kwame Nkrumah} (Ghana), \textbf{Jomo Kenyatta} (Kenya), \textbf{Hastings Banda} (Malawi).
\item \textbf{Ère des Indépendances (1957-1963) :} Indépendance du Ghana (1957, Nkrumah). Conférence d'Accra (1958). Opposition \textbf{Groupe de Casablanca} (Nkrumah, Sékou Touré, Modibo Keïta : unité politique immédiate, fédéralisme) vs \textbf{Groupe de Monrovia} (Houphouët-Boigny, Senghor, Tubman : coopération fonctionnelle, souveraineté nationale progressive).
\item \textbf{Institutionnalisation (1963-2002) :} Création de l'\textbf{OUA} (Organisation de l'Unité Africaine) à Addis-Abeba, 25 mai 1963. Charte : respect des frontières héritées (\emph{uti possidetis juris}), non-ingérence, libération des territoires restants (Afrique du Sud, Sahara Occidental). Lutte contre l'apartheid.
\item \textbf{Renouveau (2002 - aujourd'hui) :} Transformation OUA $\rightarrow$ \textbf{Union Africaine (UA)} (Durban 2002). Objectifs : intégration économique, paix/sécurité (PSC), gouvernance (APRM), diaspora comme 6\textsuperscript{e} région. Lancement de la \textbf{ZLECAF} (Zone de Libre-Échange Continentale Africaine) signée 2018 (Kigali), entrée en vigueur 2021 : plus grand marché mondial par nombre de pays (1,3 milliard d'habitants, PIB cumulé 3 400 milliards \$).
\end{enumerate}
\end{methode}

\begin{attention}[Pièges historiques à éviter]
\begin{itemize}
\item Confondre \textbf{OUA} (1963, politique, libération) et \textbf{UA} (2002, intégration économique, interventionnisme Art. 4h).
\item Oublier la \textbf{diaspora} comme matrice originelle (Du Bois, Garvey, Padmore) : le panafricanisme n'est pas né en Afrique mais \emph{pour} l'Afrique.
\item Croire que l'unité politique a été réalisée en 1963 : l'OUA a consacré la \textbf{division} en États souverains (compromis de Monrovia).
\end{itemize}
\end{attention}

\courssub{Les objectifs du panafricanisme comme philosophie du développement}

\begin{propriete}[Cinq piliers objectifs (programme officiel)]
\begin{enumerate}
\item \textbf{Unité politique et souveraineté collective :} Dépasser les frontières artificielles (Berlin 1885) pour peser dans les instances mondiales (ONU, OMC, Climat). \textbf{Nkrumah} : « Seek ye first the political kingdom » (\emph{Africa Must Unite}, 1963). Un État fédéral africain (Armée commune, Monnaie unique, Citoyenneté unique).
\item \textbf{Intégration économique et autodéveloppement :} Marché commun, transformation locale des matières premières, industrialisation par substitution aux importations, monnaie unique (Eco). \textbf{Nyerere} (\emph{Ujamaa}, 1967) : « L'autonomie réelle passe par l'autosuffisance alimentaire et la maîtrise de l'outil de production. »
\item \textbf{Renaissance culturelle et identité :} Réhabilitation de l'histoire africaine (Cheikh Anta \textbf{Diop} : \emph{Nations nègres et culture}, 1954), valorisation des langues, éducation ancrée. \textbf{Senghor} : « La Négritude est le fait d'être nègre, la prise de conscience de ce fait, l'acceptation de ce fait. » Culture comme levier de développement endogène.
\item \textbf{Solidarité et paix continentale :} Mécanismes de prévention des conflits (UA-PSC), rejet des coups d'État, responsabilité de protéger (Art. 4h Charte UA). Développement impossible sans sécurité.
\item \textbf{Place dans la mondialisation :} Voix unique aux négociations climatiques (justice climatique), commerciales (OMC), financières (réforme FMI/Banque Mondiale, dette). « L'Afrique ne doit pas être un enjeu mais un acteur. »
\end{enumerate}
\end{propriete}

\courssub{Les difficultés du panafricanisme}

\begin{definition}[Obstacles structurels et conjoncturels]
\begin{itemize}
\item \textbf{Héritage colonial :} Frontières \emph{uti possidetis} (16 000 km de frontières tracées à Berlin), États « non-nations » (ethnies divisées, groupes rivaux regroupés), langues officielles étrangères (français, anglais, portugais) barrières à l'intégration populaire.
\item \textbf{Égoïsmes nationaux et souveraineté jalousée :} Chefs d'État réticents à céder des prérogatives à une autorité supranationale (ex. : retard ratification Protocole libre circulation personnes, 2018 ; résistance à l'armée commune).
\item \textbf{Dépendance néocoloniale persistante :} Accords de défense/monnaie (zone Franc CFA), liens privilégiés avec anciennes métropoles (Françafrique), ingérences dans les élections, bases militaires étrangères. « Diviser pour régner » continue.
\item \textbf{Disparités économiques et infrastructures manquantes :} PIB Nigeria \textasciitilde 477 Mrd \$ vs Burundi \textasciitilde 3 Mrd \$ ; déficit infrastructures (transport, énergie, numérique) estimé à 100 Mrd \$/an (BAD). ZLECAF freinée par barrières non-tarifaires, corruption aux frontières, manque de connectivité.
\item \textbf{Crises de gouvernance et instabilité :} Coups d'État récurrents (Mali, Guinée, Burkina, Niger, Gabon 2020-2023), terrorisme (Sahel, Lac Tchad, Mozambique), migrations forcées, détournement ressources.
\item \textbf{Faible appropriation populaire :} Panafricanisme souvent discours d'élites, absent des programmes scolaires, peu relayé par médias locaux. Jeunesse désabusée, tournée vers l'émigration.
\end{itemize}
\end{definition}

\begin{savaistu}[Le cas de la monnaie unique : l'Eco]
Prévue pour 2000, sans cesse reportée (2020, 2027...). Zone UEMOA (8 pays, CFA lié à l'Euro, garantie France) vs ZMAO (6 pays, dont Nigeria, Ghana). Divergences sur critères de convergence (inflation < 3\%, déficit < 3\%, dette < 70\% PIB), sur le rôle de la France, sur la banque centrale. Illustre la difficulté de l'intégration monétaire sans intégration politique/fiscale préalable.
\end{savaistu}

\courssub{Synthèse : le panafricanisme peut-il développer l'Afrique ?}

\begin{aretenir}[Réponse dialectique type Bac]
\begin{itemize}
\item \textbf{Thèse (Nécessité) :} L'unité est la \textbf{condition nécessaire}. Sans marché continental (ZLECAF), pas d'industrialisation viable (échelle). Sans voix unique, pas de rééquilibrage termes de l'échange. Sans armée commune, pas de sécurité durable. Nkrumah a raison : divisés, nous sommes faibles.
\item \textbf{Antithèse (Insuffisance) :} L'unité politique formelle ne garantit pas le développement. L'UE (modèle d'intégration) montre que le marché unique précède le politique. Des pays « non-panafricanistes » (Rwanda, Botswana, Maurice, Éthiopie jusqu'à 2020) ont mieux performé que des panafricanistes historiques (Ghana, Guinée, Mali). Le développement demande \textbf{gouvernance, éducation, institutions inclusives} (Acemoglu/Robinson, \emph{Why Nations Fail}) que l'unité seule ne crée pas.
\item \textbf{Synthèse :} Le panafricanisme est \textbf{nécessaire mais non suffisant}. Il doit muter : passer du \textbf{panafricanisme d'État} (sommets, discours) au \textbf{panafricanisme des peuples} (libre circulation effective, reconnaissance diplômes, médias panafricains, startup ZLECAF, jeunesse connectée). La ZLECAF est le test décisif : si elle fonctionne (commerce intra-africain 15\% $\rightarrow$ 25\%+), le panafricanisme devient développeur. Sinon, il reste un mythe mobilisateur sans effet structurel.
\end{itemize}
\end{aretenir}

\begin{exemplebox}[Sujet de dissertation type : « L'unité africaine est-elle la condition du développement du continent ? »]
\textbf{Plan dialectique :}
\begin{enumerate}
\item \textbf{L'unité comme condition nécessaire :} Briser l'étroitesse des marchés, mutualiser les moyens, peser géopolitiquement (Nkrumah, ZLECAF, UA).
\item \textbf{Les limites de l'unité sans transformation structurelle :} Échecs de l'OUA/UA, égoïsmes nationaux, gouvernance prédatrice, dépendance extérieure persistante (Rodney, Amin).
\item \textbf{Vers un panafricanisme par le bas, économique et citoyen :} ZLECAF opérationnelle, infrastructure (PIDA), éducation panafricaine, diaspora investisseuse, démocratie comme levier (Sen).
\end{enumerate}
\end{exemplebox}

% ============================================================
% BLOC MÉTHODES ET EXERCICES CORRIGÉS (REVU ET CORRIGÉ)
% ============================================================

\courssub{Méthode 1 : Appliquer les notions à une situation concrète}

\begin{methode}[Analyser une situation concrète avec les notions de développement et de sous-développement]
Pour appliquer les notions de l'unité à une situation de la vie courante (exemple : un village sans électricité, une jeunesse diplômée sans emploi, un pays qui exporte ses matières premières brutes), suivre les étapes suivantes :
\begin{enumerate}
\item \textbf{Lire attentivement la situation} et repérer les faits précis : chiffres, acteurs, difficultés, ressources.
\item \textbf{Dégager le problème philosophique} : que montre cette situation sur le développement de l'Afrique ? Formuler le problème sous forme de question (ex. : « Pourquoi la richesse des ressources ne suffit-elle pas au développement ? »).
\item \textbf{Mobiliser les notions du cours} : définir \cle{développement} (croissance économique + progrès social + épanouissement humain) et \cle{sous-développement} (retard structurel, dépendance, pauvreté), puis vérifier quelles dimensions sont présentes ou absentes dans la situation (économique, sociale, culturelle, politique).
\item \textbf{Appliquer explicitement} : montrer, fait par fait, comment la situation illustre la notion (ex. : exportation de cacao brut $\rightarrow$ dépendance économique, caractéristique du sous-développement).
\item \textbf{Conclure par un jugement nuancé} : la situation relève-t-elle du sous-développement total ou d'un développement inachevé ? Proposer une piste (panafricanisme, transformation locale, éducation).
\end{enumerate}
\textbf{Réflexes à retenir} : toujours définir la notion avant de l'appliquer ; citer au moins deux faits précis de la situation ; ne jamais rester dans le général.
\end{methode}

\begin{exoresolu}[Analyse d'une situation concrète]
\textbf{Énoncé.} Le Cameroun exporte chaque année des centaines de milliers de tonnes de cacao brut vers l'Europe, puis importe du chocolat fabriqué à l'étranger et vendu dix fois plus cher. Dans les zones de production comme à Obala ou à Sangmélima, de nombreux planteurs vivent sans eau potable ni routes praticables en saison des pluies. À l'aide des notions de développement et de sous-développement, analysez cette situation et dégagez-en le problème philosophique.
\tcblower
\textbf{Solution détaillée.}
\begin{enumerate}
\item \textbf{Repérage des faits.} Trois faits ressortent : (a) l'exportation de matière première brute (cacao) ; (b) l'importation du produit transformé (chocolat) à prix élevé ; (c) la précarité des conditions de vie des producteurs (eau, routes).
\item \textbf{Définitions des notions.} Le \cle{développement} est un processus global qui combine la croissance économique, le progrès social (santé, éducation, infrastructures) et l'épanouissement humain. Le \cle{sous-développement} désigne un état de retard structurel marqué par la pauvreté, la dépendance extérieure et l'incapacité à transformer soi-même ses richesses.
\item \textbf{Application.} L'exportation du cacao brut illustre la \textbf{dépendance économique} : la valeur ajoutée (la transformation en chocolat) est réalisée à l'étranger, donc les profits échappent au pays. Le fait que le chocolat importé coûte dix fois plus cher montre l'inégalité de l'échange : le Cameroun vend peu cher et achète cher. Enfin, l'absence d'eau potable et de routes chez les planteurs révèle le retard des infrastructures et du niveau de vie, dimensions sociales du sous-développement.
\item \textbf{Problème philosophique dégagé.} « Pourquoi un pays riche en ressources reste-t-il pauvre ? Le développement de l'Afrique passe-t-il nécessairement par la maîtrise de ses propres richesses et par l'unité africaine ? »
\item \textbf{Conclusion nuancée.} Cette situation n'est pas une fatalité : elle manifeste un sous-développement lié à l'organisation de l'économie (absence de transformation locale), non à l'absence de richesses. La piste panafricaniste — transformation sur place, marché commun africain — apparaît comme une réponse possible.
\end{enumerate}
\textbf{Phrase de conclusion à retenir} : « Le sous-développement n'est pas l'absence de richesses, mais l'incapacité structurelle à les transformer et à en faire profiter les populations. »
\end{exoresolu}

\courssub{Méthode 2 : Rédiger une introduction de dissertation}

\begin{methode}[Construire une introduction en quatre temps]
\begin{enumerate}
\item \textbf{Amorce (accroche)} : partir d'un fait général ou d'un constat lié au sujet (ex. : l'Afrique, continent riche, reste le plus pauvre).
\item \textbf{Présentation du sujet et définition des termes} : définir brièvement les mots-clés (\cle{panafricanisme}, \cle{développement}).
\item \textbf{Problématique} : formuler UNE question principale, ouverte, qui naît d'une tension (ex. : l'unité politique suffit-elle à développer un continent ?).
\item \textbf{Annonce du plan} : présenter les deux ou trois parties (thèse, antithèse, synthèse) sans les développer.
\end{enumerate}
\textbf{Piège à éviter} : ne jamais répondre à la problématique dans l'introduction ; ne jamais recopier le sujet tel quel.
\end{methode}

\begin{exoresolu}[Rédaction d'une introduction]
\textbf{Énoncé.} Rédigez l'introduction d'une dissertation sur le sujet suivant : « Le panafricanisme est-il une voie efficace pour le développement de l'Afrique ? »
\tcblower
\textbf{Solution détaillée (modèle rédigé).}

\emph{Amorce} : « Plus d'un demi-siècle après les indépendances, l'Afrique demeure, paradoxalement, le continent le plus riche en ressources naturelles et le plus pauvre en niveau de vie. »

\emph{Définitions} : « Face à ce paradoxe est né le \cle{panafricanisme}, mouvement de pensée et d'action qui prône l'unité des peuples africains et des descendants d'Africains afin de libérer et de développer le continent. Le \cle{développement}, quant à lui, désigne le processus par lequel une société améliore durablement ses conditions économiques, sociales et humaines. »

\emph{Problématique} : « Le panafricanisme se présente ainsi comme une philosophie du développement : c'est en s'unissant que les Africains pourraient peser face aux puissances extérieures et maîtriser leurs richesses. Mais l'unité politique et culturelle prônée par le panafricanisme suffit-elle réellement à engendrer le développement économique et social de l'Afrique ? »

\emph{Annonce du plan} : « Nous montrerons d'abord que le panafricanisme constitue une réponse pertinente au sous-développement ; nous verrons ensuite les difficultés qui limitent son efficacité ; nous dégagerons enfin les conditions d'un panafricanisme réellement développeur. »

\textbf{Justification de la démarche} : chaque étape est présente (fait, définitions, question ouverte née du paradoxe richesse/pauvreté, plan en trois parties), et aucune réponse n'est donnée d'avance.
\end{exoresolu}

\courssub{Méthode 3 : Construire un paragraphe argumenté}

\begin{methode}[La règle du paragraphe unique : idée -- définition -- auteur -- exemple -- analyse]
Un paragraphe de dissertation philosophique suit cinq étapes :
\begin{enumerate}
\item \textbf{Idée} : énoncer clairement la thèse du paragraphe en une phrase.
\item \textbf{Définition} : définir la notion centrale mobilisée.
\item \textbf{Appui théorique} : citer brièvement un auteur du cours (Nkrumah, Senghor, Cheikh Anta Diop, Nyerere...) avec sa thèse, sans longue citation.
\item \textbf{Exemple concret} : illustrer par un fait précis, de préférence africain ou camerounais (création de l'OUA en 1963, ZLECAF, crise d'une frontière héritée de la colonisation).
\item \textbf{Analyse} : expliquer le lien entre l'exemple et l'idée, puis conclure le paragraphe.
\end{enumerate}
\textbf{Formule à retenir} : « Une idée, un auteur, un exemple, une explication. »
\end{methode}

\begin{exoresolu}[Rédaction d'un paragraphe argumenté]
\textbf{Énoncé.} Rédigez un paragraphe argumenté montrant que l'unité africaine est une condition du développement du continent.
\tcblower
\textbf{Solution détaillée (modèle rédigé).}

\emph{Idée} : « L'unité politique et économique des États africains apparaît d'abord comme la condition première de leur développement. »

\emph{Définition} : « En effet, le \cle{panafricanisme} peut se définir comme la volonté de dépasser les frontières artificielles héritées de la colonisation pour constituer un ensemble solidaire, capable de peser sur la scène mondiale. »

\emph{Appui théorique} : « Kwame Nkrumah, père de l'indépendance du Ghana, soutenait qu'aucun État africain isolé ne pouvait se développer seul : selon lui, seule l'union politique du continent pouvait briser la dépendance néocoloniale, car divisés, les États restent faibles face aux anciennes puissances coloniales. »

\emph{Exemple} : « La création de l'Organisation de l'unité africaine (OUA) en 1963 à Addis-Abeba, puis de la Zone de libre-échange continentale africaine (ZLECAF) en 2018, traduit concrètement cette volonté : en échangeant davantage entre eux, les pays africains réduisent leur dépendance envers l'extérieur. »

\emph{Analyse et conclusion} : « Ainsi, l'exemple de la ZLECAF montre que l'unité n'est pas un simple rêve culturel : elle crée un marché de plus d'un milliard de consommateurs, condition pour industrialiser le continent. L'unité africaine apparaît donc bien comme un levier du développement. »

\textbf{Justification} : le paragraphe respecte les cinq étapes et chaque élément est relié logiquement au suivant.
\end{exoresolu}

\courssub{Méthode 4 : Problématiser un sujet}

\begin{methode}[Transformer un sujet en problème philosophique]
\begin{enumerate}
\item \textbf{Repérer les termes du sujet} et les définir (panafricanisme, développement, efficacité, unité...).
\item \textbf{Dégager la tension} : chercher l'opposition cachée (unité culturelle vs division politique ; richesse des ressources vs pauvreté ; idéal panafricaniste vs échecs pratiques).
\item \textbf{Formuler la question} : elle doit être ouverte (on ne peut pas répondre seulement par oui ou par non sans discussion) et contenir les termes du sujet.
\item \textbf{Vérifier} : la question permet-elle au moins deux réponses opposées ? Si oui, elle est problématisante.
\end{enumerate}
\textbf{Réflexe} : une bonne problématique contient souvent « suffit-il », « peut-il », « faut-il », « dans quelle mesure ».
\end{methode}

\begin{exoresolu}[Problématisation d'un sujet]
\textbf{Énoncé.} Sujet : « Faut-il être panafricaniste pour développer l'Afrique ? » Dégagez la tension du sujet et formulez une problématique.
\tcblower
\textbf{Solution détaillée.}
\begin{enumerate}
\item \textbf{Termes} : « panafricaniste » renvoie à la doctrine de l'unité africaine ; « développer » renvoie au processus de transformation économique, sociale et humaine ; « faut-il » invite à discuter une nécessité.
\item \textbf{Tension dégagée} : d'un côté, le panafricanisme affirme que l'unité est la condition du développement (thèse de Nkrumah) ; de l'autre, certains pays se développent en privilégiant des stratégies nationales ou des alliances extérieures, ce qui suggère que le panafricanisme n'est peut-être ni suffisant ni indispensable. La tension oppose donc \emph{l'unité comme condition} et \emph{l'unité comme simple option}.
\item \textbf{Problématique formulée} : « L'unité prônée par le panafricanisme est-elle la condition nécessaire du développement de l'Afrique, ou le développement peut-il advenir par d'autres voies ? »
\item \textbf{Vérification} : deux réponses opposées sont possibles (oui : Nkrumah, l'expérience de la ZLECAF ; non : les réussites nationales ponctuelles, les échecs de l'unité politique). La question est donc bien problématisante.
\end{enumerate}
\textbf{Conclusion} : cette problématique autorise un plan dialectique en trois parties (nécessité de l'unité / limites de l'unité / conditions d'un panafricanisme efficace).
\end{exoresolu}

\courssub{Méthode 5 : Rédiger et justifier une conclusion}

\begin{methode}[Construire une conclusion justifiée]
\begin{enumerate}
\item \textbf{Rappeler la problématique} en la reformulant brièvement.
\item \textbf{Bilan du parcours} : résumer en une ou deux phrases les positions confrontées (thèse et antithèse).
\item \textbf{Réponse argumentée} : prendre clairement position et justifier ce choix par l'argument le plus fort du devoir.
\item \textbf{Ouverture} : poser une nouvelle question qui prolonge la réflexion (ex. : quelles formes nouvelles le panafricanisme peut-il prendre à l'ère du numérique ?).
\end{enumerate}
\textbf{Interdit} : introduire une notion nouvelle non étudiée dans le développement.
\end{methode}

\begin{exoresolu}[Rédaction d'une conclusion]
\textbf{Énoncé.} Rédigez la conclusion de la dissertation : « Le panafricanisme est-il une voie efficace pour le développement de l'Afrique ? », sachant que le développement a montré d'abord la pertinence du panafricanisme, ensuite ses difficultés (divisions des États, dépendance, égoïsmes nationaux).
\tcblower
\textbf{Solution détaillée (modèle rédigé).}

\emph{Rappel de la problématique} : « Nous nous demandions si le panafricanisme constitue une voie efficace pour le développement de l'Afrique. »

\emph{Bilan} : « Notre réflexion a montré que le panafricanisme, en prônant l'unité politique et économique, offre une réponse cohérente au sous-développement et à la dépendance hérités de la colonisation ; mais nous avons aussi constaté que les divisions entre États, les frontières artificielles et les égoïsmes nationaux ont longtemps freiné sa réalisation. »

\emph{Réponse justifiée} : « Nous soutenons donc que le panafricanisme demeure une voie nécessaire, mais non suffisante : il ne développera l'Afrique que s'il passe du discours culturel à l'action économique concrète, comme le montre l'espoir suscité par la ZLECAF, qui unit les marchés africains au-delà des frontières coloniales. »

\emph{Ouverture} : « Reste alors à se demander si la jeunesse africaine, à l'ère du numérique, ne sera pas la véritable force de ce panafricanisme renouvelé. »

\textbf{Justification de la démarche} : la conclusion reprend la question, fait le bilan des deux parties, tranche avec un argument précis (le passage à l'action économique) et ouvre sans introduire de notion nouvelle.
\end{exoresolu}

\courssub{Méthode 6 : Mobiliser l'histoire du panafricanisme dans un devoir}

\begin{methode}[Utiliser les repères historiques sans tomber dans le récit]
\begin{enumerate}
\item \textbf{Mémoriser les repères essentiels} : origines dans la diaspora (fin du XIX\textsuperscript{e} siècle), congrès de Manchester en 1945, indépendances (années 1960), création de l'OUA en 1963, transformation en Union africaine en 2002.
\item \textbf{Sélectionner} : ne retenir que le repère utile à l'argument (ex. : Manchester 1945 pour montrer le passage de l'idée culturelle au combat politique).
\item \textbf{Faire parler le repère} : après chaque fait, ajouter une phrase d'analyse qui le relie à la thèse du paragraphe.
\item \textbf{Ne jamais raconter} : l'histoire sert l'argumentation philosophique, elle ne la remplace pas.
\end{enumerate}
\textbf{Formule à retenir} : « Un fait historique + une analyse = un argument. Un fait seul = du hors-sujet. »
\end{methode}

\begin{exoresolu}[Mobilisation d'un repère historique]
\textbf{Énoncé.} Dans un paragraphe visant à montrer que le panafricanisme est né d'une prise de conscience, utilisez le congrès de Manchester (1945) comme argument.
\tcblower
\textbf{Solution détaillée (modèle rédigé).}

« Le panafricanisme n'est pas d'abord une théorie abstraite : il naît d'une prise de conscience historique, celle de la communauté de destin entre Africains du continent et descendants d'esclaves déportés. Le cinquième congrès panafricain de Manchester, en 1945, marque à cet égard un tournant décisif : réunissant des figures comme Kwame Nkrumah et Jomo Kenyatta, il transforme le panafricanisme, jusque-là mouvement culturel de la diaspora, en revendication politique exigeant l'indépendance des peuples africains. Ce fait montre que la conscience de l'unité africaine précède et prépare la libération politique : les indépendances des années 1960 sont ainsi le fruit d'une philosophie devenue action. Le panafricanisme apparaît donc, dès son histoire, comme une philosophie engagée au service de l'émancipation et du développement du continent. »

\textbf{Justification} : le fait (Manchester 1945) est daté, incarné par des acteurs, et immédiatement suivi d'une analyse (passage de la culture à la politique) reliée à la thèse du paragraphe. On évite ainsi le récit gratuit.
\end{exoresolu}

\begin{attention}[Les 5 erreurs qui coûtent des points au Bac]
\begin{enumerate}
\item \textbf{Confondre croissance et développement.} La croissance est une simple augmentation de la richesse produite ; le développement inclut en plus le progrès social et humain (santé, éducation, libertés). Écrire « développement = richesse » est une erreur de définition sanctionnée.
\item \textbf{Raconter l'histoire au lieu d'argumenter.} Aligner les dates (1945, 1963, 2002...) sans les relier à une thèse transforme la dissertation en exposé d'histoire : chaque fait doit être suivi d'une analyse.
\item \textbf{Oublier de définir les termes du sujet.} Un devoir qui utilise « panafricanisme » ou « sous-développement » sans les définir dans l'introduction perd des points de méthode : la définition est la première étape de toute dissertation.
\item \textbf{Donner la réponse dans l'introduction.} La problématique doit rester une question ouverte ; annoncer d'emblée « le panafricanisme est la solution » supprime toute discussion et condamne le plan dialectique.
\item \textbf{Citer un auteur sans sa thèse.} Écrire « comme le dit Nkrumah » sans préciser sa position (l'unité politique comme condition du développement) est un appui théorique vide : toujours accompagner le nom de l'idée et la relier à l'exemple.
\end{enumerate}
\end{attention}

\newpage

\coursec{Exercices}

\coursec{Panafricanisme et développement de l'Afrique}

\courssub{Je vérifie mes connaissances}

\begin{exercice}[QCM : notions fondamentales]
Pour chaque question, une seule réponse est correcte. Recopie la lettre correspondante.
\begin{enumerate}
\item Le développement se distingue de la simple croissance économique parce qu'il suppose :
\begin{enumerate}
\item une augmentation du PIB uniquement ;
\item une transformation durable des structures économiques, sociales et culturelles au profit du bien-être de tous ;
\item l'augmentation des exportations de matières premières.
\end{enumerate}
\item Le sous-développement désigne :
\begin{enumerate}
\item l'absence totale de ressources naturelles ;
\item un état de blocage des structures économiques et sociales qui empêche l'épanouissement des populations ;
\item un retard culturel irréversible des peuples.
\end{enumerate}
\item Le panafricanisme est né historiquement :
\begin{enumerate}
\item dans la diaspora africaine (Amériques, Caraïbes) à la fin du XIX\textsuperscript{e} siècle ;
\item au sein de l'Organisation des Nations unies en 1945 ;
\item lors de la création de la CEDEAO en 1975.
\end{enumerate}
\item L'OUA, créée en 1963 à Addis-Abeba, a été remplacée en 2002 par :
\begin{enumerate}
\item la ZLECAF ;
\item l'Union africaine (UA) ;
\item la CEMAC.
\end{enumerate}
\end{enumerate}
\end{exercice}

\begin{exercice}[Vrai ou faux justifié]
Indique si chaque affirmation est vraie ou fausse, puis justifie ta réponse en deux ou trois lignes.
\begin{enumerate}
\item « Un pays peut connaître une forte croissance économique tout en restant sous-développé. »
\item « Le panafricanisme est uniquement un mouvement culturel sans visée politique. »
\item « Le faible niveau du commerce intra-africain est l'une des difficultés du panafricanisme. »
\item « Le sous-développement de l'Afrique s'explique exclusivement par la colonisation. »
\end{enumerate}
\end{exercice}

\begin{exercice}[Définitions]
Définis en trois à cinq lignes chacun des concepts suivants, en donnant pour chacun un exemple concret tiré du contexte camerounais ou africain :
\begin{enumerate}
\item \cle{développement} ;
\item \cle{sous-développement} ;
\item \cle{panafricanisme}.
\end{enumerate}
\end{exercice}

\begin{exercice}[Repères chronologiques]
Replace les événements suivants sur une frise chronologique et indique en une phrase l'importance de chacun pour le panafricanisme : congrès de Manchester (1945), indépendance du Ghana (1957), création de l'OUA (1963), création de l'Union africaine (2002), entrée en vigueur de la ZLECAF (2021).
\end{exercice}

\courssub{Je m'entraîne}

\begin{exercice}[Croissance ou développement ?]
Voici la situation de deux pays africains fictifs. Le pays A voit son PIB augmenter de 6 \% par an grâce à l'exportation de pétrole brut, mais 40 \% de sa population vit sous le seuil de pauvreté, l'espérance de vie stagne et les écoles manquent. Le pays B a une croissance de 4 \% par an, investit dans l'éducation, la santé et la transformation locale de ses produits agricoles.
\begin{enumerate}
\item Lequel des deux pays est le plus engagé sur la voie du développement ? Justifie ta réponse à partir de la définition philosophique du développement.
\item Montre en cinq lignes pourquoi la croissance économique ne suffit pas au développement.
\end{enumerate}
\end{exercice}

\begin{exercice}[Manifestations du sous-développement]
À partir de tes connaissances et de l'observation de la vie courante au Cameroun, complète le tableau suivant en donnant deux manifestations concrètes du sous-développement pour chaque dimension.
\begin{center}
\begin{tabularx}{\linewidth}{|l|X|}
\hline
\rowcolor{popblueL} \textbf{Dimension} & \textbf{Manifestations concrètes} \\
\hline
Économique & \\
\hline
Sociale (santé, éducation) & \\
\hline
Politique et institutionnelle & \\
\hline
Culturelle et mentale & \\
\hline
\end{tabularx}
\end{center}
\end{exercice}

\begin{exercice}[Objectifs du panafricanisme]
Le discours de lancement de la ZLECAF (Zone de libre-échange continentale africaine) affirme : « En commerçant davantage entre nous, les Africains créeront des emplois, transformeront leurs matières premières et parleront d'une seule voix dans le monde. »
\begin{enumerate}
\item Relève dans ce texte deux objectifs du panafricanisme comme philosophie du développement.
\item Explique en quelques lignes comment l'unité africaine peut favoriser l'indépendance économique du continent.
\end{enumerate}
\end{exercice}

\begin{exercice}[Débat philosophique guidé]
Sujet : « Le sous-développement de l'Afrique est-il d'abord la faute des autres ? »
\begin{enumerate}
\item Rédige en cinq lignes la thèse (le sous-développement s'explique par les causes extérieures : traite, colonisation, domination économique actuelle).
\item Rédige en cinq lignes l'antithèse (le sous-développement s'explique d'abord par des causes intérieures : mauvaise gouvernance, corruption, divisions).
\item Rédige en huit lignes une synthèse personnelle justifiée, en t'appuyant sur un exemple camerounais ou africain précis.
\end{enumerate}
\end{exercice}

\courssub{Je me prépare au Bac}

\begin{exercice}[Dissertation — type Baccalauréat technique]
Sujet : « Le panafricanisme est-il encore une voie pertinente pour le développement de l'Afrique ? »

Tu rédigeras une dissertation complète comportant :
\begin{enumerate}
\item une introduction avec accroche, définition des termes du sujet, problématique et annonce du plan ;
\item un développement en deux ou trois parties argumentées (tu pourras mobiliser l'historique du panafricanisme, ses objectifs, ses difficultés actuelles, ainsi que des exemples comme la ZLECAF ou l'Union africaine) ;
\item une conclusion qui répond clairement à la problématique et ouvre sur une perspective.
\end{enumerate}
\end{exercice}

\begin{exercice}[Explication de texte — Kwame Nkrumah]
Texte : « L'indépendance politique n'est qu'un préalable. Sans unité, nos États resteront trop faibles pour résister aux pressions extérieures et pour construire des économies véritablement indépendantes. C'est dans l'union que l'Afrique trouvera la force de transformer ses richesses au profit de ses peuples. Cherchez d'abord le royaume politique, et tout le reste vous sera donné en plus. » (D'après Kwame Nkrumah, \emph{L'Afrique doit s'unir}, 1963.)
\begin{enumerate}
\item Dégage la thèse de l'auteur et montre comment elle s'articule en deux étapes (indépendance politique, puis unité).
\item Explique le lien que Nkrumah établit entre unité politique et indépendance économique.
\item Ce texte te paraît-il encore d'actualité à l'heure de la ZLECAF ? Justifie ta réponse en mobilisant les difficultés du panafricanisme étudiées dans la leçon.
\end{enumerate}
\end{exercice}

\begin{exercice}[Question à réponse construite — le commerce intra-africain]
Document : part du commerce intra-africain dans le commerce total de chaque continent (estimations, en \%).

\begin{center}
\begin{tabularx}{\linewidth}{|l|X|X|X|}
\hline
\rowcolor{popblueL} \textbf{Continent} & \textbf{Échanges internes (\%)} & \textbf{Échanges avec l'extérieur (\%)} \\
\hline
Europe & environ 67 & environ 33 \\
\hline
Asie & environ 59 & environ 41 \\
\hline
Amérique du Nord & environ 31 & environ 69 \\
\hline
Afrique & environ 15 & environ 85 \\
\hline
\end{tabularx}
\end{center}

\begin{enumerate}
\item Calcule l'écart (en points de pourcentage) entre la part du commerce intra-africain et celle du commerce intra-européen. Que révèle cet écart ?
\item À partir du document et de la leçon, dégage deux difficultés que rencontre le panafricanisme dans la voie du développement.
\item Propose deux solutions concrètes, inspirées des objectifs du panafricanisme, pour renforcer l'intégration économique africaine, et justifie chacune en trois lignes.
\end{enumerate}
\end{exercice}

\newpage

\coursec{Corrigés des exercices}

\coursec{Leçon 20 : Panafricanisme et développement de l'Afrique — Exercices corrigés}

\courssub{Exercice 1 — QCM : notions fondamentales}
\begin{corrige}[Exercice 1 — QCM : notions fondamentales]
\begin{enumerate}
\item \textbf{b.} Le développement est un processus global, qualitatif et durable qui transforme les structures (économiques, sociales, culturelles, politiques) pour améliorer le bien-être de l'ensemble de la population. La croissance (augmentation du PIB) n'en est qu'une condition nécessaire, non suffisante.
\item \textbf{b.} Le sous-développement n'est pas l'absence de ressources (l'Afrique en regorge), mais un état structurel de dépendance, de blocage et de mauvaise allocation des ressources qui empêche l'épanouissement humain (selon A.G. Frank, Samir Amin).
\item \textbf{a.} Le panafricanisme naît dans la diaspora (États-Unis, Caraïbes) à la fin du XIX\textsuperscript{e} siècle (conférences panafricaines de 1900, 1919, 1921, 1923, 1927, 1945) comme réaction à l'esclavage, à la colonisation et au racisme, avant de revenir sur le continent.
\item \textbf{b.} L'Organisation de l'Unité Africaine (OUA), créée le 25 mai 1963 à Addis-Abeba, a été dissoute le 9 juillet 2002 pour laisser place à l'Union africaine (UA), dotée d'institutions plus fortes (Commission, Parlement, Cour de justice) et d'une vision d'intégration plus profonde.
\end{enumerate}
\end{corrige}

\courssub{Exercice 2 — Vrai ou faux justifié}
\begin{corrige}[Exercice 2 — Vrai ou faux justifié]
\begin{enumerate}
\item \textbf{Vrai.} La croissance mesure l'augmentation quantitative de la production (PIB). Le développement exige une amélioration qualitative du niveau de vie (santé, éducation, équité, liberté). Un pays pétrolier peut voir son PIB grimper (croissance) tandis que la pauvreté, la mortalité infantile et les inégalités persistent (sous-développement) : c'est la « croissance sans développement ».
\item \textbf{Faux.} Le panafricanisme est d'abord un projet \textbf{politique} (unité continentale, souveraineté collective, siège unique à l'ONU, défense commune) et \textbf{économique} (marché commun, transformation locale, ZLECAF). La dimension culturelle (négritude, renaissance africaine) en est le socle, mais non l'unique finalité.
\item \textbf{Vrai.} Le commerce intra-africain stagne autour de 15-17\% des échanges totaux du continent (contre ~67\% en Europe, ~59\% en Asie). Ce repli sur l'extérieur (extraversion) maintient les économies africaines en position de fournisseurs de matières premières et d'importateurs de produits finis, freinant l'industrialisation et l'autonomie.
\item \textbf{Faux.} La colonisation (traite, pillage, frontières artificielles, économie d'enclave) est une cause historique majeure, mais non exclusive. Les causes \textbf{internes} pèsent tout autant aujourd'hui : mauvaise gouvernance, corruption, instabilité politique, sous-investissement dans le capital humain, démographie galopante, absence de projet industriel cohérent.
\end{enumerate}
\end{corrige}

\courssub{Exercice 3 — Définitions}
\begin{corrige}[Exercice 3 — Définitions]
\begin{enumerate}
\item \textbf{Développement :} Processus global, durable et endogène de transformation structurelle de la société (économie, social, politique, culture, environnement) visant l'amélioration continue du bien-être de \textbf{tous} les individus (libertés réelles, capacités \`a la Sen). \textit{Exemple camerounais :} La Stratégie Nationale de Développement 2020-2030 (SND30) qui vise la transformation structurelle de l'économie par l'industrialisation, l'emploi des jeunes et l'accès aux services sociaux de base.
\item \textbf{Sous-développement :} Situation structurelle de dépendance et de blocage où un pays, malgré ses potentialités, ne parvient pas à satisfaire les besoins fondamentaux de sa population (santé, éducation, emploi, souveraineté) du fait d'une économie extravertie, d'institutions faibles et d'inégalités massives. \textit{Exemple camerounais :} Le Cameroun exporte du cacao, du bois et du pétrole bruts (faible valeur ajoutée) et importe du chocolat, du mobilier et des produits pétroliers raffinés, créant un déficit commercial structurel et peu d'emplois locaux.
\item \textbf{Panafricanisme :} Idéologie et mouvement politique prônant l'unité, la solidarité et la coopération entre les peuples africains et la diaspora, pour conquérir la souveraineté politique, l'indépendance économique et la renaissance culturelle du continent. \textit{Exemple :} La création de la Zone de Libre-Échange Continentale Africaine (ZLECAF) en 2021, outil juridique et économique pour réaliser le marché unique panafricain rêvé par Nkrumah.
\end{enumerate}
\end{corrige}

\courssub{Exercice 4 — Repères chronologiques}
\begin{corrige}[Exercice 4 — Repères chronologiques]
\begin{center}
\begin{tabularx}{\linewidth}{|c|X|X|}
\hline
\rowcolor{popblueL} \textbf{Date / Événement} & \textbf{Position sur la frise} & \textbf{Importance pour le panafricanisme} \\
\hline
1945 : Congrès de Manchester & \textbf{1945} — Tournant décisif & Passage du panafricanisme « intellectuel » au panafricanisme « de masse » et d'action politique. Revendication explicite de l'indépendance immédiate. Nkrumah, Kenyatta, Padmore y posent les bases de la décolonisation. \\
\hline
1957 : Indépendance du Ghana & \textbf{1957} — Premier État libre & Le Ghana de Nkrumah devient la « Mecque » du panafricanisme : base arrière des mouvements de libération, hôte de la Conférence des États africains indépendants (1958) et de la Conférence des Peuples africains (1958). Preuve que l'indépendance est possible. \\
\hline
1963 : Création de l'OUA & \textbf{1963} — Institutionnalisation & Charte d'Addis-Abeba : consécration juridique de la souveraineté, de l'intangibilité des frontières (héritage colonial) et de la solidarité. Outil de fin de la décolonisation (apartheid), mais limité par le principe de non-ingérence. \\
\hline
2002 : Création de l'UA & \textbf{2002} — Refondation & Remplace l'OUA. Vision élargie : intégration politique accélérée, gouvernance démocratique, paix/sécurité (PSC), droits de l'homme (CADHP), développement endogène (NEPAD). Passage de « l'unité des États » à « l'unité des peuples ». \\
\hline
2021 : Entrée en vigueur ZLECAF & \textbf{2021} — Outil économique & Plus grand marché unique au monde (1,3 milliard d'habitants, PIB ~3 400 Mds \$). Vise à booster le commerce intra-africain (+52\% visé), l'industrialisation, la transformation locale. Concrétisation économique du rêve panafricain. \\
\hline
\end{tabularx}
\end{center}
\end{corrige}

\courssub{Exercice 5 — Croissance ou développement ?}
\begin{corrige}[Exercice 5 — Croissance ou développement ?]
\begin{enumerate}
\item \textbf{Le pays B est le plus engagé sur la voie du développement.}
\textbf{Justification :} Philosophiquement (Aristote, Amartya Sen, PNUD), le développement est l'expansion des \textbf{libertés réelles} et des \textbf{capacités} humaines (vivre longtemps, être instruit, participer à la vie communautaire, avoir un niveau de vie décent). Le pays B investit dans l'éducation, la santé et la transformation locale : il construit du capital humain et productif, bases du développement endogène. Le pays A subit la « malédiction de la rente » : croissance tirée par l'exportation d'une rente pétrolière non renouvelable, sans redistribution (40\% pauvres), sans services publics, sans diversification. C'est une \textbf{croissance sans développement}.
\item \textbf{Pourquoi la croissance ne suffit pas au développement (5 lignes) :}
La croissance (ΔPIB > 0) est un indicateur \textbf{quantitatif} agrégé, muet sur la \textbf{qualité} de la production, la \textbf{distribution} des revenus, la \textbf{durabilité} écologique et les \textbf{libertés} réelles. Un PIB peut croître par l'exploitation de mines polluantes, le travail des enfants, l'endettement ou la rente, sans créer d'emplois décents, d'écoles, d'hôpitaux ni d'autonomie technologique. Le développement exige une \textbf{transformation structurelle} (diversification, industrialisation), de l'\textbf{équité} (réduction pauvreté/inégalités) et de la \textbf{soutenabilité} (capital naturel préservé).
\end{enumerate}
\end{corrige}

\courssub{Exercice 6 — Manifestations du sous-développement}
\begin{corrige}[Exercice 6 — Manifestations du sous-développement]
\begin{center}
\begin{tabularx}{\linewidth}{|l|X|}
\hline
\rowcolor{popblueL} \textbf{Dimension} & \textbf{Manifestations concrètes (Cameroun / Afrique)} \\
\hline
Économique & 1. \textbf{Extraversion structurelle :} Exportation de matières premières brutes (cacao, bois, coton, pétrole, minerais) sans transformation locale ; importation massive de produits manufacturés (médicaments, machines, alimentation transformée). \\
& 2. \textbf{Économie informelle prédominante :} ~90\% des emplois en Afrique subsaharienne (BIT) ; précarité, absence de protection sociale, faible productivité, évasion fiscale massive privant l'État de ressources. \\
\hline
Sociale (santé, éducation) & 1. \textbf{Indicateurs sanitaires critiques :} Taux de mortalité maternelle (438/100 000 naissances vivantes au Cameroun en 2020) et infantile élevés ; désert médical en zone rurale ; dépendance aux importations de vaccins/médicaments (Covid-19 a révélé la vulnérabilité). \\
& 2. \textbf{Éducation de masse mais qualité faible :} Scolarisation primaire quasi universelle mais taux d'achèvement bas, surcharge des classes (80-100 élèves), inadéquation formation/emploi (chômage diplômés), fuite des cerveaux vers le Nord. \\
\hline
Politique et institutionnelle & 1. \textbf{Fragilité de l'État de droit :} Corruption systémique (Cameroun souvent mal classé à l'indice CPI de Transparency International), détournements de fonds publics, justice aux ordres, impunité des élites. \\
& 2. \textbf{Instabilité et conflits :} Crise anglophone (Nord-Ouest/Sud-Ouest) depuis 2016, insécurité au Sahel (terrorisme), coups d'État récurrents (Mali, Burkina, Guinée, Niger, Gabon) : l'État ne monopolise plus la violence légitime sur tout le territoire. \\
\hline
Culturelle et mentale & 1. \textbf{Aliénation culturelle / dévalorisation du local :} Préférence pour les diplômes, modèles de vie, produits, langues (français/anglais) et esthétiques occidentaux ; mépris des savoirs endogènes, médecines traditionnelles, langues nationales. \\
& 2. \textbf{Mentalité de rente et d'assistanat :} Attente passive de l'État ou de l'aide internationale (ONG, diaspora) pour résoudre les problèmes collectifs ; faible esprit d'initiative entrepreneuriale formelle ; clientélisme politique (vote contre services). \\
\hline
\end{tabularx}
\end{center}
\end{corrige}

\courssub{Exercice 7 — Objectifs du panafricanisme}
\begin{corrige}[Exercice 7 — Objectifs du panafricanisme]
\begin{enumerate}
\item \textbf{Deux objectifs explicites dans le texte :}
\begin{enumerate}
\item \textbf{L'intégration économique par le commerce intra-africain :} « En commerçant davantage entre nous, les Africains créeront des emplois, transformeront leurs matières premières. » C'est l'objectif de \textbf{structuration d'un marché commun} (ZLECAF) pour briser l'extraversion.
\item \textbf{L'affirmation de la souveraineté collective sur la scène mondiale :} « Parleront d'une seule voix dans le monde. » C'est l'objectif \textbf{politique et géopolitique} : poids négociateur face aux puissances (UE, Chine, USA), défense des intérêts africains (climat, dette, termes de l'échange, réforme ONU).
\end{enumerate}
\item \textbf{Comment l'unité favorise l'indépendance économique :}
Des États fragmentés (54 États, petits marchés) sont des « proies » faciles pour les multinationales et les puissances : ils négocient isolément, acceptent des termes défavorables, subissent la concurrence fiscale (course au moins-disant fiscal). L'unité crée un \textbf{marché unique de 1,3 milliard de consommateurs} : effets d'échelle pour l'industrie locale, pouvoir de négociation des contrats miniers/pétroliers, mutualisation des infrastructures (énergie, transport, numérique), monnaie commune (projet Eco) pour réduire les coûts de transaction et la volatilité. L'indépendance économique naît de la \textbf{capacité collective à transformer localement ses richesses} et à fixer ses prix.
\end{enumerate}
\end{corrige}

\courssub{Exercice 8 — Débat philosophique guidé}
\begin{corrige}[Exercice 8 — Débat philosophique guidé]
\begin{enumerate}
\item \textbf{Thèse (causes extérieures) :} Le sous-développement africain est d'abord l'héritage direct de l'histoire violente imposée par l'extérieur. La \textbf{traite négrière} (12-15 millions d'êtres humains, dépeuplement, démographie brisée, guerres inter-africaines attisées) a vidé le continent de ses forces vives. La \textbf{colonisation} a imposé des frontières artificielles (sources de conflits), une économie d'enclave (infrastructures port-chemin de fer-mine, pas de liaison inter-États), une administration prédatrice et un système éducatif aliénant. Aujourd'hui, le \textbf{néocolonialisme} (terme de Kwame Nkrumah) se perpétue par des accords de partenariat inégaux (APE), le franc CFA (pour la zone franc), la dette odieuse, le pillage des ressources par des multinationales complices d'élites locales, et les conditionnalités des institutions de Bretton Woods (FMI/Banque mondiale) qui imposent l'austérité et la libéralisation prématurée.
\item \textbf{Antithèse (causes intérieures) :} Accuser uniquement l'extérieur est une fuite en avant qui déresponsabilise les élites africaines. Depuis les indépendances (années 1960), la \textbf{mauvaise gouvernance} est le premier frein : corruption endémique (détournement de milliards \$/an), népotisme, clientélisme, absence de reddition de comptes, constitutions taillées sur mesure pour la présidence à vie. Les \textbf{divisions ethno-régionales} instrumentalisées par le pouvoir empêchent la cohésion nationale. L'\textbf{absence de vision stratégique} : pas de politique industrielle cohérente, négligence de l'agriculture vivrière, sous-investissement chronique dans l'éducation/santé/recherche (souvent < 3-5\% du PIB), fuite des capitaux (évasion fiscale/illicite > aide reçue). La démographie non maîtrisée (fécondité > 4,5 enfants/femme) annule les gains de croissance.
\item \textbf{Synthèse :} Le sous-développement est le produit \textbf{dialectique} d'une \textbf{structure historique défavorable} (héritage colonial/néocolonial) et d'une \textbf{agence politique défaillante} (choix des élites post-indépendance). L'extérieur a créé les conditions structurelles de la dépendance (extraversion, frontières, dette), mais l'intérieur a souvent \textbf{reproduit et aggravé} ces structures au lieu de les briser. \textit{Exemple camerounais :} Le pays a des atouts immenses (pétrole, gaz, bois, cacao, café, bauxite, fer, hydroélectricité, diversité agro-écologique). Pourtant, 40 ans de rente pétrolière (1977-2017) n'ont pas financé l'industrialisation ni le capital humain : routes dégradées, hôpitaux délabrés, chômage des jeunes ~60\%. La France et les multinationales ont leur part (accords secrets, évasion fiscale), mais l'État camerounais a choisi la \textbf{prédation de la rente} au lieu de la \textbf{transformation structurelle} (plan d'ajustement structurel 1990 mal négocié, absence de fonds souverains, corruption au sommet). La voie de sortie n'est ni le victimisme, ni l'autoflagellation, mais la \textbf{responsabilité lucide} : bonne gouvernance (transparence, méritocratie) + panafricanisme concret (ZLECAF, infrastructure PIDA, monnaie unique) pour changer les termes de l'échange mondial.
\end{enumerate}
\end{corrige}

\courssub{Exercice 9 — Dissertation : « Le panafricanisme est-il encore une voie pertinente pour le développement de l'Afrique ? »}
\begin{corrige}[Exercice 9 — Dissertation : « Le panafricanisme est-il encore une voie pertinente pour le développement de l'Afrique ? »]
\textbf{Barème indicatif (sur 20) :} Introduction (4) — Développement Partie I (5) — Partie II (5) — Partie III (4) — Conclusion (2). \textbf{Points de vigilance :} Définir « pertinent » (adapté aux défis actuels, efficace). Mobiliser l'historique (Manchester, OUA, UA), les objectifs (intégration, souveraineté), les difficultés (gouvernance, souveraineté nationale, infrastructures, paix). Citer des exemples précis (ZLECAF, UA, PIDA, CEDEAO/CEMAC, crises sahéliennes). Plan dialectique ou analytique.

\vspace{0,5cm}
\noindent\textbf{Introduction}
\vspace{0,2cm}
\noindent\textbf{Accroche :} « L'Afrique doit s'unir ou périr » lançait Kwame Nkrumah en 1963. Soixante ans plus tard, le continent compte 1,4 milliard d'habitants, d'immenses ressources, mais pèse moins de 3\% du PIB mondial et 2\% du commerce global.
\vspace{0,2cm}
\noindent\textbf{Définitions :} Le \cle{panafricanisme} est l'idéologie et le mouvement prônant l'unité politique, économique et culturelle de l'Afrique et de sa diaspora pour assurer son développement autonome. Le \cle{développement} est l'expansion durable des libertés et capacités humaines (Sen). La \cle{pertinence} ici signifie : capacité à répondre aux défis actuels (industrialisation, emploi, climat, géopolitique, paix).
\vspace{0,2cm}
\noindent\textbf{Problématique :} Si le panafricanisme a été le moteur des indépendances, l'échec relatif de l'OUA, la persistance des frontières coloniales, la primauté des souverainetés nationales et les crises sécuritaires actuelles ne le rendent-ils pas obsolète ? Ou bien, au contraire, la ZLECAF et l'UA marquent-elles sa réinvention nécessaire ?
\vspace{0,2cm}
\noindent\textbf{Annonce du plan :} Nous montrerons d'abord que le panafricanisme reste le \textbf{cadre philosophique et stratégique indispensable} (I), puis que ses \textbf{réalisations concrètes récentes prouvent sa vitalité} (II), pour enfin examiner les \textbf{conditions politiques strictes} sans lesquelles il reste un vœu pieux (III).

\vspace{0,5cm}
\noindent\textbf{Développement}
\vspace{0,2cm}
\noindent\textbf{I. Le panafricanisme : cadre philosophique et stratégique indispensable face à la fragmentation}
\vspace{0,1cm}
\noindent 1. \textbf{L'argument de la taille critique :} Aucun État africain seul (même le Nigeria, l'Égypte ou l'Afrique du Sud) n'atteint la taille critique (marché, capital, technologie, diplomatie) pour peser dans la mondialisation. L'unité est une \textbf{nécessité économique} (économies d'échelle, marché de 1,4 Md \$ de consommateurs via ZLECAF) et \textbf{géopolitique} (siège unique au Conseil de sécurité ONU, voix commune climat/dette).
\vspace{0,1cm}
\noindent 2. \textbf{L'héritage décolonial inachevé :} Les frontières de Berlin (1885) morcellent des peuples, créent des États non viables (16 pays enclavés), alimentent les conflits. Le panafricanisme seul propose de \textbf{dépasser l'État-nation colonial} par l'intégration régionale (CEDEAO, CEMAC, SADC, EAC) puis continentale (UA). C'est la condition de la \textbf{seconde libération} (économique), après la première (politique).
\vspace{0,1cm}
\noindent 3. \textbf{La renaissance culturelle comme socle :} Face à l'uniformisation culturelle occidentale, le panafricanisme porte la \textbf{renaissance africaine} (Agenda 2063) : valorisation des langues, savoirs endogènes, histoire commune, identité partagée. Condition \emph{sine qua non} d'une citoyenneté africaine active.

\vspace{0,2cm}
\noindent\textbf{II. Des réalisations concrètes qui attestent de sa vitalité contemporaine}
\vspace{0,1cm}
\noindent 1. \textbf{L'Union africaine (2002) et l'Agenda 2063 :} Passage de la non-ingérence à l'\textbf{non-indifférence} (droit d'intervention en crimes de guerre/génocide). Mécanismes : Conseil de paix et sécurité (CPS), Cour africaine des droits de l'homme, Mécanisme africain d'évaluation par les pairs (MAEP) pour la gouvernance. Vision long terme (2063) : « L'Afrique que nous voulons ».
\vspace{0,1cm}
\noindent 2. \textbf{La ZLECAF (2021) : le grand bond économique :} Ratifiée par 47 pays (2024). Objectif : supprimer 90\% des droits de douane intra-africains, créer un marché unique, booster le commerce intra de 15\% à 25-30\% d'ici 2035 (Banque mondiale : +450 Md \$ de revenus, 30 millions sortis de l'extrême pauvreté). Outils : Système panafricain de paiement (PAPSS), Observatoire du commerce, Protocole sur la libre circulation des personnes (32 signataires).
\vspace{0,1cm}
\noindent 3. \textbf{Intégration régionale moteur :} CEDEAO (libre circulation effective, monnaie unique \emph{Eco} en projet, force en attente), EAC (marché commun, union douanière, passeport unique), SADC. Infrastructures transversales (PIDA) : autoroute Transsaharienne, Lagos-Abidjan, corridor Nord-Sud, barrage Grand Inga (RDC), fibre optique.

\vspace{0,2cm}
\noindent\textbf{III. Les obstacles structurels : la pertinence conditionnée à la volonté politique réelle}
\vspace{0,1cm}
\noindent 1. \textbf{Le piège de la souveraineté nationale :} Les chefs d'État protègent jalousement leurs prérogatives (douanes = rente, armées, présidence). L'UA n'a pas de pouvoir supranational (budget 60\% financé par partenaires extérieurs). La ZLECAF avance mais les \textbf{barrières non tarifaires} persistent (corruption aux frontières, normes divergentes, lenteur ratification protocoles libre circulation).
\vspace{0,1cm}
\noindent 2. \textbf{Le déficit de gouvernance et de paix :} Coups d'État en série au Sahel (Mali, Burkina, Niger, Guinée, Gabon, Tchad) paralysent l'UA (suspensions sans suite efficace). Conflits armés (Soudan, RDC, Somalie, Mozambique, Cameroun anglophone) détruisent le capital humain et physique. Pas de développement sans paix, pas de paix sans gouvernance démocratique légitime.
\vspace{0,1cm}
\noindent 3. \textbf{La dépendance structurelle non résolue :} L'intégration se fait souvent \emph{par le bas} (infrastructures chinoises, financement UE/BAD, modèles occidentaux). L'industrialisation tarde (part manufacturier dans PIB < 10\%). La \textbf{transformation locale} promise par la ZLECAF suppose énergie fiable, logistique performante, main-d'œuvre qualifiée, État stratège : conditions encore largement absentes.

\vspace{0,5cm}
\noindent\textbf{Conclusion}
\vspace{0,2cm}
\noindent\textbf{Réponse :} Oui, le panafricanisme est \textbf{plus pertinent que jamais}, non comme une incantation identitaire, mais comme la \textbf{seule stratégie d'échelle} capable de briser l'extraversion et d'offrir un avenir à la jeunesse africaine. La ZLECAF et l'Agenda 2063 en sont la preuve opératoire.
\vspace{0,1cm}
\noindent\textbf{Nuance :} Sa pertinence n'est pas automatique : elle est \textbf{conditionnelle}. Elle exige le passage de la « rhétorique des chefs d'État » à la « réalité des peuples » : \textbf{gouvernance démocratique crédible}, \textbf{reddition de comptes}, \textbf{investissement massif dans le capital humain}, \textbf{financement autonome} (fiscalité, épargne, diaspora).
\vspace{0,1cm}
\noindent\textbf{Ouverture :} Le véritable test du panafricanisme au XXI\textsuperscript{e} siècle sera sa capacité à \textbf{inventer un modèle de développement propre} (écologique, inclusif, basé sur les communs et l'économie sociale), plutôt que de singer le modèle occidental obsolète. Comme le voulait Cheikh Anta Diop : « L'unité africaine n'est pas un vœu, c'est une nécessité historique. »
\end{corrige}

\courssub{Exercice 10 — Explication de texte : Kwame Nkrumah}
\begin{corrige}[Exercice 10 — Explication de texte : Kwame Nkrumah]
\begin{enumerate}
\item \textbf{Thèse de l'auteur :} L'indépendance politique formelle (drapeau, hymne, siège à l'ONU) n'est qu'un \textbf{prérequis nécessaire mais insuffisant}. La vraie souveraineté — la capacité de décider de son sort et de transformer ses richesses pour son peuple — ne s'obtient que par l'\textbf{unité politique continentale}.
\textbf{Articulation en deux étapes :}
\begin{enumerate}
\item \textbf{Étape 1 : L'indépendance politique nationale.} « Cherchez d'abord le royaume politique. » C'est la conquête de la souveraineté formelle par chaque colonie (années 1957-1960). Nkrumah y a œuvré (Ghana 1957). Mais il sait qu'elle laisse les États « trop faibles ».
\item \textbf{Étape 2 : L'unité politique continentale.} « C'est dans l'union que l'Afrique trouvera la force... » Fusion des souverainetés nationales en un \textbf{Gouvernement continental unique} (États-Unis d'Afrique) : armée commune, monnaie unique, marché commun, politique étrangère unique. Seule cette unité donne le \textbf{rapport de force} face aux puissances extérieures.
\end{enumerate}
\item \textbf{Lien unité politique $\rightarrow$ indépendance économique :}
Pour Nkrumah, l'économie suit le politique. Des États balkanisés (54 aujourd'hui) sont structurellement \textbf{dépendants} : marchés trop petits pour l'industrialisation (pas d'économies d'échelle), pouvoir de négociation nul face aux multinationales minières/pétrolières (concurrence fiscale à la baisse), dépendance aux importations (médicaments, aliment, technologie), endettement auprès des mêmes puissances. L'\textbf{unité politique} crée un \textbf{Espace économique unique} : marché de 300 millions (à l'époque), 1,4 milliard aujourd'hui ; mutualisation des ressources (énergie, eau, minerais) ; planification continentale de l'industrialisation ; monnaie unique souveraine ; capacité à fixer les prix des matières premières (cartel type OPEP). « Tout le reste vous sera donné en plus » : le développement économique est la \textbf{conséquence} de la puissance politique unie.
\item \textbf{Actualité du texte à l'heure de la ZLECAF :}
\textbf{Oui, le texte reste d'une actualité brûlante,} mais la stratégie a bifurqué.
\begin{itemize}
\item \textbf{Actualité confirmée :} La ZLECAF (2021) réalise l'objectif économique nkrumahiste : « transformer ses richesses au profit de ses peuples » par le commerce intra et l'industrialisation. L'UA (2002) incarne l'institutionnalisation politique. Les pressions extérieures (néocolonialisme, dette, climat, géopolitique multipolaire) valident le diagnostic : « trop faibles pour résister ».
\item \textbf{Difficultés persistantes (leçons du cours) :} Nkrumah voulait l'\textbf{unité politique d'abord} (État fédéral), l'économie suivant. L'histoire a imposé l'\textbf{économie d'abord} (ZLECAF), le politique restant intergouvernemental (UA). Les difficultés étudiées expliquent ce réalisme : \textbf{primauté des souverainetés nationales} (chefs d'État réticents à céder pouvoir), \textbf{hétérogénéité des régimes} (démocraties, junte, monarchies), \textbf{conflits armés} (Sahéliens, Soudan, RDC), \textbf{déficit d'infrastructures} (coût transport intra-africain 2-3x mondial), \textbf{dépendance extraversion} (modèles extractifs persistants).
\end{itemize}
\textbf{Conclusion :} Nkrumah avait raison sur le \textbf{diagnostic} (fragmentation = faiblesse) et le \textbf{but} (unité = force). La ZLECAF est l'outil économique qu'il appelait de ses vœux. Mais sans \textbf{progrès parallèle sur l'unité politique} (gouvernance démocratique, paix, fiscalité commune, défense commune), l'indépendance économique restera incomplète. Le « royaume politique » continental reste l'horizon indépassable.
\end{enumerate}
\end{corrige}

\courssub{Exercice 11 — Question à réponse construite : le commerce intra-africain}
\begin{corrige}[Exercice 11 — Question à réponse construite : le commerce intra-africain]
\begin{enumerate}
\item \textbf{Calcul de l'écart :}
Part intra-européenne $\approx 67\%$. Part intra-africaine $\approx 15\%$.
Écart = $67 - 15 = \textbf{52 points de pourcentage}$.
\textbf{Interprétation :} Cet écart abyssal de 52 points révèle l'\textbf{extraversion structurelle} des économies africaines. L'Europe a construit son développement sur l'intégration régionale (marché unique, UE) : ses échanges internes sont le double de ses échanges externes. L'Afrique, elle, exporte 85\% de son commerce vers l'extérieur (Europe, Chine, USA, Inde) et ne commerce qu'à 15\% avec elle-même. C'est la marque du \textbf{sous-développement intégré} : les économies africaines sont plus connectées à l'ancien colonisateur ou aux nouvelles puissances qu'entre elles, reproduisant le schéma colonial (matières premières contre manufacturés).
\item \textbf{Deux difficultés du panafricanisme révélées par le document :}
\begin{enumerate}
\item \textbf{Le déficit d'intégration économique réelle (infrastructures et barrières) :} Le faible chiffre (15\%) prouve que les marchés nationaux restent cloisonnés. Causes : infrastructures de transport/énergie transversales manquantes (routes, rails, électricité, internet), barrières non tarifaires (corruption aux frontières, paperasserie, normes divergentes), monnaies non convertibles (zones CFA vs autres), absence de système de paiement continental (résolu partiellement par PAPSS).
\item \textbf{La primauté de la souveraineté nationale sur l'intérêt continental :} Les États protègent leurs recettes douanières (souvent 30-50\% des recettes fiscales), leurs marchés captifs pour leurs élites économiques, et leur pouvoir politique. La ratification lente des protocoles (libre circulation des personnes, droit de résidence, d'établissement) et la mise en œuvre partielle de la ZLECAF (listes d'exclusion, règles d'origine) illustrent ce blocage politique.
\end{enumerate}
\item \textbf{Deux solutions concrètes inspirées du panafricanisme :}
\begin{enumerate}
\item \textbf{Mise en œuvre effective, rapide et irréversible de la ZLECAF (libéralisation à 90\% + règles d'origine simples + PAPSS + Observatoire du commerce).}
\textbf{Justification :} La suppression des droits de douane sur 90\% des lignes tarifaires (sur 10-13 ans) crée immédiatement de la demande intra-africaine, incite à la localisation de la production (règles d'origine « valeur ajoutée 40\% » ou « changement de position tarifaire »), réduit les coûts de transaction via le PAPSS (paiement en monnaies locales, évitant le dollar), et l'Observatoire permet le suivi transparent. C'est le \textbf{moteur industriel} promis par Nkrumah.
\item \textbf{Programme prioritaire d'infrastructures transversales « dures » et « molles » (PIDA + Corridors + Énergie + Numérique) financé par l'épargne africaine (Fonds souverain panafricain, Diaspora Bonds, fiscalité minière harmonisée).}
\textbf{Justification :} Sans routes, rails, ports, électricité fiable et internet haut débit reliant les capitales et zones de production, le commerce intra-africain reste théorique (coût transport africain $\approx$ 2-3$\times$ moyenne mondiale). Le panafricanisme impose la \textbf{mutualisation des grands ouvrages} (barrage Grand Inga, autoroute Transsaharienne, chemin de fer Dakar-Ndjamena, fibre optique continentale) pour réduire la distance économique. Le financement par des ressources africaines (prélèvement sur rentes minières/pétrolières, obligations de la diaspora) garantit l'autonomie de décision.
\end{enumerate}
\end{enumerate}
\end{corrige}

\newpage

\coursec{Fiche bilan}

\begin{popfigure}
\begin{tikzpicture}[
  mindmap,
  grow cyclic,
  every node/.style={concept, font=\small},
  root concept/.append style={concept color=popblue, font=\small\bfseries, minimum size=2.6cm},
  level 1 concept/.append style={level distance=3.6cm, sibling angle=60, font=\footnotesize},
  level 2 concept/.append style={level distance=2.2cm, font=\scriptsize, minimum size=1.4cm}
]
\node[root concept]{Panafricanisme et développement de l'Afrique}
  child[concept color=popgreen]{ node {Développement}
    child { node {Économique, social, humain} }
    child { node {IDH, PIB} }
  }
  child[concept color=poporange]{ node {Sous-développement}
    child { node {Pauvreté, dépendance} }
    child { node {Dualisme, domination} }
  }
  child[concept color=popgold]{ node {Historique}
    child { node {1900 : congrès de Londres} }
    child { node {1945 : Manchester} }
    child { node {1963 : OUA, 2002 : UA} }
  }
  child[concept color=popblue!60]{ node {Objectifs}
    child { node {Unité politique} }
    child { node {Indépendance économique} }
    child { node {Dignité culturelle} }
  }
  child[concept color=poppink]{ node {Difficultés}
    child { node {Frontières héritées} }
    child { node {Néocolonialisme} }
    child { node {Conflits, égoïsmes nationaux} }
  }
  child[concept color=popgreen!60]{ node {Auteurs}
    child { node {Nkrumah, Senghor} }
    child { node {Du Bois, Padmore} }
  };
\end{tikzpicture}
\legende{Carte mentale de la leçon : le panafricanisme comme philosophie du développement de l'Afrique.}
\end{popfigure}

\begin{bilan}
\textbf{Définitions clés}
\begin{itemize}
  \item \cle{Développement} : processus global de transformation économique, sociale, politique et culturelle qui améliore durablement les conditions de vie d'une population (croissance + changements structurels + épanouissement humain).
  \item \cle{Sous-développement} : état d'une société dont les potentialités sont bloquées ou exploitées au profit de l'extérieur ; ce n'est pas un retard naturel mais le produit d'un rapport de \cle{domination} (thèse d'A. G. Frank et S. Amin : le sous-développement est \emph{engendré} par le développement des autres).
  \item \cle{Panafricanisme} : mouvement de pensée et d'action visant l'unité des Africains du continent et de la diaspora, leur libération et leur développement solidaire.
\end{itemize}

\textbf{Repères à connaître par cœur}
\begin{center}
\begin{tabularx}{\linewidth}{|l|X|}
\hline
\rowcolor{popblueL} \textbf{Date / Auteur} & \textbf{Idée essentielle} \\
\hline
1900 & Premier congrès panafricain (Londres), W. E. B. Du Bois. \\
\hline
1945 & Congrès de Manchester : passage à l'action anticoloniale (Nkrumah, Padmore). \\
\hline
1963 & Création de l'OUA à Addis-Abeba ; 2002 : Union africaine (UA). \\
\hline
Nkrumah & « L'indépendance politique n'est rien sans l'indépendance économique » ; l'unité conditionne le développement. \\
\hline
Senghor & Unité dans la culture et la négritude ; métissage des civilisations. \\
\hline
\end{tabularx}
\end{center}

\textbf{Méthode express (dissertation)}
\begin{enumerate}
  \item Définir les deux concepts clés (développement, panafricanisme) dès l'introduction.
  \item Problématiser : l'unité africaine est-elle une condition suffisante du développement ?
  \item Construire un plan dialectique : thèse (le panafricanisme comme voie du développement) / antithèse (ses échecs et obstacles) / synthèse (une unité réaliste, économique d'abord : ZLECAf, intégration régionale).
  \item Illustrer avec des exemples camerounais et africains (CEMAC, UA, NEPAD).
\end{enumerate}
\end{bilan}

\begin{aretenir}[Checklist avant le Bac]
\begin{itemize}
  \item $\square$ Je sais définir le développement dans ses trois dimensions (économique, sociale, humaine).
  \item $\square$ Je distingue croissance économique et développement.
  \item $\square$ Je sais définir le sous-développement et citer au moins trois manifestations (pauvreté, dépendance, analphabétisme).
  \item $\square$ Je connais les grandes étapes du panafricanisme : 1900, 1945, 1963, 2002.
  \item $\square$ Je peux citer au moins trois pères du panafricanisme (Du Bois, Nkrumah, Senghor).
  \item $\square$ Je connais les objectifs du panafricanisme : unité politique, indépendance économique, dignité culturelle.
  \item $\square$ Je sais expliquer les difficultés : frontières coloniales, néocolonialisme, rivalités entre États.
  \item $\square$ Je connais la thèse de la dépendance (Frank, Amin) : le sous-développement est produit par la domination.
  \item $\square$ Je sais construire un plan dialectique sur « le panafricanisme peut-il développer l'Afrique ? ».
  \item $\square$ Je dispose d'exemples concrets : OUA/UA, CEMAC, ZLECAf, NEPAD.
  \item $\square$ Je sais rédiger une introduction avec définition des termes, problématique et annonce du plan.
\end{itemize}
\end{aretenir}

\findecours

\end{document}
